\documentclass[12pt]{article}
\usepackage{graphicx} % Required for inserting images
\usepackage{amsmath,amssymb,amsthm}
\newtheorem{teo}{Teorema} 
\newtheorem{defi}{Definicion} 
\newtheorem{eg}{Ejemplo} 
\newtheorem{cor}{Corolario} 
\newtheorem{prop}{Proposición} 
\newtheorem{lem}{Lema}
\usepackage{braket}
\newtheorem{ex}{Ejercicio} 
\newtheorem{post}{Postulado}
\usepackage{subfigure}
\usepackage{epigraph}
\usepackage[utf8]{inputenc}
\usepackage[]{mdframed}
\usepackage[spanish]{babel}
\usepackage{url}
\usepackage{wrapfig}
\usepackage{graphicx}
\usepackage{amsmath}
\usepackage{braket}
\usepackage{graphicx}
\usepackage[all]{xy}
\usepackage{fancyhdr}
\usepackage{halloweenmath}
\usepackage{url}
\usepackage{wrapfig}
\usepackage{graphicx}
\usepackage{amsmath}
\usepackage{array}
\usepackage[utf8]{inputenc}
\usepackage[export]{adjustbox}
\usepackage{float}
\usepackage[margin=3cm]{geometry}
\usepackage{amssymb}
\usepackage[all]{xy}
\usepackage{hyperref}
\usepackage{multicol}
\newcommand{\squig}{{\scriptstyle\sim\mkern-3.9mu}}
\newcommand{\lsquigend}{{\scriptstyle\lhd\mkern-3mu}}
\newcommand{\rsquigend}{{\scriptstyle\rule{.1ex}{0ex}\rhd}}
\newcounter{sqindex}
\newcommand\squigs[1]{%
  \setcounter{sqindex}{0}%
  \whiledo {\value{sqindex}< #1}{\addtocounter{sqindex}{1}\squig}%
}
\newcommand\rsquigarrow[2]{%
  \mathbin{\stackon[2pt]{\squigs{#2}\rsquigend}{\scriptscriptstyle\text{#1\,}}}%
}
\newcommand\lsquigarrow[2]{%
  \mathbin{\stackon[2pt]{\lsquigend\squigs{#2}}{\scriptscriptstyle\text{\,#1}}}%
}
\usepackage{upgreek}
\usepackage{pgfplots,tikz}
\usetikzlibrary {arrows.meta,bending,positioning}
 \usetikzlibrary{babel}
 \renewcommand{\abstractname}{Resumen}
  \renewcommand{\figurename}{Figura}


\title{Guión general hilos de grupos de Lie}
\author{Ós }
\date{\ }

\begin{document}
\maketitle
\begin{center}
    We need a super-mathematics in which the operations are as unknown as the quantities they operate on, and a super-mathematician who does not know what he is doing when he performs these operations. \\ Such a super-mathematics is the Theory of Groups.  
\end{center}
\begin{center}
     \textit{Sir Arthur Stanley Eddington}
\end{center}

\section{In the beggining there was nothing and before there was nothing there were rotaciones}
\epigraph{Before there was time, before there was anything, there was nothing. And before there was nothing, there were monsters }{\textit{The Lich en Gold Star 6.26}}
Más o menos, me gustaría que esto empezara con cosas muy básicas y he pensado que lo mejor para empezar a hablar de grupos es hablar de las rotaciones en el plano. Una rotación de un vector en el plano es lo que podéis suponer que es, mover el vector dejando su inicio fijo y, por supuesto, sin poder cambiar su módulo. Esto hace que solo pueda moverse describiendo una circunferencia. 

\begin{center}
    \begin{tikzpicture}
   \draw[very thick, red,->](0,0)--(1.5,2);
   \draw[very thick, blue,->](0,0)--(2.29,1);
    \end{tikzpicture}
\end{center}
\begin{center}
    Fig 1. Rotación del vectores, podemos ver como el vector azul es la rotación del rojo y viceversa. 
\end{center}
Esta es la forma intuitiva de ver las rotaciones y perfectamente podríamos hacer un poco de geometría y ver como son, pero vamos a pensar un poco qué reglas generales deben cumplir las rotaciones de vectores en el plano. Primero de todo, el anterior ejemplo no usaba bases, porque simplemente lo estamos haciendo de forma intuitiva, pero cuando introducimos las bases nos topamos con una sorpresa, esto es porque la anterior rotación puede ponerse como 
\begin{center}
    \begin{tikzpicture}
      \draw[very thick, black,->](0,0)--(2.5,0);
         \draw[very thick, red,->](0,0)--(1.5,2);
   \draw[very thick, blue,->](0,0)--(2.29,1);
      \draw[very thick, black,->](0,0)--(0, 2.5);
    \end{tikzpicture}
\end{center}
\begin{center}
    Fig 2. Rotación activa. 
\end{center}
En la que giramos el vector y la base se queda igual. Este tipo de rotaciones se llama activa. Aquí, por si no queda claro, estamos usando la base canónica de $\mathbb{R}^2$ que yo voy a denotar como $\{\hat x, \hat y\}$. PERO, hay otra forma de ver las rotaciones y es rotar la base\\\\

\begin{center}
    \begin{tikzpicture}
      \draw[very thick, black,->](0,0)--(2.5,0);
     \draw[very thick, blue,->](0,0)--(0.9,2.33);
         \draw[very thick, blue,->](0,0)--(-2.33,0.9);
      \draw[very thick, black,->](0,0)--(0, 2.5);
\draw[very thick, red,->](0,0)--(1.5,1.5);
    \end{tikzpicture}
\end{center}



%Dibuix de la rotación de la base 
Esta forma de rotación se llama pasiva. Ambas son formas distintas, pero equivalentes, de describir las rotaciones. Pero de ambas obtenemos una valiosa lección. De la primera, la obviedad, que es que el módulo del vector debe permanecer invariante, de la segunda, que transforma bases ortogonales en bases ortogonales. Es decir, en la rotación pasiva, el ángulo que forman los vectores de la base se preserva. Como sabéis, el ángulo entre dos vectores lo da el producto escalar por tanto, una rotación debe ser una aplicación que preserve el producto escalar. Esto es, sea $\{e_i\}$ una base, 
\begin{equation}
    R(e_i) \cdot R(e_j)=e_i\cdot e_j
\end{equation}
Veamos si desarrollando esto podemos llegar a algún tipo de ecuación que nos definan las matrices de rotación. Recordad que el producto escalar $x\cdot y$ puede expresarse matricialmente como $x^T y$ donde $X^T$ es el vector $x$ pero en horizontal. Entonces, la condición de que se preserve el producto escalar se escribe
$$(R e_i)^T Re_j = e_i^T R^T R e_j = e_i^T e_j$$
Esta igualdad solo se alcanza si
\begin{equation}
    R R^T =I
\end{equation} 
A las matrices $R$ que cumplen esto las llamaremos ortogonales. Las matrices ortogonales son aquellas que preservan el producto escalar. Esto implica, por tanto, que los módulos se conservan!! que es la condición que tenían que cumplir las rotaciones en las rotaciones activas.  Una forma de verlo es tomar determinantes a ambos lados de la condición de que una matriz sea ortogonal 
\begin{equation*}
    \begin{split}
      &  R R^T =I \\
      &  \text{det} (R R^T) = \text{det}I \\
      &  \text{det} R \ \text{det} R^T =1 \\
       & (\text{det} R)^2 =1
\end{split}
\end{equation*}
De donde se deduce que $\text{det} R =\pm 1$ o sea, las transformaciones ortogonales pueden tener o determinante $+1$ o determinante $-1$. Nos gustaría que la identidad sea una rotación porque no hacer nada sigue siendo una rotación, la de $0$ grados, por tanto, tenemos que las rotaciones son las transformaciones ortogonales de terminante unidad. 
\begin{equation}
    SO(2) =\{R\in M_2(\mathbb{R}) | R R^T=I, \text{det}R =1\}
\end{equation}
Aquí, $ M_2(\mathbb{R})$ es el conjunto de matrices cuadradas $2\times 2$ de entradas reales. Hemos llamado a este conjunto $SO(2)$ \footnote{El lector adelantado puede pensar que hemos encontrado nuestro primer grupo de Lie. PUES NO. Si bien es un grupo, no tenemos definida ninguna topología y, cosa que veremos que es totalmente necesari, por lo que no podemos hablar de que sea un grupo de Lie \textit{per se}. Estamos ante un grupo, sí, pero de momento simplemente algebraico.  }. La $O$ viene de ''Orthogonal'' y la $S$ viene de ''Special'',  significando esto que tiene determinante unidad. ¿El $2$ tengo que explicarlo también? .\\\\
Bien, hemos encontrado la forma algebraica de las matrices de rotación, pero como decíamos al principio, siempre hemos podido encontrarlas por simple trigonometría. Veamos como hacerlo.
\begin{center}
    \begin{tikzpicture}
      \draw[very thick, black,->](0,0)--(4,0);
         \draw[very thick, red,->](0,0)--(1.5,3);
   \draw[very thick, blue,->](0,0)--(3.29,1);
      \draw[very thick, black,->](0,0)--(0, 4);
      %proyecciones
      \draw[dotted, red](1.5,3)--(1.5,0);
      \draw[dotted, red](1.5,3)--(0,3);
       \draw[dotted, blue](3.29,1)--(0,1);
       \draw[dotted, blue](3.29,1)--(3.29,0);
    \end{tikzpicture}
\end{center}
\begin{center}
\end{center}
En término de los ángulos tenemos entonces dos vectores:
\begin{equation*}
    \begin{split}
        & \vec v =v_x\hat x+v_y \hat y= \cos \alpha \hat x +\sin \alpha \hat y \\
        &\vec v' =v_x'\hat x+v_y' \hat y= \cos (\alpha +\beta) \hat x +\sin (\alpha +\beta) \hat y
    \end{split}
\end{equation*}
Pero claro, nos gustaría poner $v'$ como una combinación lineal de términos que solo dependan del ángulo, puesto de rotación $\beta$ puesto que si es una rotación no puede ser de otra forma. Aplicamos entonces las formulas del coseno y seno de la suma \footnote{Estas deben de ser las únicas fórmulas del cole de las que nunca, pero nunca, me acuerdo. No sé si a vosotros también os pasa. Pero, qué casualidad, por la teoría de grupos vamos a ver como se demuestran de forma muy sencilla, tanto que siempre que uno, como yo, no se acuerde, puede hacer tres lineas y ya tendrá las fórmulas. }
\begin{equation*}
    \begin{split}
        & v'_x=\cos(\alpha +\beta) =\cos\alpha \cos\beta -\sin\alpha\sin\beta = v_x \cos\beta -\sin\beta v_y \\ 
        & v'_y=\sin(\alpha +\beta) =\sin\alpha\cos\beta +\cos\alpha\sin\beta = v_x\sin\beta + v_y\cos\beta
    \end{split}
\end{equation*}
Buscamos una matriz $\hat R$ tal que $\vec v' = R \vec v$ por tanto comparando componentes se deduce la forma general de la matriz de rotación \footnote{De aquí en adelante vamos a tomar esto como definición, el motivo es porque claramente tenemos un grado de libertad que es la orientación del giro, en este caso, y así será nuestra definición, antihorario.}
\begin{equation}
    R_\beta=\begin{pmatrix}
        \cos \beta & -\sin\beta \\ \sin\beta & \cos \beta
    \end{pmatrix}
\end{equation}
Por supuesto,hemos deducido esta matriz y de forma muy sencilla desde la vista activa de las rotaciones, pero debería ser evidente que mientras en el caso activo esta es una matriz que transforma componentes, en el punto de vista pasivo esta matriz es una matriz de cambio de base, pero es la misma matriz. Por tanto, hemos encontrado por simple geometría una forma de \textbf{representar} el conjunto algebraico y definido de forma abstracta arriba de matrices $SO(2)$. Es muy fácil comprobar que este tipo de matrices son ortogonales y de determinante unidad, como esperábamos. \\\\ Ahora que tenemos las rotaciones, podemos aprender varias cosas de ellas.\\
La primera cosa que podemos ver es que si juntas dos rotaciones vas a obtener otra rotación. Intuitivamentees evidente pero ahora que tenemos estas matrices lo podemos demostrar. Sean dos rotaciones de ángulo $\alpha$ y $\beta$ representadas por las matrices $R_\alpha$ y $R_\beta$, entonces su producto es otra rotación. Para verlo simplemente hagamos el producto 
\begin{equation*}
      \begin{pmatrix}
        \cos \beta & -\sin\beta \\ \sin\beta & \cos \beta
    \end{pmatrix} \begin{pmatrix}
        \cos \alpha & -\sin\alpha\\ \sin\alpha & \cos \alpha
    \end{pmatrix} =\begin{pmatrix}
        \cos\beta \cos \alpha -\sin\beta \sin\alpha & -\sin\alpha \cos\beta -\sin\beta\cos \alpha \\
        \sin\beta\cos\alpha +\cos\beta \sin \alpha & -\sin\beta\sin\alpha +\cos\beta\cos\alpha 
    \end{pmatrix}
\end{equation*}
Si usamos las fórmulas del seno y coseno de la suma entonces 
\begin{equation*}
      \begin{pmatrix}
        \cos \beta & -\sin\beta \\ \sin\beta & \cos \beta
    \end{pmatrix} \begin{pmatrix}
        \cos \alpha & -\sin\alpha\\ \sin\alpha & \cos \alpha
    \end{pmatrix} =\begin{pmatrix}
        \cos (\alpha +\beta) & -\sin(\alpha +\beta)\\ \sin(\alpha +\beta) & \cos (\alpha +\beta)
    \end{pmatrix} =\begin{pmatrix}
        \cos \gamma & -\sin \gamma \\ \sin \gamma & \cos\gamma
    \end{pmatrix}
\end{equation*}
Donde la última igualdad simplemente está para subrayar el hecho de que $\alpha+\beta$ es un nuevo ángulo $\gamma$ que tiene otra matriz de rotación asociada. Concluimos, efectivamente, que el producto de dos rotaciones es otra rotación. Pero además vemos una cosa interesante, y es que el producto de rotaciones se ha traducido en una suma de ángulos. \footnote{Esto será lo que usaremos para definir una topología en $SO(2)$, relacionado con la nota 2. }\\
Bien, otra cosa que podemos notar de estas rotaciones es que son asociativas, cosa que en este caso es muy fácil de ver. \\
% dibuix? formula?
Otra observación, que parece una tontería, es que tenemos dentro del conjunto una identidad. Esto es, tenemos una rotación que no hace nada, la de $0$ grados. Y de aquí vemos otra cosa, si rotamos $\alpha$ podemos también rotar $-\alpha$ o sea, el mismo ángulo en la misma dirección, lo que equivale a no hacer nada
\begin{equation*}
    \begin{pmatrix}
         \cos \alpha & -\sin\alpha\\ \sin\alpha & \cos \alpha  
    \end{pmatrix} \begin{pmatrix}
         \cos (-\alpha) & -\sin(-\alpha)\\ \sin(-\alpha) & \cos (-\alpha)
    \end{pmatrix} = \begin{pmatrix}
         \cos \alpha & -\sin\alpha\\ \sin\alpha & \cos \alpha  
    \end{pmatrix}   \begin{pmatrix}
        \cos \alpha & \sin\alpha\\ -\sin\alpha & \cos \alpha
    \end{pmatrix} =\begin{pmatrix}
        1 & 0 \\ 0 & 1
    \end{pmatrix}
\end{equation*}
Hay una cosa más. Fijaos que tenemos una rotación para cada ángulo, un ángulo no es más que un elemento $\theta \in [0,2\pi] \subset \mathbb{R}$ Concretamente, si estamos en $2\pi$ y rotamos un poco más, volvemos a empezar, esto es porque no estamos en intervalo de la recta, claro, sino en un círculo, que representamos por $S^1$. Por tanto, para cada $\theta \in S^1$ existe una rotación $R_\theta$, topológicamente hablando (formalizaremos esto después), $SO(2)$ tiene la misma topología que $S^1$. \\ La propiedad topológica que nos interesa aquí es la continuidad, esto en este caso significa que si tengo un $\theta \in S^1$ y un $\varepsilon\in S^1$ tal que $\varepsilon^2=0$ (o sea, es pequeñísimo) entonces $\theta +\varepsilon \in S^1$ los angulos, no tienen ''huecos'' entre ellos. Pero es que resultado que el ángulo $\varepsilon$ ya es interesante de por si. ¿Cómo es una rotación de ángulo $\varepsilon$? Pues, como es pequeñísimo, debe ser una transformación que sea una desviación de orden $\varepsilon$ de la identidad, o sea 
\begin{equation*}
    R(\varepsilon)= I +\varepsilon J
\end{equation*}
¿Y qué es $J$? Para ver qué es, vamos a recordar la definición de derivada (según Cauchy) :
\begin{equation*}
    f'(a)=\lim_{\varepsilon\rightarrow 0}\frac{f(a+\varepsilon)-f(a)}{\varepsilon}
\end{equation*}
Supongamos ahora una rotación $R_\theta$ y una rotación de ángulo $R_{(\theta+\varepsilon)}$ podemos entonces decir, de manera muy burda que 
\begin{equation*}
    R(\theta)' =  \lim_{\varepsilon\rightarrow 0}\frac{R(\theta+\varepsilon)-R(\theta)}{\varepsilon}
\end{equation*}
Si ahora $\theta =0$, osea, derivamos y luego $\theta =0$, (no es que estemos derivando la identidad, claro), tendremos 
\begin{equation*}
    R(0)' =   \lim_{\varepsilon\rightarrow 0}\frac{1}{\varepsilon}(R(\varepsilon)-R(0)) =\lim_{\varepsilon\rightarrow 0}\frac{1}{\varepsilon}(R(\varepsilon)-I)=\lim_{\varepsilon\rightarrow 0}\frac{1}{\varepsilon}(I+\epsilon J-I) =J 
\end{equation*}
La derivada de la matriz diremos que es la derivada de los elementos. Por tanto
\begin{equation*}
    J=R' (0) = \left. \frac{d}{d\theta}\begin{pmatrix}
        \cos \theta & -\sin\theta \\ \sin\theta & \cos\theta 
    \end{pmatrix}\right|_{\theta =0 } =\left. \begin{pmatrix}
        -\sin\theta & -\cos \theta \\ \cos \theta &-\sin\theta 
    \end{pmatrix} \right|_{\theta=0} =\begin{pmatrix}
        0 & -1 \\ 1 & 0
    \end{pmatrix}
\end{equation*}
Vamos a jugar con esta matriz y con jugar me refiero a multiplicarla por si misma
\begin{equation*}
J^2=
\begin{pmatrix}
        0 & -1 \\ 1 & 0
    \end{pmatrix}
    \begin{pmatrix}
        0 & -1 \\ 1 & 0
    \end{pmatrix}
    = \begin{pmatrix}
        -1  & 0 \\ 0 & -1 
    \end{pmatrix} =- I
\end{equation*}
Una propiedad interesante ¿os recuerda a algo? Hagamos otra potencia
\begin{equation*}
J^3=-
\begin{pmatrix}
        0 & -1 \\ 1 & 0
    \end{pmatrix}
    \begin{pmatrix}
        1 & 0 \\ 0 & 1
    \end{pmatrix}
    =-J
\end{equation*}
¿Todavía no lo ves?, bueno, hagamos dos mas...
\begin{equation*}
J^4=-
\begin{pmatrix}
        0 & -1 \\ 1 & 0
    \end{pmatrix}\begin{pmatrix}
        0 & -1 \\ 1 & 0
    \end{pmatrix}
    = - J^2 = I 
\end{equation*}
\begin{equation*}
J^5=
\begin{pmatrix}
        0 & -1 \\ 1 & 0
    \end{pmatrix}\begin{pmatrix}
        1 & 0 \\ 0 & 1
    \end{pmatrix}
    = J 
\end{equation*}
Ahora sí ¿no? \\ Las potencias de $J$ son las mismas que las del número imaginario $i=\sqrt{-1}$!!!\\ 
Claro, esto, entonces, ¿Qué implica? Los complejos (de módulo 1) se pueden construir con la asignación $\theta \mapsto e^{i\theta} $ donde $\theta$ es un ángulo, o sea, un elemento del círculo $S^1$. Supongamos ahora que hacemos la identificación $i \mapsto J$ entonces 
$$e^{\theta J}$$
Tiene sentido. Ojo. De momento, no estamos justificando nada, solo estamos haciendo suposiciones graciosas viendo que las cosas se parecen. Entramos ahora en un problema descomunal y es que ¿qué es exponenciar una matriz? Bueno, de momento no lo sabemos, pero sí sabemos que
$$e^{i\theta} =1\cdot \cos\theta + i\sin\theta $$
Que es la famosísima fórmula de Euler. 



La asignación $i\mapsto J$ nos hace que $1\mapsto I$  y tendremos
\begin{equation*}
    \begin{split}
         e^{\theta J} & = I \cos\theta+ J\sin\theta = \begin{pmatrix}
        1 & 0 \\ 0 & 1 
    \end{pmatrix} \cos \theta + \begin{pmatrix}
        0 & -1 \\ 1 & 0 
    \end{pmatrix} \sin\theta  \\
    & =  \begin{pmatrix}
        \cos\theta & 0 \\ 0 & \cos \theta 
    \end{pmatrix}+ \begin{pmatrix}
        0 & -\sin \theta \\ \sin\theta & 0 
    \end{pmatrix}= \begin{pmatrix}
        \cos \theta  & -\sin\theta \\ \sin \theta & \cos\theta 
    \end{pmatrix}
    \end{split}
\end{equation*}
Nos llevamos una sorpresa. Hemos recuperado la matriz de rotación. ¿Cómo es esto? La identidad de Euler la obteníamos del desarrollo en serie de la exponencial 
\begin{equation}
    e^x=\sum_{n=0}^{\infty} \frac{x^n}{n!}
\end{equation}
Siguiendo esto, vamos a definir la exponencial de un operador $A$ como 
\begin{equation}
    e^A=\sum_{n=0}^{\infty} \frac{A^n}{n!}
\end{equation}
El operador en nuestro caso es $\theta J$, de la misma forma podemos demostrar una fórmula de Euler, por eso, tiene sentido la asignación anterior y por eso obtenemos el mismo resultado. Pero, ¿Tiene esto sentido o es simplemente una casualidad? Veamoslo. \\\\
Estaréis de acuerdo conmigo en que hacer una rotación es rotar muchas veces poco a poco, dicho de otra forma, una rotación de ángulo $\theta$ es el producto de las $N$ rotaciones de ángulo $\theta/N$.
$$ $$
Cuanto más grande sea $N$ más pequeño será $\theta/N$ y por tanto, tendremos algo así como 
\begin{equation*}
   R(\theta) = (R(\theta/N))^N= \lim_{N\rightarrow \infty} \left[I + \left(\frac{\theta}{N} \right) J\right]^N
\end{equation*}
Vamos a expandir el binomio de destro del límite
\begin{equation*}
    \begin{split}
        \left[I + \left(\frac{\theta}{N} \right) J\right]^N & = I+ N I \left(\left(\frac{\theta}{N} \right) J\right) +\frac{N(N-1)}{2}I \left(\left(\frac{\theta}{N} \right) J\right)^2+... \\ 
        & = I+\theta J +\frac{N(N-1)}{2 N^2}\theta^2 J^2+ \dots 
    \end{split}
    \end{equation*}
Cuando tomemos el límite, cada uno de los coeficientes que multiplican a $(\theta J)^N$ tiende a $1$ por tanto 
\begin{equation*}
    \begin{split}
        \lim_{N\rightarrow \infty} \left[I + \left(\frac{\theta}{N} \right) J\right]^N & = I+\theta J+ \frac{1}{2}\theta^2 J^2 +... =e^{\theta J}
    \end{split}
\end{equation*}
Por tanto, acabamos de demostrar que
\begin{equation}
     R(\theta) = (R(\theta/N))^N= \lim_{N\rightarrow \infty} \left[I + \left(\frac{\theta}{N} \right) J\right]^N =e^{\theta J}
\end{equation}
Así que  lo que veíamos antes no era mera casualidad, sino que tiene todo el sentido del mundo. \\\
Estamos haciendo entonces la asiganación $\theta J \mapsto e^{\theta J} $ a partir de la aplicación exponencial. Podemos entonces entender $\theta J$ como un elemento de un cierto espacio vectorial donde $J$ será el elemento de la base. Por tanto, este espacio vectorial es un espacio real unidimensional y tendrá la misma forma que la recta real. Recapitulando, el conjunto de rotaciones, tenía ciertas propiedades al introducir la composición de rotaciones muy interesantes y, además, parecía presentar la misma topología que el círculo y ''al derivar'' un elemento y evaluarlo en la identidad $\theta =0$ parece que obtenemos un espacio vectorial que es $\mathbb{R}$. En resumen tenemos la siguiente situación:
\begin{center}
\begin{tikzpicture}
\draw[very thick, black](2,-3)--(2,3);
\filldraw[color=black, fill=white, very thick](0,0) circle (2);
      \filldraw[color=black, fill=black, very thick](2,0) circle (0.05);
        \node at (2,0) (A) [below left]{$\theta = 0$};
\end{tikzpicture}
\end{center} 
¿Cómo? Bueno. Quedaos con este dibujo y creedme de momento es que representa todo lo que acabamos de ver, en seguida le daremos todo el sentido del mundo.

\section{Grupos y simetría}
\epigraph{God, Thou great symmetry }{\textit{Anna Wickham. 'Envoi'}}  \noindent Para el epígrafe ver esta nota. \footnote{El poema completo dice así: \\ God, Thou great symmetry, \\ Who put a biting lust in me \\
From whence my sorrows spring, \\
For all the frittered days \\
That I have spent in shapeless ways \\
Give me one perfect thing. }
%\subsection{Qué es un Grupo}
%\textit{En esta sección introduciremos qué es un grupo como las propiedades que cumplían las matrices de rotación sin tener en cuenta cosas de topología, y daremos el ejemplo más tonto que hay de como un grupo puede representar simetrías. }\\
%Sin tapujos, enunciamos la definicón algebraica de grupo: 
En este hilo vamos a tratar de ver la relación entre los grupos y la simetría. \\\\
Simetría, dice Hermann Weyl y en mi opinión tiene razón, se entiende en la vida cotidiana principalmente de dos formas. La primera es relacionada con la armonía, la correcta proporción o el equilibrio. La simetría denota la corcondancia de las distintas partes que conforman un todo. La belleza está unida a la simetría. En este sentido, la idea de simetría no está restringida, de ninguna manera a objetos espaciales. La otra forma en la que se usa la palabra simetría, dice un poco más adelante, es para referirse a la simetría bilateral. La simetría bilateral es estrictamente geométrica y, en contraste con la anterior idea común de simetría, tiene un adefinición absolutamente precisa. Una configuración espacial es simétrica con respecto de un plano $E$ si  se llega a si misma por una reflexión por $E$. 
\begin{center}
    \begin{tikzpicture}
    \draw[thin, black](-4,0)--(4,0);
     \draw[very thick, red](0,-2)--(0,2);
     \filldraw[color=black, fill=black, very thick](2,0) circle (0.05);
      \filldraw[color=black, fill=black, very thick](-2,0) circle (0.05);
        \node at (2,0) (A) [above]{$p$};
        \node at (-2,0) (B) [above ]{$p'$};
        \node at (0,-2) (C) [above right][red]{$E$};
\end{tikzpicture}
\end{center}
Una reflexión por $E$ es una aplicación del espacio hasta el mismo  que lleva cualquier punto arbritrario $p$ hasta su imagen espejo $p'$ con respecto de $E$, $S:p\mapsto p'$. La clave de esto es que TODOS los puntos del espacio se mapean a un ÚNICO punto imagen. Existen otro tipo de asignaciones espaciales de este tipo, por ejemplo, las rotaciones. Decimos entonces que una figura tiene simetría rotacional alrededor de un eje si se llega a si misma por una rotación alrededor de $l$. La simetría bilateral entonces aparece entonces solo como el primer caso de concepto geométrico de simetría que se refiere a las operaciones de reflexión y rotación.
Este tipo de simetrías se llaman isometrías. Ambas ideas de simetría aunque no parecen la misma, acaban siéndolo si nos abstraemos un poco:\\  Diremos que una transformación es simetría de una cierta estructura si preserva dicha estructura. Ya está, tan sencillo como eso.\\ En el caso de las isometrías la estructura que se mantiene es simplemente la configuración geométrica que se tiene, sea la que sea. Pero esta noción de simetría es muy poderosa. Fijaos en que estamos diciendo que dada una 'estructura' (simétrica) hay un conjunto de transformaciones (de simetría). En este sentido, estudiando la 'estructura' obtenemos las transformaciones de simetría.  Tomemos un triángulo equilátero.  Rotar un triángulo equilátero en general siempre nos va a dar un triángulo equilátero pero, la configuración geométrica inicial, eso que llamábamos estructura, solo se preservará bajo reflexiones por las medianas del triángulo y bajo las rotaciones de $120^\circ$. 
\begin{center}
    \begin{tikzpicture}
        \draw[ultra thick, black](-2,0)--(2,0)--(0,3.5)--cycle;
        
    \end{tikzpicture}
\end{center} 
 Vemos que ciertas transformaciónes que podemos hacerle a la configuración son transformaciones de simetría, pero partíamos del triángulo. La idea clave aquí es darnos cuenta de que podíamos haberlo hecho del revés, tomar las transformaciones de simetría del triángulo y darnos cuenta de que la única estructura simétrica que las tienen es el triángulo. Las dos implican la una a la otra.  Es decir, son equivalentes
 $$\text{ESTRUCTURA}\longleftrightarrow \text{TRANSFORMACIÓN}$$
Un ejemplo de la equivalencia de derecha a izquierda es una tranformación de Lorentz. Lo explico rápido y mal aquí pues entraremos en las transformaciones de Lorentz más adelante. Por ahora es suficiente con saber que son transformaciones entre sistemas de referencias inerciales en relatividad especial. Esto es, transformaciones entre sistemas de coordenadas que se desplacen a velocidad constante teniendo en cuenta el tiempo. Esto significa, que si tenemos a un observador describiendo un cierto sistema este tendrá un sistema de referencia dado, si ahora pasamos a otro sistema de coordenadas, que será otro observador que, por ejemplo, se esté desplazando, este tiene que estar estudiando exactamente el mismo sistema. O sea, esta transformación de coordenadas es una simetría\footnote{Puede confundirse esta idea con la idea de invariante. En este caso, hablamos de invariancia Lorentz, pero la simetría es mucho más, también nos permite, por el teorema de Noether, hablar de conservaciones. } de los objetos que viven en el espacio tiempo. Este ejemplo lo estudiaremos con mucho más detalle más adelante. La cuestión es que las transformaciones de Lorentz, de cierta forma, están definiendo los posibles objetos que pueden estar en el espacio tiempo, por lo que la equivalencia anterior en este caso, va de derecha a izquierda. \\ 
Este es el motivo principal por el cual la simetría es tan importante en física. Imponiendolas a nuestros sistemas encontraremos sus soluciones (como objetos que deben respetar estas simetrías). Por ejemplo en electromagnetismo para no estar hasta navidad calculando integrales solemos argumentar por simetria la relación funcional de los campos con variables espaciales o la existencia o no de ciertas componentes del campo, en este sentido, estamos partiendo de una solución general, una estructura que respeta cierta simetría. Esta idea es el día a día de los físicos, aunque normalmente solo con el fin de simplificar resultados, ha resultado ser la pieza clave para entender la física moderna. Veremos todo esto en algún momento también, sin prisa por ahora. Ahora bien, para poder 'usar' las transformaciones como cosas en sí mismas deben, ciertamente, ser una cosa por sí misma. Y no nos vale solo un conjunto de transformaciones y ya, además deben de presentar algún álgebra, que se pueda operar entre ellas y establecer ciertas reglas. Resulta qué ¡lo son! y presentan una estructura que llamamos grupo.\\\\
% No seré yo el que discuta los problemas filosóficos que esta noción pueda llegar a acarrear, esto es, los problemas filosóficos que collevan subyugar aspectos de la realidad a una realidad abstracta como es la simetría. No se deja, por tanto, de defender una especie de platonismo, es el que ciertos aspectos pueden ser deducidos a priori a partir de esa realidad abstracta y superior que se proyecta sobre ciertos aspectos concretos.  De todos modos, se puede solucionar muy facilmente cuando recordamos que trabajamos con modelos. Cualquier parecido con la realidad es mera causalidad! 
Lo que me queda de hilo es presentaros a estos amigos y luego os demostraré que efectivamente describen simetrías. \\ 
\begin{defi}
    (Grupo). Sea $G$ un conjunto y $\star$ una aplicación\\
$$\star: G\times G\rightarrow G$$
$$(x,y)\mapsto x\star y$$
tal que\\
i) $a\cdot b\in G, \ \forall a,b\in G$ ($\star$ es cerrada en $G$).\\
ii) $\star$ es asociativa en $G$.\\
iii) $\exists e\in G| e\star x=x\star e =x, \forall x\in G$ ($\star$ tiene un elemento neutro en $G$)\\
iv) $\forall x \in G, \exists x^{-1}\in G | x\star x^{-1}= e$($\star$ tiene un elemento inverso en $G \ \forall x\in G$)\\
Llamamos \textbf{Grupo} al par $(G,\star)$ \\
\end{defi}
\noindent Existe una condición extra que \textbf{NO} es necesaria para dotar al par $(G,\star)$ de estructura de Grupo que es la conmutatividad, es decir, $x\star y= y\star x, \ \forall x,y \in G$ En este caso decimos que el par $(G,\star)$ es un \textbf{Grupo Abeliano}. \\\\
Dado un grupo, también tenemos la definición de un subgrupo de la forma más naive posible. Si $(G,\star)$ es un grupo y $H\subseteq G$ entonces $(H,\star)$ es un subgrupo de $(G, \star)$ si cumple los axiomas de grupo y escribimos $H<G$.\\\\
Demostraremos a continuación algunas propiedades básicas de los Grupos.\\
\begin{prop}
    (Unicidad del elemento neutro) Sea $G$ un grupo. Existe un único elemento $e\in G$ que sea neutro.\\
\end{prop}
\textit{Demostración}. Supongamos dos elementros neutros $e,e'$ Por ser elementos neutros se cumple 
$$e=e'e$$
$$e'e=e'$$
Juntando ambos $e=e'e=e'$ por tanto $e'=e$ Luego si hay dos elementos neutros, en realidad son el mismo, por tanto el elemento neutro es único.\\\\
\begin{lem}
    (El inverso es el mismo tanto por izquierda como por la derecha) Sea $G$ grupo y $x,y\in G | xy=e$, entonces $yx=e$.
\end{lem}

\textit{Demostración}. Demostrar esto es equivalente a probar que $(x^{-1})^{-1}=x$. 
$$x=xe=x(x^{-1}(x^{-1})^{-1})=(xx^{-1})(x^{-1})^{-1}=e(x^{-1})^{-1}=(x^{-1})^{-1}$$
\begin{prop}
    (Unicidad del elemento inverso) Sea $G$ un grupo. Entonces para un $x\in G$ existe un único elemento inverso $x^{-1}\in G$.
\end{prop} 
\textit{Demostración.} Por definición de inverso es una demostración sencilla. Sean $x^{-1},y\in G |xx^{-1}=xy=e$
$$y=ye=yxx^{-1}=ex^{-1}=x^{-1}$$
No hace falta hacer la demostración para la izquierda por el Lema que acabamos de demostrar, por lo que queda demostrado.\\\\
La proposición que hacemos es que, verdaderamente, un grupo es un objeto matemático que nos permite estudiar transformaciones de simetría, esto es, siempre que tengamos un objeto que presente una cierta simetría, el conjunto de transformaciones que respeten esta simetría conformarán un grupo. Vamos a ver por qué y luego vamos a hablar sobre las implicaciones que esto puede llegar a tener: \\\\
Para ver esto, vamos a hablar de un tipo de grupo concreto, el grupo de permutaciones. Primero, una permutación es, dado un conjunto de $n$ elementos, una biyección del conjunto sobre sí mismo, es decir, reordenaciones de los elementos del conjunto. Vamos a ver esto con algo más detalle:\\
Una biyección es una aplicación inyectiva (a dos elementos distintos del dominio le corresponden dos elementos distintos en la imagen) y sobreyectiva (todos los elementos de la imagen son imagen de algún elemento del dominio) Por tanto una biyección nos establece una relación 1 a 1 entre todos los elementos del dominio y la imagen. Las aplicaciones biyectivas, por tanto, tienen inversa. Si tenemos dos conjuntos $A$ y $B$ y podemos encontrar una aplicación biyectiva entre ellos, decimos que son isomorfos. \\\\
Esto no pasa, claro, en el caso de los grupos. En el caso de los grupos, tenemos estructura extra, que son las operaciones de grupo con sus cuatro propiedades, y por tanto, necesitamos algo que nos preserve esta estructura. La aplicación que preserva la estructura de grupo es el homomorfismo. 
\begin{defi}
    Sean $(G, \star)$, $(H, \circ)$ dos grupos y $\phi: G\longrightarrow H$ una aplicación entre ellos, decimos que es un homomorfismo si 
    $$\phi(g\star h) = \phi (g)\circ \phi (h) $$
\end{defi}
Juntando ambas cosas, diremos que dos grupos son isomorfos si existe un homomorfismo biyectivo, que llamamos isomorfismo, entre ellos. Esto será importante para luego. \\\\
Como decíamos, una permutación es una aplicación biyectiva sobre el mismo conjunto, y por lo que acabamos que ver, se puede ver que esto se refiere a las distintas posibles ordenaciones de los elementos de un conjunto. \\\\ 
% DIBUIX
Ahora, sea $X$ un conjunto, se puede definir el conjunto de biyecciones sobre el conjunto
\begin{equation}
    Sym(X)=\{f: X\rightarrow X| f \text{ biyec.}\}
\end{equation}
Dado este conjunto, se puede demostrar que es grupo junto con la operación composición de aplicaciones. A este grupo, le llamamos grupo simétrico o grupo de permutaciones de $X$ y sus elementos, claro, son permutaciones de $X$. \\
Si el conjunto es de la forma $\{1,2,3,\dots, n\}$ su grupo simétrico se escribe $S_n$. Este será, por tanto, el grupo de las permutaciones de $n$ elementos. Si recordáis del cole, tenemos $n!$ formas de permutar $n$ elementos, por tanto, este grupo tiene $n!$ transformaciones. \\\\ Ahora, tenemos un teorema super importante, y super bonito: El llamado teorema de Cayley. Dice así:
\begin{teo}
    Todo grupo es isomorfo a un subgrupo de un grupo simétrico. 
\end{teo}
El grupo simétrico tiene en su definición una estructura que está globalmente mantenida, o sea, el conjunto, si este conjunto representa 'algo', y aquí está la clave de la cuestión, claro, el grupo simétrico no es más que el grupo de transformaciones de simetría de ese 'algo', porque el grupo se preserva.  % Además, como también va para subgrupos del grupo simétrico, el conjunto puede tener ciertas restricciones que el teorema seguirá dandonos un grupo. Por lo que el teorema de Cayley nos da lo que decíamos anteriormente, que los grupos representan simetrías.  \\\\
\\\\ Vale. Puede que no me haya explicado muy bien. Vamos a verlo con un par de ejemplos. Vamos a visitar a un amigable amigo: el triángulo rectángulo, al que le marcamos sus vértices
\begin{center}
    \begin{tikzpicture}
        \draw[ultra thick, black](-2,0)--(2,0)--(0,3.5)--cycle;
        \node at (2,0)[below right]{$1$};
        \node at (0,3.5)[above]{$2$};
        \node at (-2,0)[below left]{$3$};
        
    \end{tikzpicture}
\end{center} 
Podemos hacernos la pregunta, ¿qué transformaciones deján igual geométricamente el triángulo? Es evidente que las rotaciones de $120^\circ$ grados lo dejan igual. Si giramos $120^\circ$ en sentido contrario de las agujas del reloj tendremos 
\begin{center}
    \begin{tikzpicture}
        \draw[ultra thick, black](-2,0)--(2,0)--(0,3.5)--cycle;
        \node at (2,0)[below right]{$3$};
        \node at (0,3.5)[above]{$1$};
        \node at (-2,0)[below left]{$2$};
        
    \end{tikzpicture}
\end{center} 
O si giramos el primero en sentido horario tendremos 
\begin{center}
    \begin{tikzpicture}
        \draw[ultra thick, black](-2,0)--(2,0)--(0,3.5)--cycle;
        \node at (2,0)[below right]{$2$};
        \node at (0,3.5)[above]{$3$};
        \node at (-2,0)[below left]{$1$};
        
    \end{tikzpicture}
\end{center} 
Llamamaremos al giro en sentido antihorario $\sigma_L$ y al giro en sentido horario $\sigma_R$, por darles nombre. Pero hay otras transformaciones no tan tan evidentes (siguen siendolo). Estas son las reflexiones. Si tomamos como eje las medianas del triángulo (el segmento que va de un vértice a la mitad del lado opuesto) tendremos $3$ transformaciones más. 
\begin{center}
    \begin{tikzpicture}
        \draw[ultra thick, black](-2,0)--(2,0)--(0,3.5)--cycle;
        \node at (2,0)[below right]{$1$};
        \node at (0,3.5)[above]{$2$};
        \node at (-2,0)[below left]{$3$};
        \draw[dotted, black](-2,0)--(1,1.7);
         \draw[dotted, black](2,0)--(-1,1.7);
                 \draw[dotted, black](0,3.5)--(0,0);
    \end{tikzpicture}
\end{center} 
Que son las siguientes transformaciones, que llamaré $\tau_i$ donde $i$ dependerá del vértice del que salga la mediana. 
\begin{center}
    \begin{tikzpicture}
        \draw[ultra thick, black](-2,0)--(2,0)--(0,3.5)--cycle;
        \node at (2,0)[below right]{$2$};
        \node at (0,3.5)[above]{$1$};
        \node at (-2,0)[below left]{$3$};
        \draw[dotted, black](-2,0)--(1,1.7);
    \end{tikzpicture}
\end{center} 
\begin{center}
    \begin{tikzpicture}
        \draw[ultra thick, black](-2,0)--(2,0)--(0,3.5)--cycle;
        \node at (2,0)[below right]{$1$};
        \node at (0,3.5)[above]{$3$};
        \node at (-2,0)[below left]{$2$};
         \draw[dotted, black](2,0)--(-1,1.7);
    \end{tikzpicture}
\end{center} 
\begin{center}
    \begin{tikzpicture}
        \draw[ultra thick, black](-2,0)--(2,0)--(0,3.5)--cycle;
        \node at (2,0)[below right]{$3$};
        \node at (0,3.5)[above]{$2$};
        \node at (-2,0)[below left]{$1$};;
                 \draw[dotted, black](0,3.5)--(0,0);
    \end{tikzpicture}
\end{center} 
Además hay una transformación más que podemos hacer al triángulo y es, exacto, no hacerle nada, llamamos a esta transformación identidad, y escribimos $e$. Decimos ahora que es un grupo con la composición de transformaciones y que será el grupo de simetría del triángulo. 
$$G_{\Delta}=\{e,\tau_1,\tau_2,\tau_3,\sigma_1,\sigma_2\}$$
Habría que demostrar que esto efectivamente es un grupo, pero es sencillo. Sobretodo nos podemos dar cuenta de que este grupo resulta ser el grupo de permutaciones de un conjunto de $3$ elementos, o, de otra forma, el grupo simétrico del conjunto $\{1,2,3\}$. Podéis comprobar que, las anteriores transformaciones son equivalentes a las $6$ (recordad, veíamos, $S_n$ tenía $n!$ elementos) de $S_3$. Aquí aparece el teorema de Cayley, pero quizás en este ejemplo no se ve tanto la gracia. Vamos a ver otro. Otra figurita amigable: el cuadrado!
\begin{center}
    \begin{tikzpicture}
        \draw[ultra thick, black](0,0)--(4,0)--(4,4)--(0,4)--cycle;
        \node at (4,0)[below right]{$1$};
        \node at (4,4)[above right]{$2$};
        \node at (0,4)[above left]{$3$};
        \node at (0,0)[below left]{$4$};
    \end{tikzpicture}
\end{center} 
De nuevo, vamos a encontrar el grupo de simetría del cuadrado. Impulsados por lo visto en el triángulo podemos pensar que es el grupo simétrico de $\{1,2,3,4\}$, $S_4$, pero este no es el caso. El motivo es porque no hay reflexiones de dos elementos que dejen dos elementos contiguos inmóviles. O sea, tomad el cuadrado con la disposición de índices de arriba, si permutamos $3\leftrightarrow 2$ entonces necesariamente debemos permutar $4\leftrightarrow 1$ porque, ojo, aquí está la clave, la propia estructura del cuadrado nos obliga, así que, así de primeras, podemos decir que, seguro $S_4$ no es el grupo de simetría del cuadrado. Podemos ponernos a calcular los elementos por observaciones similares a las del triángulo los elementos concretos que serán transformaciones de simetría del cuadrado. Este grupo, se llama grupo diédrico $4$ y escribimos $D_4$ y será, claro, un subgrupo de $S_4$. Aquí encontramos otro ejemplo donde entra el teorema de Cayley, concretamente, que es grupo de simetrías del cuadrado es isomorfo a $D_4 < S_4$ \\\\
Decía al principio que iba a tratar mostraros como estas estructuras que hemos llamado grupo representan simetrías. Todo está en el teorema de Cayley. Vamos a volver al cuadrado y vamos a hacerlo poco a poco. Partimos del cuadrado 
\begin{center}
    \begin{tikzpicture}
        \draw[ultra thick, black](0,0)--(4,0)--(4,4)--(0,4)--cycle;
    \end{tikzpicture}
\end{center} 
Lo que hacíamos ahora es nombrar los vértices 
\begin{center}
    \begin{tikzpicture}
        \draw[ultra thick, black](0,0)--(4,0)--(4,4)--(0,4)--cycle;
        \node at (4,0)[below right]{$1$};
        \node at (4,4)[above right]{$2$};
        \node at (0,4)[above left]{$3$};
        \node at (0,0)[below left]{$4$};
    \end{tikzpicture}
\end{center} 
Esto nos permitía representar los vértices en términos de algo así como un conjunto subyacente $\{1,2,3,4\}$ que sabemos tiene un grupo simétrico, $S_4$. Sin embargo, este no es el grupo de simetría del cuadrado ¿por qué? porque $S_4$ no incluye la estructura del cuadrado, sin embargo, $D_4$, un subgrupo de $S_4$ sí la inclye, de cierta forma $D_4$ codifica la estructura del cuadrado, su simetría. Por Cayley, todo grupo es grupo de permutaciones pero un grupo de permutaciones preserva una cierta estructura por tanto todo grupo preserva una cierta estructura.
% Nuestro querídisimo teorema de Cayley para nosotros es que todo grupo (y por tanto, de transformaciones de simetría) es isomorfo a un subgrupo del grupo simétrico, pero debido a que en la idea de que sea subgrupo está codificada la propia estructura que hace que algo presente una simetría o no, entonces, los grupos. por tanto, atienden a la estructura de los objetos (de simetría) y contienen para si las transformaciones que respetan tal estructura.  
\\\\ Espero con esto, de verdad lo espero, haberos convencido, por tanto, de la bonita idea de que los grupos representan simetrías. 
\\\\
Intentemos volver al primer hilo, para ello, vamos a estudiar el círculo 

\begin{center}
\begin{tikzpicture}
\filldraw[color=black, fill=white, very thick](-2,0) circle (2);
\end{tikzpicture}
\end{center} 
 El círculo tiene infinitas transformaciones distintas que lo dejan igual. Cada diámetro es un eje de simetría y cada ángulo es una posible rotación. Si definimos un ángulo inicial, todas las rotaciones de ángulo $\theta$ serán transformaciones de simetrías del círculo.\\\\
En el caso del círculo, como veíamos, siempre que tengamos un ángulo $\theta$ un ángulo inifinitamente pequeño $\varepsilon$  puede sumarsele y será otro ángulo $\theta +\varepsilon \in S^1$. Esto lo hablábamos en el hilo anterior y era debido a la continuidad del círculo. El conjunto de transformaciones de simetría del círculo forman un conjunto infinito, además, dado que para cada punto del círculo existe una transformación de simetría, no es muy descabellado decir que el conjunto de transformaciones de simetría del círculo es, efectivamente, el propio círculo. \\\\
Este conjunto, de la misma forma que antes, conforma un grupo con la composición de transformaciones. Fijaos lo que tenemos pues, tenemos un espacio (el círculo) que es a la vez un grupo! Esto es lo que llamamos un grupo de Lie. Y aquí lo dejamos, en el próximo hilo vamos a ver por fin que es esto de los grupos de Lie. 
\section{Los grupos de Lie no son tan solo grupos de infinitos elementos}

Existe una diferencia fundamental entre un grupo de infinitos elementos como podría ser el grupo de permutaciones de infinitos elementos y el grupo de simetría del círculo. ¿La diferencia? La idea de continuidad, es decir, las transformaciones de simetría del círculo podemos entenderlas como una función continua en función del parámetro que es el ángulo. Vamos a repasar la definición de función continua:\\
Una función $f:\mathbb{R}\rightarrow\mathbb{R}$ es continua en un punto $x_0$ si $$\forall \varepsilon >0, \exists \delta >0 \text{ tal que }|x-x_0|<\delta \Rightarrow |f(x)-f(x_0)|<\varepsilon$$. 
Esta definición está muy bien, pero ocurre que nosotros no tenemos una $f:\mathbb{R}\rightarrow\mathbb{R}$ tenemos una asignación $\theta\mapsto T_\theta$ donde $\theta\in[0,2\pi]$ recorre el círculo y $T_\theta$ son las transformaciones de simetría del círculo. Esto nos va a obligar a extender nuestra definición de continuidad \\\\
Para ello, necesitamos una teoría nueva, esta teoría se llama Topología. Hoy no vamos a ver nada de topología como tal, pero sí que vamos a dar su idea básica. La topología, en términos muy simples, es una teoría que nos permite generalizar el concepto de distancia. Por tanto, dos espacios topológicos serán iguales o mejor dicho, homeomorfos, si puede transformarse el uno en otro por medio de transformaciones continuas. No se preocupa por la forma de los objetos, sino por, por ejemplo, como se unen, o sus agujeros. El ejemplo más típico es el donut y la taza. Esto objetos son evidentemente distintos geométricamente, pero topológicamente, son el mismo espacio porque yo de la taza puedo llegar al donut y viceversa. \\\\
Volvamos al círculo, el léctor avispado se habrá dado cuenta de que, realmente $SO(2)$ no es el grupo de simetría del círculo, el grupo de simetría del círculo es cualquier transformación ortogonal $O(2)$, esto es porque los elementos de determinante $-1$ son reflexiones, por definción. Pero, como decíamos, el el subconjunto de $O(2)$ con determinante $-1$ no podía formar un grupo, porque no tendría la identidad. Sin embargo, $O(2)$ sí es un grupo pues, como espacio topológico, tiene dos componentes conexas, una que tiene la identidad y otra que no. Esta idea de "conectidad" merece un poco más de concrección: \\
Diremos que un espacio es conexo si puede ponerse de una sola pieza, simplemente. Pero dentro de ser conexo un espacio puede ser o conexo por caminos o simplemente conexo. \\
Diremos que es conexo por caminos si dos puntos cualquiera del espacio pueden unirse por un camino. Por otro lado, diremos que un espacio es simplemente conexo si no tiene agujeros. Esto significa que todo bucle que defina sobre el espacio puede contraerse a un punto. Estas son ambas ideas muy intuitivas que podéis ver en el dibujo. \\\\

Nosotros estamos interesados en un tipo muy concreto de espacio topológico, una variedad. Variedad no es más que la palabra fancy para decir hipersuperficie e hipersuperficie no es más que la palabra fancy para decir superficie de más dimensiones. Ahora, ¿qué caracteriza una superficie? Pues para nosotros una $n$- variedad será un espacio topológico que es localmente homeomorfo (recordad, isomorfos como espacios topológicos) a $\mathbb{R}^n$. Esto significa, que en la variedad siempre que yo coja un abierto $U$, existirá una aplicación $x$ tal que $x(U)\subseteq \mathbb{R}$. En la imagen tenéis un ejemplo.  \\
Volvamos a nuestro queridísimo círculo. Ahora, podemos coger cualquier abierto del círculo y va a existir una aplicación que nos lo manda a algún subconjunto de $\mathbb{R}$. Si llenamos la variedad de estos parches, de manera que puedan definirse unas coordenadas en ella y tal que las funciones definidas en esta variedad sean infinitamente diferenciables, decimos que es una variedad suave. Esto parece muy complicado pero no lo es para nada. \\ 

El hecho de que estemos pasando a $\mathbb{R}^n$ nos permite derivar. El término suave simplemente se refiere a que toda función definida sobre la variedad sea infinitamente derivable.\\\\

Un ejemplo de un uso de una variedad sin darnos cuenta físico es el llamado espacio de configuraciones. Para quien no lo sepa, un espacio de configuraciones es el espacio que conforman las infinitas posibles posiciones de un sistema físico. Este espacio tiene una estructura natural de variedad suave, porque las posiciones de un sistema físico, clásico, he de decir, cambian en el tiempo suavemente. Cuando hacemos teoría aquí, básicamente mecánica Lagrangiana definimos en general las llamadas 'coordenadas generalizadas' que usualmente se escriben $q^i$. Estas son unas coordenadas abstractas que no sirven para hacer cuentas, pero sí para desarrollar toda la teoría. Estos serían los propios 'puntos' de la variedad.  A la teoría no le importa las coordenadas que uses, lo cual está bastante bien como teoría. Esto se llama invariancia bajo difeomorfismos (y es por cierto uno de esos pequeños problemitas que nos impide cuantizar la gravedad xd) Cuando vamos a resolver un problema, nombramos estas coordenadas, por ejemplo, para el péndulo simple $q^1=x, q^2=y$ esto es equivalente a coger la variedad y ponerle un atlas, es decir, hacer un mapa de ella, con estas coordenadas. Pero también podríamos haber hecho otro mapa, otras coordenadas, $q^1=\theta, q^2=r$ pero como $r$ en el péndulo simple es constante, entonces $r=cte.$ o dicho de otra manera, NO NOS IMPORTA el radio ¿qué espacio es característico de esto? exacto, el círculo. El espacio de configuraciones del péndulo simple es el círculo, $S^1$. Este atlas en el que todo el círculo se puede caracterizar por un único angulo $\theta$ es el mismo que usábamos cuando hablábamos del círculo como grupo de simetría. Impulsados por esto, nos atrevemos a decir:
\\\\
Un grupo de transformaciones continuas, o de Lie, es una variedad suave a la que además le damos la estructura de grupo, esto es, una variedad con una aplicación binaria suave.
Además por motivos de suavidad, vamos a pedir que la aplicación que mapea cada elemento con su inverso sea suave.
\\\\
Bueno, si tomamos esta definición, nuestro resultado del círculo tenía todo el sentido del mundo, como veíamos antes, pero es que hay más, porque las variedades suaves tienen una característica muy  importante. Como decíamos antes, el término 'suave' se refería a que las funciones definidas sobre la variedad eran infinitamente diferenciables. Esto podemos aprovechar para construir el llamado espacio tangente a la variedad en un punto. \\\\
Dado un punto, pasan por él inifinitas curvas suaves distintas, el espacio de las velocidades de estas infinitas curvas es un espacio vectorial que llamamos espacio vectorial tangente y escribimos $T_p M$. Este espacio os sonará a algunos, pues es en el espacio tangente el espacio en el que hacemos las cuentas en relatividad. Usando todos los espacios tangentes y como se conectan entre ellos podemos recuperar la idea de geometría, pero eso para otro momento. Recordad cuando estudiábamos $SO(2)$ que derivábamos y evaluábamos en $\theta =0$ esto es equivalente a hacer el espacio tangente en la identidad!! $T_e M$ así, tomábamos un elemento cualquiera de $SO(2)$ y derivábamos en la identidad
\begin{equation*}
    J=R' (0) = \left. \frac{d}{d\theta}\begin{pmatrix}
        \cos \theta & -\sin\theta \\ \sin\theta & \cos\theta 
    \end{pmatrix}\right|_{\theta =0 } =\left. \begin{pmatrix}
        -\sin\theta & -\cos \theta \\ \cos \theta &-\sin\theta 
    \end{pmatrix} \right|_{\theta=0} =\begin{pmatrix}
        0 & -1 \\ 1 & 0
    \end{pmatrix}
\end{equation*}
Lo que estamos consiguiendo aquí, entonces, es el vector base del espacio vectorial tangente a la variedad en la identidad. Así, cualquier elemento de este espacio vectorial se escribe como $\theta J$ con $\theta$ es un real. Este es el motivo por el cual el hilo anterior lo acabamos diciendo que el dibujo
\begin{center}
\begin{tikzpicture}
\draw[very thick, black](2,-3)--(2,3);
\filldraw[color=black, fill=white, very thick](0,0) circle (2);
      \filldraw[color=black, fill=black, very thick](2,0) circle (0.05);
        \node at (2,0) (A) [below left]{$\theta = 0$};
\end{tikzpicture}
\end{center} 
representaba todo lo que acabábamos de ver. Nos falta por ver una cosa, la matriz exponencial. Fijaos que mientras que $SO(2)$ son las matrices ortogonales de determinante unidad, su álgebra está generada por $J$ y por tanto serán todas matrices antisimétricas de traza nula. Tenemos el resultado elemental de teoría espectral que dice que
\begin{equation}
    \det (e^A)=e^{\text{tr}A}
\end{equation}
Aunque es general, para verlo en el caso más sencillo, asumid que $A$ es diagonalizable y que $A$ está en su forma diagonal, entonces su traza es la suma de valores propios y  $e^A$ será una matriz cuya diagonal sean exponenciales de vectores propios, entonces $\det (e^A) =e^{m_1}e^{m_2}\dots=e^{m_1+m_2+\dots} =e^{\text{tr}A}$. No nos importa la demostración de esto ahora, nos importa el resultado. Esto es debido a que si $\det (e^A)=1$ entonces $\text{tr}A=0$ necesariamente. Así que... si $A$ perteneciera a $T_eG$ entonces $e^A$ parece pertenecer a $G$! Esta es la intuición, pero resulta ser general. Vamos a verlo.\\
Si $X$ pertenece al algebra, entonces será la derivada de alguna curva evaluada en la identidad porque pertenece al espacio tangente a la identidad. Sea la curva en el grupo parametrízada $\gamma: \mathbb{R}\rightarrow G$ y tal que $\gamma(0)=id$  entonces, $X$ pertenece a
$\mathfrak{g}$ si $X=\gamma(0)'$. Esta ecuación es la que tendremos que resolver. Fijaos que una curva parametrizada por $t$ de la forma $\exp(tX)$ satisface la ecuación puesto que su derivada en la identidad será $Xe^{0X}=X$ Por lo que podemos decir que si $X$ pertenece al álgebra, entonces $e^X$ pertenecerá al grupo, es decir, para pasar del álgebra al grupo usamos la aplicación exponencial. En general uno puede demostrar que si $\{X_1,X_2,\dots, X_n\}$ es un conjunto de elementos del álgebra y $G$ es un grupo conexo entonces 
$e^{X_1}e^{X_2}\dots e^{X_n}$ es un elemento del grupo.\\\\ 
Bueno, vamos a poner nuestra maquinaria a jugar, en este hilo vamos a estudiar todo $SO(n)$ en un momento. Vamos a empezar con $SO(2)$, que es nuestro amigo que ya conocemos tan bien. \\\\
Vamos ahora, en uso de estos resultados, encontrar, sin geometría, la matriz de rotaciones en el plano, es decir, el grupo $SO(2)$. \\ Lo que haremos será encontrar que matrices pueden formar una base de $\mathfrak{so(2)}$ y luego, encontrar la matriz de rotación en el plano por exponenciación.\\\\
El álgebra de $SO(2)$ será su espacio tangente en la identidad, que no es más que la matriz identidad. puesto que estamos tratando un grupo de Lie matricial
Supongamos que la matriz
$$A= \begin{pmatrix}
	a & b \\
	c & d
\end{pmatrix}$$
pertenece a este álgebra . Como pertenece a este álgebra, debe existir una curva en el grupo que pase por la identidad y tal que la matriz que buscamos sea tangente en la identidad.  \\
Podemos parametrizar esta curva $\gamma$ 
$$\gamma:\mathbb{R}\rightarrow SO(2)$$
Además debe cumplirse que en el 0 la curva es la identidad de $SO(2)$ 
$$\gamma(0)=\mathbb{I}$$
La curva esta debera ser claramente una matriz 

$$\gamma(\lambda)= \begin{pmatrix}
	x(\lambda) & y(\lambda) \\
	z(\lambda) & w(\lambda)
\end{pmatrix}$$
Esto es, porque quede clarísimo. Esto es una curva definida sobre el grupo $SO(2)$. 
Recordemos que condiciones debía tener una matriz de rotación.\\ 
Dimos una definición algebraica para $SO(2)$, que era:  
\begin{equation}
    SO(2) =\{R\in M_2(\mathbb{R}) | R R^T=I, \text{det}R =1\}
\end{equation}
Pasemos ahora a formalizar estas dos condiciones:\\
La ortogonalidad se expresa en términos de la curva de arriba como
\begin{equation}
    \begin{pmatrix}
        x(\lambda) & y(\lambda) \\ z(\lambda) & w(\lambda) 
    \end{pmatrix}
        \begin{pmatrix}
        x(\lambda) & z(\lambda)  \\ y(\lambda) & w(\lambda) 
    \end{pmatrix} =\begin{pmatrix}
        1 & 0 \\ 0 & 1 
    \end{pmatrix}
\end{equation}
Realizando el producto matricial y comparando términos se traduce en las condiciones 
$$x^2+y^2=1 $$
$$xz+yw=0$$
$$z^2+w^2=1$$
Por otro lado, determinante unidad en este caso es evidentemente la condición
$$xw-yz=1$$
Además como para $\lambda=0$, $\gamma(t)$ pasa por la identidad y, por tanto, esa matriz debe ser la identidad 
$$x(0)=w(0)=1, y(0)=z(0)=0$$
Esto era, repito, una curva definida sobre el grupo. En términos del álgebra tenemos que
$$\gamma'(0)= \begin{pmatrix}
	x'(0) & y'(0) \\
	z'(0) & w'(0)
\end{pmatrix}= \begin{pmatrix}
	a & b \\
	c & d
\end{pmatrix}$$
porque el álgebra es el espacio tangente en la identidad. \\\\
Con todo esto podemos encontrar condiciones sobre $a, b, c, d$. \\\\
Si derivamos la condición de determinante unidad y evaluamos en la identidad 
$$\left.\frac{d}{d\lambda}\right|_{0}[xw-yz]=0 $$
Por simple regla del producto
$$x'(0)w(0)+x(0)w'(0)-y'(0)z(0)-y(0)z'(0)=0$$
sustituyendo los valores que sabemos obtenemos que 
$$x'(0)+w'(0)=0$$ y por tanto 
$$a=-d $$
Ahora derivamos las ecuaciones que encontramos por ser ortogonales
$$x'(0)z(0)+x(0)z'(0)+ y'(0)w(0)+y(0)w'(0)=0 $$
obtenemos que 
$$z'(0)+y'(0)=0$$
$$b=-c $$
finalmente de las otras dos condiciones si derivamos en la identidad 
$$2x(0)x'(0)=0 $$
$$2w(0)w'(0)=0$$
$$x'(0)=w'(0)=0 $$
$$a=d=0 $$
por tanto las matrices que pertenecen a $\mathfrak{so(2)}$ son las matrices 
$$A= \begin{pmatrix}
	0 & b \\
	c & 0
\end{pmatrix}$$ tales que$ b=-c $ o sea, las matrices de la forma
$$A= \begin{pmatrix}
	0 & -c \\
	c & 0
\end{pmatrix}$$
es decir que están generadas por  
$$J= \begin{pmatrix}
	0 & -1 \\
	1 & 0
\end{pmatrix}$$
Que es el elemento de la base del álgebra que ya habíamos encontrado antes, haciendo el trabajo inverso, yendo del grupo al álgebra. Lo increíble es que lo que acabamos de hacer es totalmente general. 
\\
Podemos pasar ahora al caso importante: rotaciones en tres dimensiones. Las codiciones son las mismas pero en dos dimensiones. El determinante unidad nos da la condición:
\begin{equation}
	\varepsilon_{ijk}a_{1}^ia_{2}^ja_{3}^k=1
\end{equation}
Donde estamos sumando para cada índice repetido arriba y abajo.
Si derivamos
\begin{equation}
	\varepsilon_{ijk}\left[(a_2^ja_3^k)x_1^i+(a_1^ia_3^k)x_2^j+(a_1^ia_2^j)x_3^k\right]=0
\end{equation}
Donde el elemento del álgebra será $X$ y sus componentes $a'^i_j(0)=x^i_j$
Ahora, como en la identidad debe ser, la identidad, $a^i_j(0)=\delta^i_j$ por tanto
\begin{equation}
	\varepsilon_{ijk}\left[(\delta_2^j\delta_3^k)x_1^i+(\delta_1^i\delta_3^k)x_2^j+(\delta_1^i\delta_2^j)x_3^k\right]=0
\end{equation}
Aplicamos las deltas y nos queda
\begin{equation}
	\varepsilon_{i23}x_1^i+\varepsilon_{1j3}x_2^j+\varepsilon_{12k}x_3^k=0
\end{equation}
Como la epsilon con índices repetidos es nula solo nos sobrevive.
\begin{equation}
	x_1^1+x_2^2+x_3^3=0
\end{equation}
Es decir, $\text{tr}(X)=0$ las matrices serán de traza nula como pasaba en el caso dos dimensional. \\
La condición de ortonormalidad se traduce en
\begin{equation}
	a^i_j(a^T)^j_k=\delta^i_k
\end{equation}
Si derivamos y en uso de la misma notación de antes
\begin{equation}
	(a^T)^j_k(0)x^i_j+a^i_j(0)(x^T)^j_k=\delta^k_jx^i_j+\delta^i_jx^k_j=0
\end{equation}
Actuando las deltas
\begin{equation}
	x^i_k+x^k_i=0
\end{equation}
Encontrando la condición
\begin{equation}
	x^i_k=-x^k_i
\end{equation}
Es decir, son matrices antisimétricas, de nuevo, como veíamos en $SO(2)$ y $so(2)$. \\
Con esto tenemos todas las condiciones para encontrar una base del álgebra. Para $\mathfrak{s}\mathfrak{o}(3)$ Es claro que una posible base será
\begin{equation}
	L_1=\begin{pmatrix}
		0 & -1 & 0 \\ 1 & 0 & 0\\ 0&0&0
	\end{pmatrix}
\end{equation}
\begin{equation}
	L_2=\begin{pmatrix}
		0 & 0 & 1 \\ 0 & 0 & 0\\ -1&0&0
	\end{pmatrix}
\end{equation}
\begin{equation}
	L_3=\begin{pmatrix}
		0 & 0 & 0 \\ 0 & 0 & -1\\ 0&1&0
	\end{pmatrix}
\end{equation}
De nuevo, tan solo tendríamos que tomar la exponencial para obtener una rotación. Ahora tenemos tres generadores por lo que le corresponde una exponencial a cada uno. Así, un elemento general de $SO(3)$ se escribe
\begin{equation}
    R=e^{\theta L_1}e^{\varphi L_2}e^{\psi L_3}
\end{equation}
Donde $\theta, \varphi, \psi \in [0,2\pi]$ Bueno y cual es el sentido de esto. Como esto ya es larguísimo vamos a estirarnos más y a calcular las exponenciales. La primera será
$e^{\theta L_1}$ que es
\begin{equation}
    \exp {\theta \begin{pmatrix}
		0 & -1 & 0 \\ 1 & 0 & 0\\ 0&0&0
	\end{pmatrix}}
\end{equation}
\\\\
Podemos ahora calcular el conmutador. Acóstumbra a definirse $K_i=iL_i$
\begin{equation}
	[K_i,K_j]=K_iK_j-K_jK_i
\end{equation}
Es un buen ejercicio comprobar que obtenemos
\begin{equation}
	[K_i,K_j]=i\epsilon_{ijk}K_k
\end{equation}
Que es el álgebra de rotaciones, $\mathfrak{so(3)}$. Con este factor $i$, la aplicación exponencial se define con un factor $-i$.
Lo hecho aquí es generalizable a cualquier dimensión! Para dimensión $N$
\begin{equation}
	\sum_i x^i_i=0, \ \ \ i=1,2,3,...,N
\end{equation}
\begin{equation}
	x^i_k=-x^k_i, \ \ \  i=1,2,3,...,N
\end{equation}
En general podemos decir que $\mathfrak{s}\mathfrak{o}(n)$ es el álgebra de matrices antisimétricas de traza nula. \\ La matriz $N\times N$ tendrá $N^2$ elementos de los cuales establecemos restricciones en $N$ de ellos (la diagonal) y en los elementos no diagonales, pero solo la mitad, esto será un número cualquiera $D$ que será el número de generadores y por tanto la dimensión del álgebra. No es muy dificil ver que se cumple que
\begin{equation}
	N^2-N-D=D
\end{equation}
De donde despejamos 
\begin{equation}
	D=\text{dim} (\mathfrak{so(n)})=\frac{N(N-1)}{2}
\end{equation}
(Este resultado suele formalmente demostrarse por inducción.)\\
Y así, tan fácil, hemos encontrado como funcionan las rotaciones en cualquier dimensión finita. Fijaos que esto es una generalización para cualquier dimensión del famoso teorema sobre rotaciones que nos dice que, en tres dimensiones, podemos descomponer cualquier rotación en tres. Ahora podemos decir:
\\
En un espacio $n-$dimensional, cualquier rotación se puede descomponer en $n(n-1)/2$ rotaciones.\\
Como $\mathfrak{s}\mathfrak{o}(n)$ tiene la misma dimensión que $SO(n)$ hemos encontrado la dimensión de la variedad. 
No tiene por qué sorprendernos este resultado, pero sí que es un resultado bonito y sorprendente, y, además, lo hemos sacado de una forma muy sencilla (nos falta encontrar que la ecuación 18 es real, pero ya digo que por inducción no nos debería quitar el sueño). Esto ya es señal del poder que tiene la teoría de grupos de Lie.\\ 

\section{Representaciones}
%NUEVA FORMA
Con tanta cosa de variedad de espacio tangente de exponencial y de cosas nos hemos podido perder. No os preocupéis, no voy a haceros estudiar (aunque sí animaros a ello), iremos repasando conceptos claves cada vez. Pero lo que sí está claro es que el formalismo matemático, aun por bonito que es, ha acabado emborronando nuestra idea de simetría. Recordad que vimos en la segunda parte de la introducción a Grupos de Lie, que una simetría es un conjunto de transformaciones que dejan una cierta estructura invariante. Vimos en este hilo también, que conformaban un grupo, y que esto nos lo daba el teorema de Cayley, que decía que todo grupo era isomorfo a un grupo de permutaciones, y vimos que estas, en esencia, lo que acababan haciendo era siempre respetar una cierta estructura. La pregunta ahora es obvia ¿cómo aplicamos eso a grupos continuos? ¿Tiene sentido si quiera? \\\\
No es muy conocido, y con razón, que Bertrand Rusell antes de destruir los sueños de Frege o de perder el tiempo con Whitehead dedicándole 400 páginas a la demostración lógica de la proposición $1+1=2$, o de ser un icono pop en general,  fue un filósofo cuyo interés estaba especialmente en la geometría. En 1897 se publica 'An essay on the foundations of geometry'. En este libro Rusell hace un repaso de distintas filosofías de la geometría. Lo que nos importa a nosotros es que critica, mejor dicho, se atreve a criticar, la teoría de las variedades de Riemaann porque esta no podía estipular el espacio en el que se encontraban las variedades. Es curioso porque aquí ya aparece la gran preocupación de Rusell: meter cosas en cosas. Pero bueno. La cuestión es que todos los matemáticos que entendían bien a Riemann respondieran diciendo que precisamente la capacidad de formular la teoría sin necesidad de un espacio ambiente es una de sus virtudes. Y en este caso, esta propiedad brilla de sobremanera. \\
Pensemos en las rotaciones de nuevo. Ya hemos visto que son un circulo $S^1$ ahora ¿Cual es su espacio ambiente? Pues el que sea, siempre y cuando incluya al círculo, claro. Por ejemplo, recordad cuando os puse el ejemplo del péndulo simple, decíamos que podíamos usar las coordenadas $q^1=x, q^2=y$ esto es porque son unas coordenadas de $\mathbb{R}^2$ y luego simplemente nos restringimos al círculo dándonos la paremetrización usual $x=\cos\theta, y=\sin\theta$. En este caso, tenemos pues que los elementos del círculo son un vector de dos dimensiones $(\cos\theta, \sin\theta)$ o equivalentemente el ángulo $\theta$ como ya hemos visto. \\ Cuando un espacio está dentro de otro de esta forma, lo llamamos embediemiento o encaje. Concretamente, son aplicaciones injectivas que preservan la estructura. $S^1$ puede tener como espacio ambiental también $\mathbb{R}^3$ aunque en este caso, el tema se complica pues a diferencia de en $\mathbb{R}^2$ que solo había una forma de tener el círculo, en $\mathbb{R}^3$ hay muchísimas formas (dadas por un eje, claro). O también podríamos haber tomado $\mathbb{R}^4$ ¡Incluso peor! Sí, pero $\mathbb{R}^4$ oculta un secreto. Fijaos, los elementos de este espacio son vectores  de la forma $(a,b,c,d)$ pero podríamos reordenar un poquito esto y tendríamos 
\begin{equation*}
    \begin{pmatrix}
        a & b \\ c & d
    \end{pmatrix}
\end{equation*}
Entonces podríamos encajar el círculo en un espacio matricial. Queremos además que sea un grupo así que no puede ser una matriz cualquiera, debe ser al menos invertible y tener identidad (que el producto de matrices es asociativo y cerrado es un hecho conocido) Así que tomamos ese conjunto de matrices, que llamamos grupo general lineal
\begin{equation}
    GL(2,\mathbb{R})
\end{equation}
En este caso es el grupo general lineal de dimensión dos y entradas reales. Hay muchas formas de probar lo siguiente, y os voy a dar la más fea de todas. Si un vector de $\mathbb{R}^2$ es $(\cos\theta, \sin\theta)$
podemos decir que un vector de $\mathbb{R}^4$ es $(\cos\theta, \sin\theta, 0, 0)$ pero si lo ponemos en forma matricial tendremos algo así como
\begin{equation*}
    \begin{pmatrix}
        \cos\theta & \sin\theta \\  0& 0
    \end{pmatrix}
\end{equation*}
Esta matriz es evidentemente no invertible. Podéis comprobar que una forma de que esta matriz sea inverible y tenga inversa es
\begin{equation*}
    \begin{pmatrix}
        \cos\theta & -\sin\theta \\ \sin\theta & \cos\theta
    \end{pmatrix}
\end{equation*}
Vaya qué sorpresa si es nuestra amiga rotación otra vez! Bueno la cuestión es que observamos lo siguiente.\\
Si consideramos como espacio ambiental del círculo $\mathbb{R}^2$, entonces un elemento general del círculo se representa como $(\cos\theta,\sin\theta)$ que podríamos expresar como $\cos\theta+i\sin\theta$ (tan solo es una base diferente) que sabemos que es $e^{i\theta}$ por tanto, de alguna forma la aplicación que va de $SO(2)\cong S^1$ a $\mathbb{R}^2$ es
\begin{equation*}
    \exp(i \cdot) 
\end{equation*}
Por otro lado, cuando el espacio ambiental es $\mathbb{R}^4$ hemos obtenido que los elementos del círculo se representan como una matriz de rotación que ya vimos se podía escribir como $e^{\theta J}$ donde $J$ es el llamado 'generador de rotaciones' y que era el elemento de la base del álgebra de Lie, según el anterior hilo. Por tanto, la aplicación que va de $SO(2)\cong S^1$ a $\mathbb{R}^4\cong GL(2,\mathbb{R})$ (este isomorfismo no es verdad, no importa en este caso pero no es verdad) es 
\begin{equation*}
    \exp(J \cdot) 
\end{equation*}
En el hilo sobre rotaciones os hablaba de una equivalencia en $i$ y $J$ que no expliqué mucho, tan solo nos basábamos en que tenían las mismas potencias. Ahora podemos hacer más preciso esto, con el concepto de representación. Una representación es una aplicación que preserva la estructura de grupo, esto es, un homomorfismo del grupo de Lie al grupo general lineal de un cierto espacio vectorial.

\begin{equation}
    \Pi: G\longrightarrow GL(V)
\end{equation}
El grupo general lineal de un espacio vectorial $V$ es el grupo de aplicaciones lineales invertibles de ese espacio vectorial. Si $V$ un espacio vectorial real(complejo) de dimensión n entonces escribimos 
$GL(V)=GL(\mathbb{R},n)$ ($GL(V)=GL(\mathbb{C},n)$ )\\
En este sentido, tenemos una especie de teorema de Cayley para grupos continuos, puesto que todos los grupos de Lie definidos como actuación sobre cierto espacio vectorial, $V$, serán subgrupos de este. De esta manera volvemos a recuperar la idea de simetría. Sin embargo, la idea de simetría continua se manifiesta en física de una forma mucho más bonita y precisa, por medio de las cantidades conservadas. Aunque eso es otra historia.\\
La representación de grupos inducce una representación de álgebras
\begin{equation}
    \pi : \mathfrak{g}\rightarrow\mathfrak{g}\mathfrak{l}(V)
\end{equation}
Aquí la condición de preservar la estructura es un homomorfismo entre álgebras que se refiere a
\begin{equation}
    \pi([|X,Y|])=[|\pi(X),\pi(Y)|]
\end{equation}
En representaciones matriciales el paréntesis de Lie se define de una forma particularmente sencilla y conocida:
$$[|\pi(X),\pi(Y)|]:=[\pi(X),\pi(Y)]=\pi(X)\pi(Y)-\pi(Y)\pi(X)$$
Aplicación que llamamos Conmutador. Si $A,B$ son dos matrices y $[A,B]=0$ o sea $AB=BA$ y decimos que conmutan. \\\\
Ahora podemos volver y revisar lo de antes y vemos que lo que teníamos no eran más que dos representaciones de $SO(2)$ 
\begin{equation}
    \Pi_1(x)=e^{\pi_1(J) \theta}=e^{i\theta}
\end{equation}
\begin{equation}
    \Pi_2(x)=e^{\pi_2(J) \theta}=e^{J\theta}
\end{equation}
La primera es una representación de una dimensión que actúa sobre los complejos y la segunda es de dos dimensiones y actúa sobre los vectores de dos dimensiones. En la segunda $\pi_2(J)=J$ porque por alguna representación hay que empezar. Esta es la llama representación definitoria (defining representation) aunque también en física se le llama fundamental, que resulta ser el nombre para otra cosa, pero bueno, es más facil que escribir que definitoria. La representación fundamental es una definición aunque normalmente siempre está dada por ese conjunto algebraico preciso que teníamos antes de que todo se convirtiera en variedades y cosas. Por ejemplo en el caso de $SO(2)$ estas eran las matrices ortogonales de determinante unidad, y por tanto esas tomamos como representación fundamental. Cuando estamos hablando de $U(1)$, la representación fundamental es $e^{i\theta}$ de manera que $\pi_1{i}=i$ claro. La teoría es más general para representaciones complejas así que así nos vamos a quedar de momento. En general se puede demostrar que una representación de $U(1)\cong SO(2)\cong S^1$ general  
$$\Pi: U(1)\rightarrow GL(n, \mathbb{C})$$ 
es de la forma
\begin{equation}
    \Pi(e^{i\theta})= \begin{pmatrix}
        e^{im_1\theta} & &  \\
        & \ddots &   \\
        & & e^{im_n\theta}
    \end{pmatrix}
\end{equation}
Donde $m_i$ son enteros y se llaman pesos  de la representación y veremos que, de hecho, la definen, lo cual está bastante bien. 
\\\\
%FORMA VIEJA
\noindent Los anteriores hilos fueron todos para llegar al concepto central que es el de grupo de Lie. Eso debería ser suficiente para al menos haceros despertar algún gusanillo y a empezar a ver grupos de Lie en más sitios.  A partir de ahora empezamos con cosa concretas y para ir a lo concreto hay que empezar por lo genearl. En esta ocasión vamos a ver las ideas más generales de la teoría de representación. Concretamente, veremos qué es una representación, qué es una representación irreducible y daremos los primeros pasos a ver como se clasifican. \\\\
\noindent Imaginad que hemos estudiado detenidamente un grupo y su álgebra, por ejemplo, de nuevo, el grupo de rotaciones. Para $SO(2)$, estudiamos como actuaba el grupo sobre vectores de $SO(2)$ y concluimos, por simple geometría la forma que debían tener. Pero fijaos en el detalle: para calcular $SO(2)$ (o $SO(n)$, ya que lo hicimos en general) no usamos ni mencionamos nada sobre que tuvieran que actuar en ningún espacio. Lo hicimos sin querer, dando una forma inicial a la curva definida sobre el grupo y después lo encontramos todo por la relación entre el grupo y su álgebra. Esto es interesante, porque da pie a la pregunta ¿Sobre qué objetos actúa entonces realmente el grupo? \\\\
La respuesta nos la da nuestra forma de representar el grupo, por ejemplo, en el caso de $SO(2)$ usamos matrices cuadradas $2\times 2$ de entradas reales, por lo que está claro que actúan sobre $\mathbb{R}^2$. Pero, ¿ y si hubieramos usado otras matrices para representarlas? ¿y si, cuando hicimos nuestro desarrollo de $SO(2)$, en vez de partir de una curva 
$$\gamma(\lambda)= \begin{pmatrix}
	x(\lambda) & y(\lambda) \\
	z(\lambda) & w(\lambda)
\end{pmatrix}$$
hubieramos partido de una matriz $3\times3$ o $4\times4 $ o $10 \times 10$ o $2\times 2$ con componentes complejas? ¿Podemos si quiera hacerlo?  Aunque también te puedes preguntar otra cosa:
¿Por qué me iba a preguntar esto? ¿qué narices me importa a mi? Pues bien, por ejemplo, imagina que has estudiado como rota una partícula que vive en $\mathbb{R}^3$ ¿cómo aplicas lo visto para como rotan dos partículas juntas? ¿Cómo podrías hacerlo? La respuesta está en las representaciones sobre $\mathbb{R}^3\otimes \mathbb{R}^3$ (luego aclaramos este producto, no os preocupéis). O, ¿Cómo rota una función de onda? esto es, ¿cómo rota un elemento de un cierto espacio de Hilbert?. Estas son cuestiones importantes que necesitamos saber en física y la respuesta nos la da, por supuesto, la teoría de representaciones de grupos de Lie. \\\\
Dado un espacio vectorial real de dimensión $n$ finita V, el maxímo grupo de transformaciones sobre los vectores de este espacio es el grupo general lineal de $V$ escrito como $\text{GL}(V)$. Dicho de otra forma, este es el espacio de matrices invertibles $A: V\rightarrow V$. Como es real y de dimensión $n$ resulta ser $\text{GL}(\mathbb{R},n)$. Este grupo, es, de nuevo, un grupo de infinitos elementos y continuo, es el mayor grupo que contiene a todos los grupos de Lie que se puedan definir como transformaciones sobre un espacio vectorial y por ser sus elementos matrices, se llamará Grupo Matricial.\\
Los grupos matriciales son especialmente sencillos porque sabemos cómo calcular la exponencial de una matriz, de otra cosa ya tendríamos que pararnos a pensar ¿verdad?.\\\\ 
Lo que nos interesará a nosotros es poder representar cualquier grupo de Lie como transformaciones lineales sobre un cierto espacio vectorial, es decir, nos interesa representar, válgame, el Grupo utilizando matrices que actuen sobre determinados espacios vectoriales. Para ello, debemos mapear el grupo de Lie $G$ al grupo general lineal de un espacio vectorial $\text{GL}(V)$. \\\\
El homomorfismo $\Pi: G\rightarrow\text{GL}(V)$. se llama representación de $G$.\\\\
Homomorfismo recordad, como estamos tratando dos grupos se refiere a que 
$$\Pi (gh)= \Pi(g)\Pi(h)$$ De la misma forma, debemos hacer lo mismo con el álgebra:
Una representación finitamente dimensional del álgebra de Lie $\mathfrak{g}$ sobre un espacio vectorial $V$ de dimensión finita es un homomorfismo $\pi$: $\mathfrak{g}\rightarrow \mathfrak{g}\mathfrak{l}(V)$ y escribimos $(\pi, V)$. \\\\
Y homomorfismo aquí se refiere a que se preserva el paréntesis de Lie:
$$\pi([|X,Y|])=[|\pi(X),\pi(Y)|]$$
En representaciones matriciales el paréntesis de Lie se define de una forma particularmente sencilla y conocida:
$$[|\pi(X),\pi(Y)|]:=[\pi(X),\pi(Y)]=\pi(X)\pi(Y)-\pi(Y)\pi(X)$$
Aplicación que llamamos Conmutador. Si $A,B$ son dos matrices y $[A,B]=0$ o sea $AB=BA$ y decimos que conmutan
\subsection*{Representaciones reducibles}
\noindent Como hemos visto, representar nos permite, de alguna forma, generalizar la actuación de ciertos grupos. Por ejemplo, estudiar las representaciones del momento angular nos permitirá entender como "giran" las cosas en mecánica cuántica, como empezaba a comentar antes. También la llamada representación adjunta tiene unas propiedades que nos permiten clasificar todos los grupos de Lie, aunque eso es para otro momento, si acaso. \\ La estrategia a seguir para tratar con representaciones será descomponerlas en más pequeñas. Una descomposición de una representación es separar el espacio vectorial sobre el que actúa como suma directa de subespacios
$$V=V_1 \oplus V_2\oplus \dots \oplus V_k$$ Esto se refiere a que un vector en $V$ se puede escribir como un vector $(v_1, v_2, \dots, v_k)$ donde $v_i \in V_i$ Por tanto, en este caso la representación  de un elemento cualquier del grupo $g \in G$ tiene la llamada forma de diagonal en bloque (tendría que explicar por qué? es obvio lol): 
\\
\begin{equation}
    \begin{pmatrix}
        (v_1) \\ (v_2) \\ \vdots \\ (v_k)
    \end{pmatrix}
\end{equation}
\\ 
$$\Pi(g)=\begin{pmatrix}
    \Pi_1(g) \\ & \ddots \\ & & \Pi_k(g)
\end{pmatrix}$$
Cada una de las representaciones $$\Pi_i: G\rightarrow GL(V_i)$$ se le llama subrepresentación de $\Pi$. Si una representación no tiene subrepresentaciones entonces será irreducible. Encontrar las representaciones irreducibles de un álgebra será la idea principal.\\ Para acabar con las representaciones debemos de acabar con el producto tensorial de estas. Los tensores son el terror de los alumnos de primero de física pero en este caso se refiere a algo muy sencillo. Si $\Pi_1: G\rightarrow GL(V), \ \ \Pi_2:H\rightarrow GL(W)$ son representaciones, el producto tensorial de estas es la representación $$\Pi_1 \otimes \Pi_2: G\times H \rightarrow GL(V\otimes W)$$ definida como
$$(\Pi_1 \otimes \Pi_2)(g,h)=\Pi_1(g) \otimes \Pi_2(h)$$
Donde $g\in G$ y $H\in h$. Esta representación nunca será irreducible y admitirá una descomposición de la forma anteior.\\
El espacio $V\otimes W$ es el producto tensorial de $V$ con $W$ y es, como digo, el terror de los alumnos de primero de física. Pero en realidad se refiere a algo muy sencillo. Si $\{e_i\}$ es una base de $V$ y $\{f_j\}$ es una base de $W$, una base de $V\otimes W$ será $\{e_i \otimes f_j\}$ es decir, cogemos las dos bases y las juntamos para hacer otra nueva que serán como bloques. Para más información podéis ver el video sobre tensores de MatesMike... Ah no, que aún no lo ha hecho...\\\\
\\ Finalmente, a nivel del álgebra la representación tensorial actúa como 
$$(\pi_1\otimes \pi_2)(X,Y)=\pi_1(X)\otimes I+I\otimes\pi_2(Y)$$
\\
$$\vec J =\vec S \otimes I+I\otimes\vec L$$
\\
Es decir, actúa primero una sobre un hueco y luego la otra sobre el otro. Esto lo usaremos luego, así que espero que haya quedado claro. 
\subsection{¿Qué esconde SO(3)?}
En el anterior hilo introducimos el concepto de representación. Vimos que una representación no era más que la forma de actuar del grupo sobre un cierto espacio vectorial. Definimos también el concepto de representación irreducible y encontramos las representaciones irreducibles de $U(1)$ que dijimos, como fact, que era el grupo del electromagnetismo, aunque eso lo veremos, en algún momento, muuuuuuucho más tarde. Acabamos también diciendo que solo si $G$ era simplemente conexo existía una correspondencia $1$ a $1$ entre las representaciones del grupo y del álgebra. Nuestro interés ahora que hemos visto $SO(2)$ es ver $SO(3)$ y dado este importante colorario, deberíamos de preguntarnos primero ¿Es $SO(3)$ simplemente conexo? O dicho de otra manera, ¿Tiene $SO(3)$ algún agujero? Y si no tiene ¿Cual es su grupo recubridor universal? \\
Lo primero de todo es construirnos $SO(3)$, siguiendo el estilo que ya he establecido, aunque serán conocidas ya por todos las matrices de rotación en el espacio.\\
Veíamos que $SO(3)$ estaba generado por un álgebra con base 
\begin{equation}
	L_1=\begin{pmatrix}
		0 & -1 & 0 \\ 1 & 0 & 0\\ 0&0&0
	\end{pmatrix}
\end{equation}
\begin{equation}
	L_2=\begin{pmatrix}
		0 & 0 & -1 \\ 0 & 0 & 0\\ 1&0&0
	\end{pmatrix}
\end{equation}
\begin{equation}
	L_3=\begin{pmatrix}
		0 & 0 & 0 \\ 0 & 0 & -1\\ 0&1&0
	\end{pmatrix}
\end{equation}
En términos de representación estamos en la representación fundamental de $SO(3)$ que actúa sobre el espacio $\mathbb{R}^3$. Es interesante cálcular el conmutador entre ellas obteniendo
\begin{equation}
    \begin{split}
        & [L_1,L_2]=L_3 \\
        & [L_2,L_3]=L_1 \\
        & [L_3,L_1]=L_2 
    \end{split}
\end{equation}

Este resulado se suele encapsular en una sola línea como
\begin{equation}
    [L_i, L_j]=\epsilon_{ijk}L_k
\end{equation}
Dónde $\epsilon_{ijk}$ es el símbolo de Levi-Civita.Por ejemplo $[L_1,L_2]$ es $i=1, j=2$ tenemos $\epsilon_{12k}$ $k$ debe necesariamente ser $3$ así que nos da $L_3$. Si $[L_2,L_1]$ tenemos $i=2, j=1$ por tanto $\epsilon_{21k}$ la permutación nos da un $-1$ así que tenemos $-L_3$, claro, recordad que el corchete de Lie es antisimétrico. \\\\
Ahora, por exponenciación, recuperamos las matrices de rotación, que para la rotación en sentido horario son
\begin{equation}
    \begin{split}
        R_z= \begin{pmatrix}
            \cos \theta & -\sin\theta& 0 \\ 
            \sin \theta & \cos \theta &0 \\
            0 &0 & 1 
        \end{pmatrix} \\
        R_y = \begin{pmatrix}
            \cos\theta & 0 &-\sin\theta \\
            0 & 1 & 0\\
            \sin \theta &0 & \cos\theta
        \end{pmatrix}\\
        R_x =\begin{pmatrix}
            1 & 0 &0   \\
            0 & \cos\theta & -\sin\theta \\
            0& \sin\theta & \cos\theta 
        \end{pmatrix}
    \end{split}
\end{equation}
Así, una rotación general de $SO(3)$ es el producto de estas 3 matrices
\begin{equation}
    R_xR_yR_z
\end{equation}
Como decía, lo que vamos a intentar hacer hoy es ver si $SO(3)$ es simplemente conexo o no. Es decir, cualquier camino cerrado puede deformarse continuamente en un punto. . Esto matemáticamente pues es en realidad muy sencillo, porque hay una teoría por ahí que se encarga precisamente de esto, pero esto aunque mis compañeros matemáticos lo saben bien, los físicos no tanto. Para mis compañeros matemáticos, $SO(3)$  es homeomorfo al espacio real proyectivo y su primera homotopía es $\mathbb{Z}/2\mathbb{Z}$  y se acabó, habría que calcular estas cosas y tal pero vamos el statement es eso. Pero vamos a verlo esto más físicamente, porque es lo mío, claro.
Así que, ¿Qué es un camino cerrado en $SO(3)$? Pues es cualquier curva que empiece donde acabe, y en nuestro caso la parametrizaremos por $\lambda$ de manera que será una curva
$$\gamma: [0,2\pi]\rightarrow SO(3)$$
De la forma
$$\gamma(\lambda) =R_x(\lambda)R_y(\lambda)R_z(\lambda)$$
%$$\gamma(\lambda)= \begin{pmatrix}
%	a_{11}(\lambda) & a_{12}(\lambda) & a_{13}(\lambda)\\
%	a_{21}(\lambda) & a_{22}(\lambda) & a_{23}(\lambda) \\
 %   a_{31}(\lambda) & a_{32}(\lambda) & a_{33}(\lambda)
%\end{pmatrix}$$
con $\gamma(0)=\gamma(1)$. Así pues, para cada $\lambda$ tendremos una matriz de rotación. Usaremos un tipo de descomposición muy conocida: las rotaciones de Euler. En general podemos usar cualquier descomposición que queramos pues dependerán de la base, para una rotación normal y corriente los generadores de $SO(3)$ que deducimos ya están bien. Sin embargo,  las rotaciones de Euler son interesantes porque más que rotar cosas nos ayudan a ver la orientación de algún objeto, esto es, respecto de otro sistema de referencia. Mis baldufa simétrica enjoyers saben de lo que hablo. \\
Como digo, los ángulos de Euler son una forma de caracterizar la orientación del sólido rígido en el espacio, consiste en tres ángulos con sus respectivas rotaciones. \\\\
Partimos del espacio 3d cartesiano. Empezamos con él ángulo de precesión  ($\phi$) con una rotación antihoraria alrededor del eje $z$ cartesiano. Así, una rotación $(x,y,z)\rightarrow (x_1,y_1,z_1)$ es 
\begin{center}
    \begin{tikzpicture}
         \draw[very thick, black,->](0,0)--(3,0);
         \draw[very thick, black,->](0,0)--(0,3);
          \draw[very thick, black,->](0,0)--(-1.5,-1.5);
   \node at (3,0) [above] {$y$};
   \node at (0,3) [above left] {$z\equiv z_1$};
   \node at (-1.5,-1.5) [above left ] {$x$};
   \draw[very thick, red,->](0,0)--(2.81907786, 1.02606043);
   \draw[very thick, red,->](0,0)--(-0.89650872, -1.92256915);
   \draw[very thick, red,->](0,0)--(0,3);
      \node at (2.81907786, 1.02606043) [above right] {$y_1$};
      \node at (-0.89650872, -1.92256915) [below right] {$x_1$};
      \draw[thick, ->,red] (1,0) arc (0:20:1);
      \draw[thick, ->,red] (-0.5,-0.5) arc (200:215:1);
        \node at (1,0) [red, above right] {$\phi$};
        \node at (-0.5, -0.5) [red, above left] {$\phi$};
    \end{tikzpicture}
\end{center}
En términos matriciales, supongamos que la base del sistema de referencia inicial es $\{\hat i, \hat j, \hat k\}$ y del rotado $\{\hat i_1, \hat j_2 ,\hat k_3\}$ relacionados por la matriz 
\begin{equation}
    \begin{pmatrix}
        \hat i_1 \\ \hat j_1 \\ \hat k_1 
    \end{pmatrix} =
    \begin{pmatrix}
        \cos \phi & \sin \phi & 0 \\ 
        -\sin \phi & \cos \phi  & 0 \\
        0 & 0 & 1 
    \end{pmatrix}
    \begin{pmatrix}
        \hat i \\ \hat j \\ \hat k
    \end{pmatrix}
\end{equation}
El siguiente ángulo de Euler es el de notación $\theta$ que es un giro antihorario alrededor del nuevo eje $x_1$  
\begin{center}
    \begin{tikzpicture}
         \draw[very thick, black,->](0,0)--(3,0);
         \draw[very thick, black,->](0,0)--(0,3);
          \draw[very thick, black,->](0,0)--(-1.5,-1.5);
   \node at (3,0) [above] {$y$};
   \node at (0,3) [above ] {$z\equiv z_1$};
   \node at (-1.5,-1.5) [above left ] {$x$};
   \draw[very thick, red,->](0,0)--(2.81907786, 1.02606043);
   \draw[very thick, red,->](0,0)--(-0.89650872, -1.92256915);
   \draw[very thick, red,->](0,0)--(0,3);
      \node at (2.81907786, 1.02606043) [above right] {$y_1$};
      \node at (-0.89650872, -1.92256915) [below right] {$x_1\equiv x_2$};
      \draw[thick, ->,red] (1,0) arc (0:20:1);
      \draw[thick, ->,red] (-0.5,-0.5) arc (200:215:1);
        \node at (1,0) [red, above right] {$\phi$};
        \node at (-0.5, -0.5) [red, above left] {$\phi$};

        %angulo 2
         \draw[very thick, green,->](0,0)--(2.29813333, 1.92836283);
          \draw[thick, ->,green] (1.127631144, 0.410424172) arc (0:25:1);
           \node at (1.127631144, 0.410424172) [green, above right] {$\theta$};
           %
        
           \draw[very thick, green,->](0,0)--(-1.02606043,  2.81907786);
          \draw[thick, ->,green] (0,1) arc (90:110:1);
           \node at (0,1) [green, above left] {$\theta$};
           \node at (-1.02606043,  2.81907786)[above left] {$z_2$};
           \node at (2.29813333, 1.92836283)[above right] {$y_2$};
              \draw[very thick, green,->](0,0)--(-0.89650872, -1.92256915);
    \end{tikzpicture}
\end{center}
Matricialmente tendremos
\begin{equation}
    \begin{pmatrix}
        \hat i_2 \\ \hat j_2 \\ \hat k_2
    \end{pmatrix} =
    \begin{pmatrix}
        1 & 0&0 \\
        0 & \cos\theta & \sin\theta \\
        0&-\sin\theta & \cos \theta
    \end{pmatrix}
    \begin{pmatrix}
        \hat i_1 \\ \hat j_1 \\ \hat k_1
    \end{pmatrix}
\end{equation}
Finalmente, tendremos el ángulo de spin $\psi$ que es una rotación antihoraria por el eje $z_2$
\begin{center}
    \begin{tikzpicture}
         \draw[very thick, black,->](0,0)--(3,0);
         \draw[very thick, black,->](0,0)--(0,3);
          \draw[very thick, black,->](0,0)--(-1.5,-1.5);
   \node at (3,0) [above] {$y$};
   \node at (0,3) [above ] {$z\equiv z_1$};
   \node at (-1.5,-1.5) [above left ] {$x$};
   \draw[very thick, red,->](0,0)--(2.81907786, 1.02606043);
   \draw[very thick, red,->](0,0)--(-0.89650872, -1.92256915);
   \draw[very thick, red,->](0,0)--(0,3);
      \node at (2.81907786, 1.02606043) [above right] {$y_1$};
      \node at (-0.89650872, -1.92256915) [below left] {$x_1\equiv x_2$};
      \draw[thick, ->,red] (1,0) arc (0:20:1);
      \draw[thick, ->,red] (-0.5,-0.5) arc (200:215:1);
        \node at (1,0) [red, above right] {$\phi$};
        \node at (-0.5, -0.5) [red, above left] {$\phi$};

        %angulo 2
         \draw[very thick, green,->](0,0)--(2.29813333, 1.92836283);
          \draw[thick, ->,green] (1.127631144, 0.410424172) arc (0:25:1);
           \node at (1.127631144, 0.410424172) [green, above right] {$\theta$};
           %
        
           \draw[very thick, green,->](0,0)--(-1.02606043,  2.81907786);
          \draw[thick, ->,green] (0,1) arc (90:110:1);
           \node at (0,1) [green, above left] {$\theta$};
           \node at (-1.02606043,  2.81907786)[above left, blue] {$z_2\equiv z_3$};
           \node at (2.29813333, 1.92836283)[above right] {$y_2$};

           %angulo 3
            \draw[very thick, blue,->](0,0)--(1.5, 2.59807621);
          \draw[thick, ->,blue] (0.919253332, 0.771345132) arc (0:25:1);
           \node at (0.919253332, 0.971345132) [blue, above right] {$\psi$};
           %
        \draw[very thick, blue,->](0,0)--(-1.02606043,  2.81907786);
        %
        -0.18488525 -2.11324808
        \draw[very thick, blue,->](0,0)--(-0.18488525, -2.11324808);
          \draw[thick, ->,blue] (-0.358603488, -0.76902766) arc (240:260:1);
           \node at (-0.058603488, -0.76902766) [blue, above right] {$\psi$};
           \node at (-0.18488525, -2.11324808) [blue, above right] {$x_3$};
           \node at (1.5, 2.59807621) [blue, above right] {$y_3$};
            \draw[very thick, green,->](0,0)--(-0.89650872, -1.92256915);
    \end{tikzpicture}
\end{center}
Y matricialmente  tenemos
\begin{equation}
    \begin{pmatrix}
        \hat i_3 \\ \hat j_3 \\ \hat k_3 
    \end{pmatrix} =
    \begin{pmatrix}
        \cos \psi & \sin \psi & 0 \\ 
        -\sin \psi & \cos \psi  & 0 \\
        0 & 0 & 1 
    \end{pmatrix}
    \begin{pmatrix}
        \hat i_2 \\ \hat j_2 \\ \hat k_2
    \end{pmatrix}
\end{equation}
Los ángulos de Euler, como digo, nos permiten caracterizar por completo un la orientación de un objeto, esto es, otra base que está rotada respeto de la base cartesiana y, por tanto, cualquier rotación  puede ponerse en términos de ángulos de Euler. En vez de utilizar una rotación para cada eje, podemos utilizar estas. Por tanto, una rotación conjunta de los ángulos de Euler $R_\psi R_\theta R_\phi$ será un elemento de $SO(3)$. No son más que otras 'coordenadas' de $SO(3)$ pero nos ayudan a ver algo interesante. El motivo de de usar estas rotaciones es que son naturales del movimiento del sólido rígido, por ejemplo, el ángulo de precesión describe, efectivamente, la precesión, el decir 'el giro del giro'. Este fenómeno que uno puede observar en peonzas. El ángulo de nutación describe, lo habéis adivinado, el movimiento de nutación. Y el de spin, así es, el movimiento de spin, o sea, el giro per se. Una peonza es un caso concreto de un tipo de cuerpo que es un 'giroscopio' que presenta estos 3 tipos de movimientos. 
%Aquí explicar los angulos de Euler.
\\\\
Podemos usar un giroscopio para mantener la orientación si el giroscopio se encuentra en el llamado soporte de Cardan (Gimbal). El soporte de Cardan está hecho de manera que los posibles cambios de orientación, descritos por ángulos de Euler, que pudiera sufrir el giroscopio, pasen a ser rotaciones en los marcos del soporte. Esta es la idea detrás de los estabilizadores y de muchos sistemas de navegación guiados. Los marcos rotan pero el giroscopio, que da la orientación no, por lo que así pueden entenderse variaciones de orientación, como la desestabilización de todo el sistema. La idea se usa también en programación de objetos 3D. Si alguna vez habéis visto a alguien hacer algo en Maya o en Blender podéis comprobar que para las rotaciones aparecen 3 marcos círculares, cada uno representando un ángulo de Euler.\\ 
Pero fijaos en el aparato ¿Qué ocurre si dos ejes se superponen? ¡Entonces perdemos un grado entero de libertad! Esto es un fenómeno conocido llamado 'Gimbal Lock'. En animación 3D es algo que debe tenerse en cuenta, pues si esto ocurre, da lugar a resultados extraños. 
%Aquí os dejo un fragmento del video: 'UNDERSTANDING GIMBALS & GIMBAL LOCK SOLUTIONS - 3D Animation Tutorial' https://www.youtube.com/watch?v=z3dDsz4f20A&t=176s&ab_channel=SkittyAnimate.  
De la misma forma, esto es algo que debe tenerse en cuenta en sistemas guiados por giroscopio, no es tan solo un problema teórico y causado cuando tenemos acceso a cualquier rotación como en el caso de la animación 3d, el Gimbal Lock tiene consecuencias reales. El caso más famoso fue durante la misión del Apolo 11 en el cual la nave tuvo que dar un giro completo de 180º grados, durante la misión del Apolo 13 también ocurrió, pero en esa ocasión la operación fue mucho más complicada. Hay mucho escrito sobre ello en internet. Aquí os dejo dos videos: 
%https://www.youtube.com/watch?v=oj7v3MXJL3M&t=0s&ab_channel=DavidCambridge
%https://www.youtube.com/watch?v=OmCzZ-D8Wdk&ab_channel=TheVintageSpace 
Puedes incluso vivir el Gimbal Lock en tu brazo, sí, sí. Si intentas rotar tu brazo 360º en horizontal, verás que tu mano se dará la vuelta, cambiará de orientación,vamos. Intentadlo. No os estoy engañando para que hagáis el tonto con el brazo, es totalmente verdad. 
\begin{figure}
    \centering
    \includegraphics[width=0.5\linewidth]{image.png}
    \caption{Enter Caption}
    \label{fig:enter-label}
\end{figure}


Vamos a ver esto, explícitamente una matriz de rotación en términos de ángulos de Euler es
\begin{equation}
   R= \begin{pmatrix}
        \cos \psi & \sin \psi & 0 \\ 
        -\sin \psi & \cos \psi  & 0 \\
        0 & 0 & 1 
    \end{pmatrix}     \begin{pmatrix}
        1 & 0&0 \\
        0 & \cos\theta & \sin\theta \\
        0&-\sin\theta & \cos \theta
    \end{pmatrix} \begin{pmatrix}
        \cos \phi & \sin \phi & 0 \\ 
        -\sin \phi & \cos \phi  & 0 \\
        0 & 0 & 1 
    \end{pmatrix}
\end{equation}
Supongamos que $\theta=\pi/2$ 
\begin{equation}
   R= \begin{pmatrix}
        \cos \psi & \sin \psi & 0 \\ 
        -\sin \psi & \cos \psi  & 0 \\
        0 & 0 & 1 
    \end{pmatrix}     \begin{pmatrix}
        1 & 0&0 \\
        0 & \cos\theta & \sin\theta \\
        0&-\sin\theta & \cos \theta
    \end{pmatrix} \begin{pmatrix}
        0 & 1 & 0 \\ 
        -1 & 0  & 0 \\
        0 & 0 & 1 
    \end{pmatrix}
\end{equation}
Ahora podemos realizar la multiplicación de estas tres matrices y aplicando fórmulas trigonométricas obtenemos
\begin{equation}
R =
\begin{pmatrix}
\cos(\psi + \theta) & \sin(\psi + \theta) & 0 \\
-\sin(\psi + \theta) & \cos(\psi + \theta) & 0 \\
0 & 0 & 1
\end{pmatrix}
\end{equation}
Esto es una rotación antihoraria en el eje z cartesiano, pero está bien porque la hemos obtenido de una curva general dándonos cuenta de un fenómeno físico.  Si ahora definimos $\psi +\phi \in[0,2\pi]$ esta matriz da una curva en $SO(3)$ que es cerrada pero de ninguna manera es contractible a un punto ¿Por qué no? El motivo es porque para $\psi +\phi =\pi$ la orientación cambia.

Imaginad por un momento que $SO(3)$ fuera tan sencillo de describir como una línea, entonces tenemos
\begin{center}
    \begin{tikzpicture}
    \draw[thin, black](-4,0)--(4,0);
     \filldraw[color=black, fill=black, very thick](2,0) circle (0.05);
      \filldraw[color=black, fill=black, very thick](-2,0) circle (0.05);
      \filldraw[color=black, fill=black, very thick](0,0) circle (0.05);
        \node at (2,0) (A) [above]{$R(\pi)$};
        \node at (-2,0) (B) [above ]{$R(\pi)$};
        \node at (0,0) (C) [above][red]{$R(2\pi)$};
        \node at (0,0) (C) [below][red]{$R(0)$};
        \node at (-1,0) (C) [below][black]{$-$};
        \node at (1,0) (C) [below][black]{$+$};
\end{tikzpicture}
\end{center}
La razón por lo que lo he expresado de esa manera es la siguiente: Cuando llegamos a $\lambda=\pi$ la orientación ha cambiado totalmente, tenemos una orientación negativa, por así decirlo. Pensad en vuestra mano, los que hayáis intentado lo que he dicho antes. Este hecho hace que la curva no pueda reducirse a un punto. Para verlo, la historia le ha atribuido a Dirac 'El truco del cinturón'. Podéis comprobar que si cogéis un cinturón y le dais un giro por unpunto, y ahora dejáis uno de los ejes fijos, jamás vais a poder deshacer ese giro. El cinturón es precisamente nuestra curva en $SO(3)$, es el cambio de orientación (el giro) el que impide que sea contractible a un punto. 

En resumen, existen en $SO(3)$ al menos una curva que no es contractible a un punto, y por tanto, $SO(3)$ no es simplemente conexo!
Recordemos las dos implicaciones que tenía esto:
1) $SO(3)$ entonces tiene un grupo simplemente conexo que tiene su misma álgebra $\bar G$.
2) No todas las representaciones del álgebra de $\mathfrak{s}\mathfrak{o}(3)$ darán una representación de $SO(3)$ pero sí de su recubridor universal $\bar G$.\\\\
Podríamos esperar que $SO(3)$ fuera la $3-$esfera y no lo ha sido, pero sí que podemos esperar que su recubridor universal sea la $3-$esfera. Hay muchos motivos para pensar esto, claro, aunque nunca la vida es tan sencilla pues $S^{n(n-1)/2}$ no recubre en general $SO(n)$ de hecho las únicas esferas que admiten estructura de grupo son $S^0$, $S^1$ y $S^3$ y esto está muy relacionado con la existencia de solo 4 álgebras de división normadas. \\\\
Quizás, y solo quizás, los objetos que estábamos usando para describir las rotaciones no son los 'buenos' o dicho de otra forma, no es la representación más natural. Esto es un poco raro pero pensad en $SO(2)$, era más general y natural pensar en representaciones sobre los complejos y, por tanto, representaciones por medio de los complejos negativos.Y decimos que es natural en tanto que pasamos a trabajar con el grupo $U(1)$ y no con $SO(2)$. Resultan ser isomorfos, pero es importante considerar ese detalle.\\
Dicho de otra forma, la forma más 'natural' de entender las rotaciones en el plano era como complejos de módulo $1$ actuando sobre complejos. Resulta que fue Hamilton (ni el del musical ni el de la F1), se supone, el primero en darse cuenta de que los complejos no eran más que equivalentes a un par de reales $(a,b)$ y que podía desarrollarse todo el formalismo de los complejos de esta manera, sin requerir a escribirlos como $a+bi$. Esta equivalencia ya, por supuesto, muy entendida, es muy interesante realmente. Dado que un álgebra de dos números reales describía las rotaciones en el plano, Hamilton pasó muchos años intentando encontrar el álgebra de tres números reales que describieran las rotaciones en el espacio, hoy diríamos que estaba buscando un álgebra de división normada de tres dimensiones.\\\\
\textit{Every morning in the early part of (October) on my coming down to breakfast, your (then) little brother William Edwin, and yourself, used to ask me: 'Well, Papa, can yo multiply triplets?' Whereto I was obliged to reply, with a sad sgake of the head: 'No, I can only add an substract them'.}{William Rowan Hamilton a su hijo }\\\\
Finalmente, el 16 de Octubre de 1843, mientras paseaba con su mujer por el puente de Brougham, en Dublin, le vino la inspiración y se dio cuenta, no podían ser $3$ sino cuatro y se le ocurrió además las reglas de multiplicación que para que no se le olvidarán en el ese mismo momento rayó en el puente. 
$$i^2=j^2=k^2=ijk=-1$$
Y así Hamilton descubrió los cuaterniones. Hamilton se obsesionó con ellos y entre él y sus compañeros póstumos comenzaron a ver varias cosas sobre ellos, incluida la que veremos hoy, que, como sospechábamos, vamos a poder describir rotaciones tridimensionales con ellos. Por la época en muchas escuelas los cuaterniones eran una asignatura más. Unos años más tarde, Gibbs opuso resistencia con su notación del producto escalar $\cdot$ y del producto vectorial $\times$ y hubo una especie de disputa en su momento que supongo que ya sabéis quien acabó ganando. Bueno, al lio. Un cuaternión general es un número de la forma:
\begin{equation}
    q=t+xl+yj+zk
\end{equation}
Donde 
$$l^2=j^2=k^2=ljk=-1$$ conforman un álgebra de división normada con el producto
\begin{equation}
    \begin{split}
        q_1 \cdot q_2  =& (t_1+x_1 l + y_1 j +z_1 k) (t_2+x_2l+y_2 j +z_2 k) \\
        & (t_1t_2 - x_1x_2 - y_1y_2 - z_1z_2) +  \\
        &(t_1x_2 + x_1t_2 + y_1z_2 - z_1y_2)l + \\
        &(t_1y_2 - x_1z_2 + y_1t_2 + z_1x_2)j +\\
        &(t_1z_2 + x_1y_2 - y_1x_2 + z_1t_2)k
    \end{split}
\end{equation}
Que acostumbra ahora a escribirse como
\begin{equation}
    q_1 \cdot q_2 = (t_1t_2 - \vec{v}_1 \cdot \vec{v}_2) + (t_1 \vec{v}_2 + t_2\vec{v}_1 + \vec{v}_1 \times \vec{v}_2)l
\end{equation}
Donde $\vec v_1 =+x_1 l + y_1 j +z_1 k,  \vec v_2=+x_2 l + y_2 j +z_2 k$ Usando la notación de Gibbs.\\
La norma de un cuaternion es
$$|q|^2=q\bar q=(t+xl+yj+zk)(t-xl-yj-zk)=t^2+x^2+y^2+z^2$$
donde se define el complejo conjugado de un cuaternión como
$$ \bar q =t-xl-yj-zk$$
Los  elementos $l,j,k$ están normalizados a la unidad, esto es que su norma es $1$. Alguien que haya cursado relatividad puede llegar a sorprenderse por la aparente similitud entre esta expresión y la métrica de Minkowski, pero a pesar de eso, no hay ninguna forma evidente de relacionarlas. Sin embargo, resulta que en la complexificación de los cuaterniones, se en encuentran las tres representaciones del grupo de Lorentz que aparecen en el modelo estándar. Hablaremos de esto en algún momento, es más fácil de lo que suena. \\\\\
De la misma forma que para los complejos tenemos parte real y parte imaginaria, vamos a definir la parte real como el número múltiplicado por $1$ y la parte imaginaria como todo lo demás, que cambiará de signo bajo la cojugación cuaterniónica, es decir
$$Im(q)=xl+yj+zk$$
Y su conjugado lo unico que hará es un cambio global de signo. 
$$\bar {Im(q)}=-Im(q)=-(xl+yj+zk)$$
Decimos que es de división por el hecho de que podemos dividir, o sea, si 
$q_1 q_2 =0$ o $q_1=0$ o $q_2=0$. 
Esto lo cuento como detalle, el álgebra de cuaterniones, $\mathbb{H}$ (por su descubridor) es la tercera álgebra de división normada (siendo las otras dos los reales ($\mathbb{R}$) y los complejos ($\mathbb{C}$) ).\\\\
Como los complejos, tenemos varias formas posibles de representar los quaterniones. La primera es la que acabamos de decir 
\begin{equation}
    q=t+xl+yj+zk
\end{equation}
Siendo cuatro números reales, nos inspira a pensar que una representación son matrices 2x2 complejas
\begin{equation}
    q=t I+x K_1+yK_2+zK_3
\end{equation}
Donde ha de cumplirse que
\begin{equation}
K_1 K_2 =K_3
\end{equation}
\begin{equation}
K_2K_3 =K_1
\end{equation}
\begin{equation}
    K_3 K_1 =K_2
\end{equation}
hmmm curiosas relaciones si las vemos así no creen? Y si calculamos ahora unos cuantos conmutadores (recordad que estamos en una representación matricial)
\begin{equation}
    [K_1,K_2]=K_3-(-K_3)=2K_3 
\end{equation}
\begin{equation}
    [K_3,K_1]=K_2-(-K_2)=2K_2
\end{equation}
\begin{equation}
    [K_2,K_3]=K_1-(-K_1)=2K_1
\end{equation}
Estas son precisamente las relaciones de conmutación de $\mathfrak{s}\mathfrak{o}(3)$ salvo un factor $2$ (que es factor de escala así que nos da igual o sea, las matrices $K_i/2$ nos dan una representación de $\mathfrak{s}\mathfrak{o}(3)$). Por lo que esto pinta bastante bien, al menos de momento.\\\\
¿Bueno y cuales son estas matrices? Vamos además a pensar que como $K_i^2=-1$ podemos al menos decir que una es nuestra $J$ de $SO(2)$, la tomaremos como la primera
\begin{equation}
    K_2 =\begin{pmatrix}
        0 & -1 \\ 1 & 0
    \end{pmatrix}
\end{equation}
(Decimos que es las dos porque lo vamos a elegir así, simplemente).  Primero, vamos a ver si está normalizada a la unidad, esto es, $|K_2|=K_2 \bar K_2=I$. Como hemos visto antes, en la representación matricial tendremos que $\bar K_i =-K_i$
\\
Como $K_2^2=-I$ tenemos que
$$|K_2|=K_2\bar K_2=K_2(-K_2)=-K_2^2=-(-I)=I$$
Por lo que, mirad qué bien, este amigo que ya habíamos encontrado antes está normalizado ya! \\
Como los $K_i$ son por definición cuaterniones imaginarios, si cada uno de las componentes de la matriz fueran complejos imaginarios puros, entonces si la matriz es la forma
\begin{equation}
    K_i = \begin{pmatrix}
        a & b \\ c & d
    \end{pmatrix} 
\end{equation}
El conjugado cuaterniónico será 
\begin{equation}
    \bar K_i =-K_i= \begin{pmatrix}
        a^* & b^* \\ c^* & d^*
    \end{pmatrix}
\end{equation}
Consideraremos que las matrices son de esta forma, porque así la condición $K_i^2=-I$ nos da además que están normalizadas a la unidad, como acabamos de comprobar. Así, estas serán las dos condiciones que pediremos para encontrar estas matrices. Tomemos una $K_i$ cualquiera, su cuadrado debe ser $-I$ 
\begin{equation}
    \begin{pmatrix}
        a & b \\ c &d
    \end{pmatrix}\begin{pmatrix}
        a & b \\ c &d
    \end{pmatrix} =-\begin{pmatrix}
        1 & 0 \\ 0 & 1 
    \end{pmatrix}
\end{equation}
De donde se deducen las ecuaciones
\begin{equation}
    a^2+bc=-1
\end{equation}
$$ab+bd=0$$ $$ca+dc=0$$ $$cb+d^2=1$$
De las nos ecuaciones del medio deducimos la misma idenitdad $$a=-d$$
Lo que significa que tienen traza nula,cosa que tamibén cumple $K_2$. Supongamos el caso más evidente $a=d=0$ 
la normalización a la unidad nos da
\begin{equation}
    K_i\bar K_i = \begin{pmatrix}
        0 & b \\ c & 0
    \end{pmatrix} \begin{pmatrix}
        0 & b^* \\ c^* & 0
    \end{pmatrix}=\begin{pmatrix}
        bc^* &0  \\  0& cb^*
    \end{pmatrix}=\begin{pmatrix}
        1 & 0 \\ 0 & 1 
    \end{pmatrix}
\end{equation}
Con las ecuaciones de antes y estas tenemos que
\begin{equation}
    bc^*=1
\end{equation}
\begin{equation}
    bc=-1
\end{equation}
Que tiene como solución $b=c=-i$. Por tanto
\begin{equation}
    K_1=\begin{pmatrix}
        0 & -i \\ -i & 0
    \end{pmatrix}
\end{equation}
Y la tercera matriz la obtenemos porque $K_1 K_2 =K_3$ 
\begin{equation}
    K_3=\begin{pmatrix}
        0 & -i \\ -i & 0
    \end{pmatrix} \begin{pmatrix}
        0 &-1\\1 &0
    \end{pmatrix} =\begin{pmatrix}
        -i & 0\\ 0 & i
    \end{pmatrix}
\end{equation}
Lo definimos así porque normalmente se elige la componente $3$ (que corresponderá al eje $z$) como la matriz que es ya diagonal. \\\\
He de dejar claro que estas NO son las matrices de Pauli. Es evidente porque estas matrices son antihermitianas y las matrices de Pauli lo son, pero para conseguirlo como siempre multiplicamos todo por $i$. Definimos entonces las matrices $\sigma_m=iK_m, \ \ m=1,2,3$ 
$$ \sigma_1 = \begin{pmatrix}
    0 & 1 \\ 1 & 0
\end{pmatrix}$$
$$ \sigma_2 = \begin{pmatrix}
    0 & -i \\ i & 0
\end{pmatrix}$$
$$ \sigma_3 = \begin{pmatrix}
    1 & 0 \\  0& -1
\end{pmatrix}$$
Ahora son matrices hermíticas, esto es $A^\dagger = A $ conocidas por todas y todos en este mundo, las matrices de Pauli. De todos modos usaremos las nuestras porque para eso las hemos encontrao.\\\\
Podemos ver primero de todo que son matrices sin traza antihermitianas, esto es $A^\dagger=-A$. Antes veíamos que tenían las mismas relaciones de conmutación que $\mathfrak{s}\mathfrak{o}(3)$ (salvo constante) vemos ahora que tiene todo el sentido del mundo, puesto que una matriz antisimétrica es antihermitiana (fijaos que en el otro sentido no es cierto) Pero estas matrices ¿Qué grupo generan? Pues generarán el recubridor universal de $SO(3)$ (cosa que veremos en seguida) Supongamos que $A$ es una matriz del álgebra, entonces una matriz del grupo será la exponencial $U=e^A$ Consideremos la daga de esta matriz entonces es $$U^\dagger= (e^A)^\dagger= e^{A^\dagger}=e^{-A}$$
$$U U^\dagger =e^Ae^{-A}=I$$
Las matrices $U$ que cumplen esto, es decir, que su daga son su inversa, decimos que son unitarias. Además, vimos la identidad 
\begin{equation}
  \det U= \det (e^A)=e^{\text{tr}A}=e^0=1
\end{equation}
Entonces tenemos que el álgebra de Lie que hemos encontrado genera el grupo de matrices unitarias $2\times 2$ de determinante unidad, $SU(2)$ y su álgebra, siguiendo el mismo estilo que antes, será el álgebra $\mathfrak{s}\mathfrak{u}(2)$. Por lo que tenemos el isomorfismo 
\begin{equation}
    \mathfrak{s}\mathfrak{u}(2) \cong \mathfrak{s}\mathfrak{o}(3)
\end{equation}
Ahora vamos a demostrar que efectivamente $SU(2)$ es el recubridor universal de $SO(3)$. Hay muchas formas de ver que esto es así, vamos a hacerlo de la forma más sencillita claro. Primero de todo ¿qué forma tiene un elemento de $SU(2)$? Como hemos aprendido, para ver esto deberiamos tomar la exponencial, ahora bien, vamos a hacer otra cosa. En términos de NUESTRAS matrices y no las del tal Pauli ese un cuaternion general se escribe
\begin{equation}
    Q=t\begin{pmatrix}
        1 & 0 \\ 0 & 1
    \end{pmatrix} + x\begin{pmatrix}
        0 & -i \\ -i & 0
    \end{pmatrix}+y\begin{pmatrix}
        0 & -1 \\ 1 & 0 
    \end{pmatrix}+z \begin{pmatrix}
        -i & 0 \\ 0 & i
    \end{pmatrix} = \begin{pmatrix}
        t-iz & -ix-y \\ -ix+y & t+iz
    \end{pmatrix}
\end{equation}
Parece que ahora es todo mucho más complicado, pero fijaos, consideremos ahora su traspuesta conjugada 
\begin{equation}
    Q^\dagger= \begin{pmatrix}
        (t-iz)^* & (ix+y)^* \\  (-ix-y)^* & (t+iz)^* 
    \end{pmatrix} =\begin{pmatrix}
        t+iz & -ix +y \\ +ix-y & t-iz
    \end{pmatrix}
\end{equation}
Os podéis entretener si queréis en comprobar que 
\begin{equation*}
    QQ^\dagger =\begin{pmatrix}
        t-iz & -ix-y \\ -ix+y & t+iz
    \end{pmatrix}\begin{pmatrix}
        t+iz & -ix +y \\ +ix-y & t-iz
    \end{pmatrix} =\begin{pmatrix}
        t^2+x^2+y^2+z^2 &0 \\ 0 & t^2+x^2+y^2+z^2
    \end{pmatrix}
\end{equation*}
Esto no parece nada destacable, hasta que hacemos el determinante 
\begin{equation}
    \det Q =(t-iz)(t+iz)-(-y-ix)(-ix+y) =t^2+x^2+y^2+z^2 =Q\bar Q =|Q|^2
\end{equation}
Por tanto si los cuaterniones están normalizados a la unidad, tendremos que 
$$\det Q =1$$
$$  QQ^\dagger= I$$
Por tanto, los cuaterniones normalizados a la unidad son el grupo $SU(2)$. Además la condición 
$$t^2+x^2+y^2+z^2=1$$
Describe una esfera en el espacio $4-$dimensional, por tanto, es la $3-$esfera así que hemos encontrado que
\begin{equation}
    \mathbb{H}(1)\cong SU(2)\cong S^3
\end{equation}
Donde $\mathbb{H}(1)$ son los cuaterniones  normalizados a la unidad. \\\\
Como $    \mathfrak{s}\mathfrak{u}(2) \cong \mathfrak{s}\mathfrak{o}(3)$ y $S_3$ es una $3-$esfera y por tanto es simplemente conexa, $SU(2)\cong S^3$ es el recubridor universal de $SO(3)$. Resulta además un recubrimiento doble, eso significa que $2$ elemento de $SU(2)$ tan el mismo elemento de $SO(3)$. Sencillamente, una vuelta completa de $SU(2)$ en vez de darnos la identidad, nos cambima de signo, esto significa que $U(0)=I\neq U(2\pi)=-I$ mientras que en $SO(3)$ ocurre que $R(0)=R(2\pi)=I$ por lo tanto, la identidad de $SO(3)$ toma dos valores en $SU(2)$ por eso decimos que es su recubrimiento doble. Los recubridores universales de los grupos de rotaciones se llamán grupos de spin $\text{Spin}(n)$. Por tanto en este caso decimos que
$$\text{Spin}(3)\cong SU(2)$$
Y hasta aquí $SO(3)$, en el siguiente empezamos con $SU(2)$.
\subsection*{Pesos y raíces}
Decimos que una matriz $A$ es diagonal si solo tiene elementos en su diagonal principal 
$$\begin{pmatrix}
    a_1 &  &  & &  \\ 
     &a_2 & & & \\
     &  & \ddots& \\ 
    &  &    &a_k \\
\end{pmatrix}$$
Estas matrices cumplen la llamada ecuación de autovalores 
$A v_i = a_i v_i$ 
Donde $v_i$ se llaman autovectores de $A$ y los $a_i$ son los autovalores, que serán los elementos de la diagonal. \\
Ahora bien, si dos matrices conmutan, entonces existe una base en la que ambas son diagonales. No lo vamos a ver formalmente pero observalo se puede hacer de forma sencilla. Supongamos que $A, B$ conmutan y que $u_i$ es un autovector común de ambas
$$A u_i = m_i u_i$$
multiplicamos por $B$ entonces
$$B (A u_i) = m_i B(u_i)$$
Pero $u_i$ también es autovector de $B$ por tanto
$$B (A u_i)= m_i m_i' u_i$$
Donde $m_i'$ es el autovalor de $B$.
Ahora, como conmutan tenemos aedmás que
$$A (B u_i) =m_i m_i' u_i= m_i'A(u_i)$$ 
por tanto 
$$B u_i = m_i' u_i$$
Por tanto será diagonal con mismo autovector. 
\\
Supongamos 
$$A\ket{a'}=a'\ket{a'}$$
$$\bra{a''}[A,B]\ket{a'}= 0$$
$$\bra{a''}AB\ket{a'}-\bra{a''}BA\ket{a'}=(a''-a')\bra{a'}B\ket{a'}=0$$
$$A\ket{a,b}=a\ket{a,b}$$
$$B\ket{a,b}=b\ket{a,b}$$
Un álgebra de Lie típicamente tiene un conjunto de elementos que conmutan entre sí y que conforman una subálgebra que llamamos subálgebra de Cartan $\mathfrak{h}$) Llamamos a los elementos generadores de Cartan. Así,  si $H_1, H_2, \dots, H_n$ son generadores de Cartan entonces $[H_i, H_j]=0$. 
El número de elementos de la subalgebra de Cartan se llama rango del álgebra $\mathfrak{g}$. \\ 
Sabemos que entonces este conjunto puede entero diagonalizarse simultáneamente. Nuestro interés será trabajar con todos los autovalores a la vez, esto es lo que se llama un peso. \\\\
Sea $(\pi, V)$ una representación del álgebra de Lie $\mathfrak{g}$ entonces la tupla $\mu=(m_1, m_2, \dots)$ es un peso de $\pi$ si existe un $v$ no nulo en $V$ tal que
$$\pi(H_1)\ket{v}=m_1\ket{v}, \ \ \  \pi(H_2)\ket{v}=m_2\ket{v}, \ \ \  ..., \ \ \  \pi(H_n)\ket{v}=m_n\ket{v}.$$
Es decir, los pesos son los autovalores simultáneos de $H_1, H_2, \dots, H_n$ en una representación $\pi$ dada. Los vectores $v$ que cumplen esto son los vectores propios de los generadores de Cartan que tienen un nombre propio: vector de pesos y el espacio que generan estos vectores se llama espacio de pesos.\\ Esta definición es importantísima! como uno podría esperar si el teorema se llama Teorema del peso máximo.
\begin{equation}
    \Pi(i)= \begin{pmatrix}
        im_1 & &  \\
        & \ddots &   \\
        & & im_n
    \end{pmatrix}
\end{equation}




\\\\ Para verlo, tomemos por ejemplo el álgebra  
$\mathfrak{s}\mathfrak{u}(2)$ . El grupo $SU(2)$ son las matrices $2\times 2$ unitarias de determinante $1$, es decir las matrices $U$ tales que $U U^\dagger=I$. Su álgebra será $\mathfrak{s}\mathfrak{u}(2)$ que son las matrices $2\times2$ hermíticas $(H=H^\dagger)$ de traza (suma de la diagonal principal) nula. Ésta será un álgebra tridimensional y su base más común son las llamadas matrices de pauli 
$$\sigma_1= \begin{pmatrix}
    0 & 1 \\ 1 & 0
\end{pmatrix}$$
$$\sigma_2=\begin{pmatrix}
    0 & -i \\ i & 0
\end{pmatrix}$$
$$\sigma_3=\begin{pmatrix}
    1 & 0 \\ 0 & -1
\end{pmatrix}$$
Aunque esto es una base perfectamente válida, se usa mucho también la base $J_i=\sigma_i/2 $ como veremos es porque los pesos de $\mathfrak{s}\mathfrak{u}(2)$ se llaman en física spin y por motivos históricos el spin adopta semienteros. Para no confundir a los pobres físicos que esten viendo esto usaremos esta nueva base. Sus relaciones de conmutación son muy fáciles de obtener:
$$[J_i,J_j]=i\epsilon_{ijk} J_k$$
Todas las matrices de $\mathfrak{s}\mathfrak{u}(2)$ no conmutan entre ellas por tanto podemos elegir cualquiera como nuestro generador de Cartan. Elegimos $J_3$ por ser ya diagonal. Sus pesos serán $1/2$ con vector de pesos $(1,0)$ y $-1/2$ con vector de pesos $(0,1)$. Porque 
$$\begin{pmatrix}
    1 & 0 \\ 0 & -1\\
\end{pmatrix} \begin{pmatrix}
    1 \\ 0\\
\end{pmatrix} =\frac{1}{2}\cdot \begin{pmatrix}
    1 \\ 0\\
\end{pmatrix}$$
$$\begin{pmatrix}
    1 & 0 \\ 0 & -1\\
\end{pmatrix} \begin{pmatrix}
    0 \\ 1\\
\end{pmatrix} =-\frac{1}{2}\cdot \begin{pmatrix}
    0 \\ 1\\
\end{pmatrix}$$
Podemos representarlos en el llamado diagrama de pesos:

\begin{center}
    \begin{tikzpicture}
        \draw[ultra thick, black](-2,0)--(2,0);
        \filldraw[color=black, fill=white, very thick](-1,0) circle (0.1);
        \filldraw[color=black, fill=white, very thick](1,0) circle (0.1);
        \node at (-1,0) (A) [below left]{$-1/2$};
        \node at (1,0) (B)[below right]{$1/2$};   
        \node at (2,0)[below]{$J_3$};
\draw [-{>[sep=0pt]}] (-1,0.2) to [bend left] (1,0.2);
\draw [-{>[sep=0pt]}] (1,-0.2) to [bend left] (-1,-0.2);
\node at (0,1) [below]{$+1$}; 
\node at (0,-1) [above]{$-1$}; 
    \end{tikzpicture}
\end{center}
De esta forma podemos concluir algo muy obvio pero de vital importancia: Si estamos en el peso $-1/2$ para pasar al $1/2$ deberemos de sumar $1$. De la misma forma, si estamos en el peso $1/2$ para pasar al $-1/2$ deberemos restar $1$. Además, si estamos en el peso $1/2$ sumar $1$ no nos devuelve un peso, no nos da nada y para $-1/2$ igual, si estamos en este peso y restamos $-1$ nos devuelve nada.\\ Ahora bien, esto es general y lo podemos formalizar. Para ello necesitamos la siguiente definición: \\ 
Una tupla ordenada $\alpha=(a_1, a_2, \dots,a_n)$ es una raíz si $a_1, a_2, \dots$ son no nulos y existe un elemento $Z$ de $\mathfrak{g}$ tal que
$$[H_1,Z]=a_1Z, \ \ \ [H_2,Z]=a_2Z , \ \ \ ...,[H_n,Z]=a_nZ$$
$Z$ es el vector de raíces correspondiente a la raíz $\alpha$. 
Si definimos ahora la aplicación $\text{ad}(H_1)(\cdot)=[H_1,\cdot]$ La asignación $H_1\rightarrow \text{ad}(H_1)(\cdot)$  podemos tomarla como una representación que llamamos representación adjunta. Esta representación es muy importante porque actúa sobre el propio espacio vectorial del álgebra porque $\text{ad}(H_1)(\cdot)$ es una transformación lineal del álgebra (actúa sobre ella). Por tanto, esta dimensión tendrá la misma dimensión que el álgebra. Dada esta definición, entonces, \textbf{las raíces son los pesos de la representación adjunta}.\\
Si $$\pi(X) =ad(X)=[X,\cdot]$$

$$\pi (H_1)\ket{v}=m_1 \ket {v}$$
$$\pi(Z_\alpha) \pi (H_1)\ket{v}=\pi(Z_\alpha) m_1 \ket{v}$$
Por otro lado, por la definición de raíz
$$[\pi(H_1),\pi(Z_\alpha)]=\pi(H_1)\pi(Z_\alpha)-\pi(Z_\alpha)\pi(H_1)=a_1\pi(Z_\alpha)$$
$$\pi(Z_\alpha)\pi(H_1)=\pi(H_1)\pi(Z_\alpha)-a_1\pi(Z_\alpha)$$
$$\pi(H_1)\pi(Z_\alpha)\ket{v}-a_1\pi(Z_\alpha)\ket{v}=\pi(Z_\alpha) m_1 \ket{v}$$
$$\pi(H_k)\pi(Z_\alpha)\ket{v}=(m_k+a_k) \pi(Z_\alpha)\ket{v}$$
Me gustaría subrayar que la dimensión del álgebra y la de sus matrices son cosas distintas.. Por ejemplo el grupo $SU(2)$ de matrices unitarias $2\times 2$ de determinante unidad tiene su álgebra en realidad dimensión $3$. Y el grupo $SU(3)$ del mismo tipo de matrices pero $3\times 3$ tiene como dimensión su álgebra $8$!. Entenderemos bien la diferncia más adelante. \\\\

\noindent Volviendo a $\mathfrak{s}\mathfrak{u}(2)$,  raíces serán las diferencias de pesos, pues como hemos visto deben pasarnos de uno a otro, es decir
$$-1/2 - 1/2 =-1$$
$$1/2 - (-1/2)=1$$
Además, por definición los generadores de Cartán tienen raíz nula así que deberá de corresponderles una raíz en el centro $(0,0)$ siempre. 
El diagrama de raíces, es decir, de pesos de la representación adjunta será
\begin{center}
    \begin{tikzpicture}
        \draw[ultra thick, black](-2,0)--(2,0);
        \filldraw[color=black, fill=white, very thick](-1,0) circle (0.1);
        \filldraw[color=black, fill=white, very thick](1,0) circle (0.1);
        \filldraw[color=black, fill=white, very thick](0,0) circle (0.1);
        \node at (-1,0)[below left]{$-1$};
        \node at (1,0)[below right]{$1$};
        \node at (0,0)[below right]{$0$};
        \node at (2,0)[below]{$J_3$};
    \end{tikzpicture}
\end{center}\
%Volvamos a $\mathfrak{s}\mathfrak{u}(2)$. Tiene tres generadores, que son las tres componentes de  momento angular, la dimensión de su álgebra es tres y por tanto la representación adjunta será la representación tridimensional, que es la de spin 1 con la tercera componente dada por $$L_3=\begin{pmatrix}
   % 1 & 0 &0 \\
   % 0& 0 & 0 \\
   % 0 & 0 &-1
%\end{pmatrix}$$
%de manera que tiene tres vectores de pesos:
%(1 0 0) con peso (1), (0 1 0) con peso (0) y (0 0 1) con peso (-1). De la misma forma que para la representación de spin 1/2, podemos representar los pesos en un diagrama
%\begin{center}
%\begin{tikzpicture}
        %\draw[ultra thick, black](-2,0)--(2,0);
      %  \filldraw[color=black, fill=white, very thick](-1,0) circle (0.1);
       % \filldraw[color=black, fill=white, very thick](1,0) circle (0.1);
       % \node at (-1,0)[below]{$-1$};
       % \node at (1,0)[below]{$1$};     
 %   \end{tikzpicture}
%\end{center}
%Aquí están nuestros tan preciados unos y menos unos!
La idea, entonces, es que "la actuación de las raíces" nos pasa de un peso a otro.  Tenemos la intuición pero resulta ser un teorema: \\
Teorema. (Es muy facil de demostrar, tendría que poner el desarrollo) Sea $ \alpha=(a_1, a_2, ...,a_n)$ una raíz y $Z_a$ en $\mathfrak{g}$ el correspondiente vector de raíces. Sea $\pi$ una representación de $\mathfrak{g}$ y $\mu=(m_1, m_2, \dots)$ el peso asociado a la representación y sea $v$ su vector de pesos correspondiente entonces: 
$$\pi(H_1)\pi(Z_a)v=(m_1+a_1)v, \pi(H_2)\pi(Z_a)v=(m_2+a_2)v, ...$$
entonces, o $\pi(Z)v=0$ o $\pi(Z)$ es un nuevo vector de pesos con peso $\mu+\alpha$. \\ Si $\pi(Z)v=0$ entonces $\mu$ se dice que es el peso máximo de la representación. En física, los elementos $\pi(Z_a)$ se llaman operadores escalera.  \\\\
Ahora podemos decir, por ejemplo, que la representación fundamental de $\mathfrak{s}\mathfrak{u}(2)$ tiene peso máximo $1/2$. Puesto que las raíces son $1$ y $-1$ tendremos que 
\begin{center}
    \begin{tikzpicture}
        \draw[ultra thick, black](-2,0)--(2,0);
        \filldraw[color=black, fill=white, very thick](-1,0) circle (0.1);
        \filldraw[color=black, fill=white, very thick](1,0) circle (0.1);
        \node at (-1,0) (A) [above left]{$-1/2$};
        \node at (1,0) (B)[below right]{$1/2$};   
        \node at (2,0)[below]{$J_3$};
\draw [-{>[sep=0pt]}] (-1,0.2) to [bend left] (1,0.2);
\draw [-{>[sep=0pt]}] (1,-0.2) to [bend left] (-1,-0.2);
\node at (0,1.2) [below]{$\pi(Z_{1})$}; 
\node at (0,-1.2) [above]{$\pi(Z_{-1})$}; 
\draw [-{>[sep=0pt]}] (1,0.2) to [bend left] (3,0.2);
\draw [-{>[sep=0pt]}] (-1,-0.2) to [bend left] (-3,-0.2);
\node at (2,1.2) [below]{$\pi(Z_{1})$}; 
\node at (-2,-1.2) [above]{$\pi(Z_{-1})$}; 
    \node at (3,0.2)[below]{$0$}; 
  \node at (-3,-0.2)[above]{$0$};
    \end{tikzpicture}
\end{center}
Entonces el peso máximo es $1/2$. Esta representación tendrá como peso máximo $1/2$ y ninguna otra más lo tendrá. Esto define de forma única la representación, lo cual es impresionante.\\\\
Ouf, eso fue duro. Pero ahora ya tenemos todos los ingredientes para entender nuestro tan ansiado teorema:\\\\
(Me gustaría como animar algun tipo de brillo o subrayado a los principales objetos del teorema explicados anteriormente, esto es, el algebra de lie, el peso maximo, representación irrep...)\\\\ 
Sea $\mathfrak{g}$ un álgebra de Lie semisimple de rango $n$:\\
1. Toda representación irreducible finitamente dimensional de $\mathfrak{g}$ tiene un peso máximo.\\
2. Dos representaciones irreducibles finitamente dimensionales con mismo peso máximo son isomorfas. \\
3. Si $(\pi,V)$ es una representación irreducible finitamente dimensional de $\mathfrak{g}$ con peso máximo $\mu$, entonces $\mu$ es una n-tupla de enteros positivos.\\
4. Si $\mu$ esuna n-tupla de enteros positivos, entonces existe una representación irreducible finitamente dimensional de $\mathfrak{g}$ con mayor peso $\mu$.\\\\
Lo que este teorema nos dice, por tanto es que si tenemos una representación de un álgebra de Lie (comp. semisimp.) entonces sabemos que tendrá un peso máximo. Y, a la inversa, si tenemos un peso (una $n-$tupla ordenada de enteros positivos) entonces este será el peso máximo de una representación. Con este teorema podemos construir todas las representaciones irreducibles de un álgebra. \\

Ahora, el teorema nos permite construir todas las representaciones de $\mathfrak{s}\mathfrak{u}(2)$. Usualmente se define un peso máximo y mínimo y por medio de actuaciones continuadas de los operadores escalera encontramos todos los pesos. No nos interesa esto a nosotros ahora, lo que nos interesa es la idea. Así, la representación de peso $4$ es decir spin $2$ tendrá como diagrama de pesos: 

\begin{center}
    \begin{tikzpicture}
        \draw[ultra thick, black](-5,0)--(5,0);
        \filldraw[color=black, fill=white, very thick](-1,0) circle (0.1);
        \filldraw[color=black, fill=white, very thick](1,0) circle (0.1);
        \filldraw[color=black, fill=white, very thick](0,0) circle (0.1);
         \filldraw[color=black, fill=white, very thick](-2,0) circle (0.1);
          \filldraw[color=black, fill=white, very thick](2,0) circle (0.1);
        \node [scale=0.7] at (-1.2,0)[below left]{$-1$};
        \node [scale=0.7] at (-2.2,0)[below left]{$-2$};
        \node [scale=0.7] at (0.8,0)[below left]{$1$};
        \node [scale=0.7] at (1.8,0)[below left]{$2$};
        \node [scale=0.7] at (-0.2,0)[below left]{$0$};
        \node [scale=0.7] at (3.8,0)[below left]{$J_3$};

    \draw [-{>[sep=0pt]}] (-1,0.2) to [bend left] (0,0.2);
    \draw [-{>[sep=0pt]}] (0,-0.2) to [bend left] (-1,-0.2);
    \draw [-{>[sep=0pt]}] (0,0.2) to [bend left] (1,0.2);
    \draw [-{>[sep=0pt]}] (-1,-0.2) to [bend left] (-2,-0.2);
      \draw [-{>[sep=0pt]}] (1,0.2) to [bend left] (2,0.2);
          \draw [-{>[sep=0pt]}] (-2,0.2) to [bend left] (-1,0.2);
    \draw [-{>[sep=0pt]}] (-2,-0.2) to [bend left] (-3,-0.2);
    \draw [-{>[sep=0pt]}] (1,-0.2) to [bend left] (0,-0.2);
        \draw [-{>[sep=0pt]}] (2,-0.2) to [bend left] (1,-0.2);
    \draw [-{>[sep=0pt]}] (2,0.2) to [bend left] (3,0.2);
\node [scale=0.6] at (-1.5,-1) [above]{$\pi(Z_{-1})$}; 
\node [scale=0.6] at (-2.5,-1) [above]{$\pi(Z_{-1})$}; 
\node [scale=0.6] at (-1.5,-1) [above]{$\pi(Z_{-1})$}; 
\node [scale=0.6] at (-0.5,-1) [above]{$\pi(Z_{-1})$};
\node [scale=0.6] at (0.5,-1) [above]{$\pi(Z_{-1})$};
\node [scale=0.6] at (1.5,-1) [above]{$\pi(Z_{-1})$};

\node [scale=0.6] at (-1.5,+1) [below]{$\pi(Z_{-1})$};  
\node [scale=0.6] at (-1.5,+1) [below]{$\pi(Z_{-1})$}; 
\node [scale=0.6] at (-0.5,+1) [below]{$\pi(Z_{-1})$};
\node [scale=0.6] at (0.5,+1) [below]{$\pi(Z_{-1})$};
\node [scale=0.6] at (1.5,+1) [below]{$\pi(Z_{-1})$};
\node [scale=0.6] at (2.5,+1) [below]{$\pi(Z_{-1})$};
\end{tikzpicture}
\end{center}
Lo que hacemos es simplemente tomar el peso máximo, spin $2$, por simetría el peso mínimo será spin $-2$ y por la actuación de las raíces encontramos todos los demás. Esto nos vale siempre, he aquí el poder del teorema del peso máximo. Así, solo con este teorema podemos decir: \\ 
Hemos encontrado así cuales son todas las posibles irreps de $\mathfrak{s}\mathfrak{u}(2)$. \\\\
Ahora, nos interesa encontrar como descomponer una representación en irreps. Si tomamos el producto tensorial de una irrep con otra sabemos que será una representación reducible. Vamos a encontrar como descomponerla en suma de irreducibles. \\ Supongamos el producto de irreps:
$1/2\otimes s$ Un vector de pesos de este estado será $\ket{1/2, m_{1/2}}\otimes \ket{s,m_s}$ donde el primer número de cada vector se refiere al spin de la representación (peso máximo) y el segundo a los posibles valores de spin (posibles pesos), La actuación del generador de Cartan sobre este vector de pesos será la representación tensorial, de acuerdo a lo anterior, entonces:
$$J_3=J_3^{1/2}\otimes I+I\otimes J_3^s $$
que al actuar sobre el estado es 

$$J_3 (\ket{1/2, m_{1/2}}\otimes \ket{s,m_s}) =(J_3^{1/2}\otimes I+I\otimes J_3^s)(\ket{1/2, m_{1/2}}\otimes \ket{s,m_s})=$$
$$=J_3^{1/2} \ket{1/2, m_{1/2}} \otimes I  \ket{s,m_s}+I\ket{1/2, m_{1/2}}\otimes J_3^s\ket{s,m_s}=$$
$$=m_{1/2}\ket{1/2, m_{1/2}} \otimes  \ket{s,m_s}+m_s\ket{1/2, m_{1/2}}\otimes\ket{s, m_s}$$
Por tanto



$$J_3 (\ket{1/2, m_{1/2}}\otimes \ket{s,m_s}) =(m_s+m_{1/2}) \ket{1/2, m_{1/2}}\otimes \ket{s,m_s}$$ 
\\ 

Supongamos que la representación que actua sobre este estado es irreducible con pesos $(m_s+m_{1/2})$ ¿cuales son los pesos máximos posibles?. Encontrar esto será equivalente a encontrar la descomposición en irreps. \\ 
Ahora, $m_s$ puede tomar cualquier valor desde $-s$ hasta $s$ pero $m_{1/2}$ solo puede tomar $1/2, -1/2$ así pues, el valor máximo de $m_j=m_s+m_{1/2}$ es $j=s+1/2$ y $j=s-1/2$ por tanto decimos que 
$s\otimes 1/2 = (s+1/2)\oplus (s-1/2)$. Podemos seguir este razonamiento para cualquier spin y tendremos al final el resultado:
$$j_1\otimes j_2 = (j_1+j_2)\oplus (j_1+j_2-1)\oplus\dots\oplus|j_1-j_2|$$
Esta descomposición se llama descomposición de Clebsch-Gordan. 
$$(\pi_i \otimes \pi_j)(X)(\ket{\psi}\otimes\ket{\phi})$$
$$(\pi_i\otimes\pi_j)(J_3)(\ket{m_i, i}\otimes\ket{m_j,j})=\pi_i(J_3)\ket{m_i,i}\otimes \ket{m_j,j}+\ket{m_i, i}\otimes \pi_j(J_3)\ket{m_j,j}$$
$$(\pi_i\otimes\pi_j)(J_3)(\ket{m_i, i}\otimes\ket{m_j,j})=(m_i+m_j)(\ket{m_i,i}\otimes\ket{m_j,j})$$
$$(\pi_i\otimes\pi_j)(J_3)(\ket{i, i}\otimes\ket{j,j})=(i+j)(\ket{i,i}\otimes\ket{j,j})$$
$$\ket{1/2,1/2}\otimes\ket{1/2,1/2}\rightarrow1/2+1/2=1$$
$$\ket{1/2,1/2}\otimes\ket{-1/2,1/2}\rightarrow1/2-1/2=0$$
$$\ket{-1/2,1/2}\otimes\ket{1/2,1/2}\rightarrow-1/2+1/2=0$$
$$\ket{-1/2,1/2}\otimes\ket{-1/2,1/2}\rightarrow-1/2-1/2=-1$$

\begin{equation}
    \begin{split}
        &\ket{1/2,1/2}\otimes\ket{1/2,1/2}\rightarrow1/2+1/2=1 \\
        &\ket{1/2,1/2}\otimes\ket{-1/2,1/2}\rightarrow1/2-1/2=0 \\
        &\ket{-1/2,1/2}\otimes\ket{1/2,1/2}\rightarrow-1/2+1/2=0 \\
        &\ket{-1/2,1/2}\otimes\ket{-1/2,1/2}\rightarrow-1/2-1/2=-1
    \end{split}
\end{equation}
\subsection*{$\mathfrak{s}\mathfrak{u}(3)$}
Visto $\mathfrak{s}\mathfrak{u}(2)$ no debe ser muy difícil entender $\mathfrak{s}\mathfrak{u}(3)$.\\
A diferencia de $\mathfrak{s}\mathfrak{u}(2)$ no hay forma fácil de expresar las relaciones de conmutación de este álgebra. $\mathfrak{s}\mathfrak{u}(3)$ es una álgebra $8-$dimensional debido a que es el álgebra de Lie de las matrices $3\times 3$ unitarias, que tomamos como representación fundamental. La base más común de este álgebra es la llamada base de Gell-Mann o matrices de Gell-Mann:\\
$$
\lambda_1 =\begin{pmatrix}
    0 & 1 & 0 \\ 1 & 0 & 0\\0 &0& 0\\
\end{pmatrix}
\lambda_2 =\begin{pmatrix}
    0 & -i & 0 \\ i & 0 & 0\\ 0 &0& 0\\
\end{pmatrix}
\lambda_3 =\begin{pmatrix}
    1 & 0 & 0 \\ 0 & -1 & 0\\ 0 &0& 0\\
\end{pmatrix}
\lambda_4 =\begin{pmatrix}
    0 & 0 & 1 \\ 0 & 0 & 0\\ 1 &0& 0\\
\end{pmatrix}
$$

$$
\lambda_5 =\begin{pmatrix}
    0 & 0 & -i \\ 0 & 0 & 0\\i &0& 0\\
\end{pmatrix}
\lambda_6 =\begin{pmatrix}
    0 & 0 & 0 \\ 0 & 0 & 1\\ 0 &1& 0\\
\end{pmatrix}
\lambda_7 =\begin{pmatrix}
    0 & 0 & 0 \\ 0 & 0 & -i\\ 0 &i& 0\\
\end{pmatrix}
\lambda_8 =\frac{1}{\sqrt{3}}\begin{pmatrix}
    1 & 0 & 0 \\ 0 & 1 & 0\\ 0 &0& -2\\
\end{pmatrix}
$$

$\mathfrak{s}\mathfrak{u}(3)$ tiene dos matrices que conmutan entre sí $\lambda_2$ y $\lambda_8$ y el resto no. Por tanto, la subalgebra de Cartan de $\mathfrak{s}\mathfrak{u}(3)$ es la generada por estos dos elementos y por tanto $\mathfrak{s}\mathfrak{u}(3)$ es un álgebra de Lie de rango 2, lo que nos dice que su espacio de pesos será bidimensional. Esto, como veremos, nos complicará mucho los diagramas de pesos pero no os preocupéis porque nos toparemos con un ya conocido y agradable amigo.\\\\
La representación fundamental que estamos tratando aquí sera la de dimensión $3$ así que al encontrar los pesos de esta rerpresntación estaremos encontrando los pesos de la primera representación irreducible de $\mathfrak{s}\mathfrak{u}(3)$.  Sin embargo, en esta base los pesos tendrán valores extraños no enteros de la misma forma que en $\mathfrak{s}\mathfrak{u}(2)$ los pesos eran semienteros. Como no nos importan los cálculos, vamos a tomar una base en la que los generadores de Cartan tienen esta forma:  \\
$$T_3=\begin{pmatrix}
    1 &0 & 0 \\ 0 & -1 & 0 \\ 0 & 0 & 0\\
\end{pmatrix}$$
$$T_8=\begin{pmatrix}
    0 &0 & 0 \\ 0 & 1 & 0 \\ 0 & 0 & -1\\
\end{pmatrix}$$
De esta forma los pesos y sus respectivos vecotres de pesos son evidentes:
 
$$\begin{pmatrix}
    1 \\ 0 \\ 0 
\end{pmatrix} \longrightarrow (1,0) $$
$$\begin{pmatrix}
     0\\ 1\\ 0 
\end{pmatrix} \longrightarrow (-1,1) $$
$$ \ \ \ \ \  \begin{pmatrix}
    0 \\ 0 \\ 1 
\end{pmatrix} \longrightarrow (0,-1) $$
Siempre representaremos los pesos en el llamado diagrama de pesos en el que los ejes son $(m,n)$ las componentes de los pesos. En este caso tendremos: 
\begin{center}
    \begin{tikzpicture}
        \draw[thick, gray](-2,0)--(2,0);
        \draw[thick, gray](0,-2)--(0,2);
        
        \draw[very thick, black](1,0)--(-1,1)--(0,-1)--cycle;
        \filldraw[color=black, fill=white, very thick](1,0) circle (0.1);
        \filldraw[color=black, fill=white, very thick](-1,1) circle (0.1);
        \filldraw[color=black, fill=white, very thick](0,-1) circle (0.1);
        \node at (1,0)[above right]{$(1,0)$};
        \node at (-1,1)[above left]{$(-1,1)$};  
        \node at (0,-1)[below left]{$(0,-1)$}; 
    \end{tikzpicture}
\end{center}
Sin embargo, aún podemos realizar un cambio más y es definir las cantidades definir los ejes $k,l$ tales que un peso de la forma $(m,n)$ en el diagrama de pesos se escriba como $s=km+ln$. Estos ejes se definirán tal que el eje $l$ describa un ángulo de  $60^\circ$ con el de $k$. Entonces, los ejes ahora serán los ejes de simetría del triángulo equilátero, como veíamos antes: 
\begin{center}
    \begin{tikzpicture}
        \draw[thick, gray](-3,0)--(3,0);
  \draw[thick, gray](-2,3.46410)--(2,-3.46410);
        \draw[thick, gray](-2,-3.46410)--(2,3.46410);
        \draw[very thick, black](1,0)--(-0.5,0.8660254)--(-0.5,-0.8660254)--cycle;
        \filldraw[color=black, fill=white, very thick](1,0) circle (0.1);
        \filldraw[color=black, fill=white, very thick](-0.5,0.8660254) circle (0.1);
        \filldraw[color=black, fill=white, very thick](-0.5,-0.8660254) circle (0.1);
        \node at (1,0)[below]{$(1,0)$};
        \node at (0,1)[above]{$(-1,1)$};  
        \node at (-1,-1)[below]{$(0,-1)$}; 
        \node at (2,3.46410)[above]{$l$}; 
        \node at (3,0)[below]{$k$}; 








        
    \end{tikzpicture}
\end{center}
Las raíces serán las diferencias de los pesos
$$
(1,0)-(-1,1)=(2,1)
$$
$$
(1,0)-(0,-1)=(1,1)
$$
$$
(-1,1)-(1,0)=(-2,1)
$$
$$
(-1,1)-(0,-1)=(-1,2)
$$
$$
(0,-1)-(1,0)=(-1,-1)
$$
$$
(0,-1)-(-1,1)=(-1,-2)
$$
\begin{center}
    \begin{tikzpicture}
        \draw[thick, gray](-3,0)--(3,0);
  \draw[thick, gray](-2,3.46410)--(2,-3.46410);
        \draw[thick, gray](-2,-3.46410)--(2,3.46410);
        \draw[very thick, black](1,0)--(-0.5,0.8660254)--(-0.5,-0.8660254)--cycle;
        \filldraw[color=black, fill=white, very thick](1,0) circle (0.1);
        \filldraw[color=black, fill=white, very thick](-0.5,0.8660254) circle (0.1);
        \filldraw[color=black, fill=white, very thick](-0.5,-0.8660254) circle (0.1);
        \node at (1,0)[below]{$(1,0)$};
        \node at (0,1)[above]{$(-1,1)$};  
        \node at (-1,-1)[below]{$(0,-1)$}; 
        \node at (2,3.46410)[above]{$l$}; 
        \node at (3,0)[below]{$k$}; 


        \draw [-{>[sep=5pt]}, red] (-0.5,0.8660254) to [bend right] (1,0);
        \draw [-{>[sep=5pt]}, green] (1,0) to [bend right] (-0.5,0.8660254);

         \draw [-{>[sep=5pt]}, red] (-0.5,-0.8660254) to [bend right] (-0.5,0.8660254);
            \draw [-{>[sep=5pt]}, green] (-0.5,0.8660254) to [bend right] (-0.5,-0.8660254);

            \draw [-{>[sep=5pt]}, red] (1,0) to [bend right] (-0.5,-0.8660254);
            \draw [-{>[sep=5pt]}, green] (-0.5,-0.8660254) to [bend right] (1,0);






        
    \end{tikzpicture}
\end{center}
por tanto el diagrama de raíces, es decir, de pesos de la representación adjunta será
\begin{center}
    \begin{tikzpicture}
        \draw[thick, gray](-3,0)--(3,0);
        \draw[thick, gray](-2,3.46410)--(2,-3.46410);
        \draw[thick, gray](-2,-3.46410)--(2,3.46410);
        \draw[very thick, black](1.5,0.8660254)--(0,1.73200)--(-1.5,0.8660254)--(-1.5,-0.8660254)--(0,-1.73200)--(1.5,-0.8660254)--cycle;
        \filldraw[color=black, fill=white, very thick](1.5,0.8660254) circle (0.1);
         \filldraw[color=black, fill=white, very thick](-1.5,0.8660254) circle (0.1);
        \filldraw[color=black, fill=white, very thick](-1.5,-0.8660254) circle (0.1);
        \filldraw[color=black, fill=white, very thick](1.5,-0.8660254) circle (0.1);
        \filldraw[color=black, fill=white, very thick](0,1.73200) circle (0.1);
        \filldraw[color=black, fill=white, very thick](0,-1.73200) circle (0.1); 
        \filldraw[color=black, fill=white, very thick](0,0) circle (0.1);
        \filldraw[color=black, fill=none, very thick](0,0) circle (0.2);
           \node at (2,3.46410)[above]{$l$}; 
        \node at (3,0)[below]{$k$}; 
    \end{tikzpicture}
\end{center}
Donde también debemos poner los pesos de los generadores de Cartan que son $(0,0)$ en esta representación por definición.\\\\
Podemos ver que ambos diagramas de pesos presentan una simetría. Esto no es casualidad, y se puede demostrar que pasará siempre.  La simetría del diagrama de pesos de  $\mathfrak{s}\mathfrak{u}(3)$ es la simetría del triángulo rectángulo, $S_3$ ¡Qué casualidad que ya lo habíamos estudiado antes!\\
En general, el grupo de simetría de un diagrama de pesos en general se llama grupo de Weyl, así pues, diremos que el grupo de Weyl de $\mathfrak{s}\mathfrak{u}(3)$ es $S_3$. \\
\begin{center}
    \begin{tikzpicture}
        \draw[thick, gray](-3,0)--(3,0);
        \draw[thick, gray](-2,3.46410)--(2,-3.46410);
        \draw[thick, gray](-2,-3.46410)--(2,3.46410);
        \draw[very thick, black](1.5,0.8660254)--(0,1.73200)--(-1.5,0.8660254)--(-1.5,-0.8660254)--(0,-1.73200)--(1.5,-0.8660254)--cycle;
        \draw[very thick, black](1,0)--(-0.5,0.8660254)--(-0.5,-0.8660254)--cycle;
          \node at (2,3.46410)[above]{$l$}; 
        \node at (3,0)[below]{$k$}; 
    \end{tikzpicture}
\end{center}
Si empezamos con un peso, entonces todos los pesos creados como reflejos por los tres ejes de simetría del triángulo, es decir, la acción del grupo de Weyl, estarán en el diagrama de pesos. \\ Este es el último ingrediente que nos quedaba para construir cualquier representación $\mathfrak{s}\mathfrak{u}(3)$. Recapitulemos:\\
(i) El teorema de peso máximo nos asegura que dada una representación tendremos un peso máximo salvo (up to) isomorfismo. \\
(ii) Por acción del grupo de Weyl, podemos construir el resto de pesos del diagrama.
(iii) La acción de las raíces construye nuevas representaciones hasta el peso máximo.\\
Así obtenemos la receta de cocina para construcción de representaciones irreducibles de $\mathfrak{s}\mathfrak{u}(3)$:\\
1. Usamos el grupo de Weyl esto nos definirá un polígono.\\
2. Construimos la red de raíces y tomamos pesos que estén dentro de la red centrada en cada uno de los pesos. \\
Lo último que nos queda es ver la multiplicidad de estos espacios, esto ya es mucho más complicado y está muy fuera de lo que quiere o puede hacer este video, pero como regla general:\\   (i) La capa externa siempre tiene multiplicidad 1.\\
(ii) Si el perimetro es un hexágono  lo de dentro tendrá multiplicidad $+1$. \\
(iii) Si lo de fuera es un triángulo, lo de dentro tendrá la misma multiplicidad.
\\
Ejemplo: (se puede quitar este ejemplo porque no es nada luego) Diagrama de pesos de $(1,2)$. 
El paso uno es dibujar este peso en nuestro diagrama
\begin{center}
    \begin{tikzpicture}
        \draw[thick, gray](-3,0)--(3,0);
        \draw[thick, gray](-2,3.46410)--(2,-3.46410);
        \draw[thick, gray](-2,-3.46410)--(2,3.46410);
         \filldraw[color=black, fill=white, very thick](2,1.73200) circle (0.1);
           \node at (2,3.46410)[above]{$l$}; 
        \node at (3,0)[below]{$k$}; 
    \end{tikzpicture}
\end{center}
El siguiente paso es construir los pesos exteriores por acción del grupo de Weyl
\begin{center}
    \begin{tikzpicture}
        \draw[thick, gray](-3,0)--(3,0);
        \draw[thick, gray](-2,3.46410)--(2,-3.46410);
        \draw[thick, gray](-2,-3.46410)--(2,3.46410);
         \filldraw[color=black, fill=white, very thick](2,1.73200) circle (0.1);
         \filldraw[color=black, fill=white, very thick](0.5,2.598024) circle (0.1);
          \filldraw[color=black, fill=white, very thick](-2.5,0.866024) circle (0.1);
        \filldraw[color=black, fill=white, very thick](-2.5,-0.866024) circle (0.1);
         \filldraw[color=black, fill=white, very thick](0.5,-2.598024) circle (0.1);
         \filldraw[color=black, fill=white, very thick](2,-1.73200) circle (0.1);
           \node at (2,3.46410)[above]{$l$}; 
        \node at (3,0)[below]{$k$}; 
    \end{tikzpicture}
\end{center}
Construimos el polígono resultante
\begin{center}
    \begin{tikzpicture}
        \draw[thick, gray](-3,0)--(3,0);
        \draw[thick, gray](-2,3.46410)--(2,-3.46410);
        \draw[thick, gray](-2,-3.46410)--(2,3.46410);
        \draw[very thick, black](2,1.73200)--(0.5,2.598024)--(-2.5,0.866024)--(-2.5,-0.866024)--(-2.5,0.866024)--(-2.5,-0.866024)--(0.5,-2.598024)-- (2,-1.73200)--cycle;
         \filldraw[color=black, fill=white, very thick](2,1.73200) circle (0.1);
         \filldraw[color=black, fill=white, very thick](0.5,2.598024) circle (0.1);
          \filldraw[color=black, fill=white, very thick](-2.5,0.866024) circle (0.1);
        \filldraw[color=black, fill=white, very thick](-2.5,-0.866024) circle (0.1);
         \filldraw[color=black, fill=white, very thick](0.5,-2.598024) circle (0.1);
         \filldraw[color=black, fill=white, very thick](2,-1.73200) circle (0.1);
           \node at (2,3.46410)[above]{$l$}; 
        \node at (3,0)[below]{$k$}; 
    \end{tikzpicture}
\end{center}
Lo siguiente es construir los pesos debido a la acción de las raíces, es decir, ponemos los pesos que pertenezcan a la red de raíces centrada en cada uno de los pesos 
\begin{center}
    \begin{tikzpicture}
        \draw[thick, gray](-3,0)--(3,0);
        \draw[thick, gray](-2,3.46410)--(2,-3.46410);
        \draw[thick, gray](-2,-3.46410)--(2,3.46410);
        \draw[very thick, black](2,1.73200)--(0.5,2.598024)--(-2.5,0.866024)--(-2.5,-0.866024)--(-2.5,0.866024)--(-2.5,-0.866024)--(0.5,-2.598024)-- (2,-1.73200)--cycle;
        \draw[very thick, black](0.5,0.8666)--(-1,0)--(0.5,-0.8666)--cycle;
         \filldraw[color=black, fill=white, very thick](2,1.73200) circle (0.1);
         \filldraw[color=black, fill=white, very thick](0.5,2.598024) circle (0.1);
          \filldraw[color=black, fill=white, very thick](-2.5,0.866024) circle (0.1);
        \filldraw[color=black, fill=white, very thick](-2.5,-0.866024) circle (0.1);
         \filldraw[color=black, fill=white, very thick](0.5,-2.598024) circle (0.1);
         \filldraw[color=black, fill=white, very thick](2,-1.73200) circle (0.1);
         \filldraw[color=black, fill=white, very thick](2,0) circle (0.1);
         \filldraw[color=black, fill=white, very thick](-1,1.732) circle (0.1);
         \filldraw[color=black, fill=white, very thick](-1,-1.732) circle (0.1);
         
         \filldraw[color=black, fill=white, very thick](0.5,0.8666) circle (0.1);
          \filldraw[color=black, fill=white, very thick](-1,0) circle (0.1);
           \filldraw[color=black, fill=white, very thick](0.5,-0.8666) circle (0.1);
            \node at (2,3.46410)[above]{$l$}; 
        \node at (3,0)[below]{$k$};

       % \draw [-{>[sep=5pt]}, green] (2,1.73200)  to [bend right] (0.5,2.598024);
       % \draw [-{>[sep=5pt]}, green] (0.5,2.598024)  to [bend right] (-1,1.732);
      %  \draw [-{>[sep=5pt]}, green] (-1,1.732)  to [bend right] (-2.5,0.866);
        %\draw [-{>[sep=5pt]}, green] (-2.5,0.866)  to [bend right] (-2.5,-0.866);
      %  \draw [-{>[sep=5pt]}, green] (-2.5,-0.866)  to [bend right] (-1,-1.732);
       % \draw [-{>[sep=5pt]}, green] (-1,-1.732)  to [bend right] (0.5,-2.598024);
        % \draw [-{>[sep=5pt]}, green] (0.5,-2.598024)  to [bend right] (2,-1.73200);
        %  \draw [-{>[sep=5pt]}, green] (2,-1.73200) to [bend right] (2,0);
           %\draw [-{>[sep=5pt]}, green] (2,0) to [bend right] (2,1.73200);

        %\draw [-{>[sep=5pt]}, green] (2,1.73200)  to [bend right] (0.5,0.8666);
        %\draw [-{>[sep=5pt]}, green] (0.5,0.8666)  to [bend right] (2,0);
        %\draw [-{>[sep=5pt]}, green] (2,0)  to [bend right] (0.5,-0.8666);
        % \draw [-{>[sep=5pt]}, green] (0.5,-0.8666)  to [bend right] (2,-1.73200);
            

    \end{tikzpicture}
\end{center}
Lo que nos deja el diagrama de pesos siguiente:

\begin{center}
    \begin{tikzpicture}
        \draw[thick, gray](-3,0)--(3,0);
        \draw[thick, gray](-2,3.46410)--(2,-3.46410);
        \draw[thick, gray](-2,-3.46410)--(2,3.46410);
        \draw[very thick, black](2,1.73200)--(0.5,2.598024)--(-2.5,0.866024)--(-2.5,-0.866024)--(-2.5,0.866024)--(-2.5,-0.866024)--(0.5,-2.598024)-- (2,-1.73200)--cycle;
        \draw[very thick, black](0.5,0.8666)--(-1,0)--(0.5,-0.8666)--cycle;
         \filldraw[color=black, fill=white, very thick](2,1.73200) circle (0.1);
         \filldraw[color=black, fill=white, very thick](0.5,2.598024) circle (0.1);
          \filldraw[color=black, fill=white, very thick](-2.5,0.866024) circle (0.1);
        \filldraw[color=black, fill=white, very thick](-2.5,-0.866024) circle (0.1);
         \filldraw[color=black, fill=white, very thick](0.5,-2.598024) circle (0.1);
         \filldraw[color=black, fill=white, very thick](2,-1.73200) circle (0.1);
         \filldraw[color=black, fill=white, very thick](2,0) circle (0.1);
         \filldraw[color=black, fill=white, very thick](-1,1.732) circle (0.1);
         \filldraw[color=black, fill=white, very thick](-1,-1.732) circle (0.1);
         
         \filldraw[color=black, fill=white, very thick](0.5,0.8666) circle (0.1);
          \filldraw[color=black, fill=white, very thick](-1,0) circle (0.1);
           \filldraw[color=black, fill=white, very thick](0.5,-0.8666) circle (0.1);
            \node at (2,3.46410)[above]{$l$}; 
        \node at (3,0)[below]{$k$};



    \end{tikzpicture}
\end{center}
Finalmente, tendremos que poner la multiplicidad, como hemos dicho, la externa tendrá siempre multiplicidad 1 porque es una irrep. y el polígono interior por ser triángulo en un hexágono será de multiplicidad 2. \\
Finalmente encontramos que la $(1,2)$ de $\mathfrak{s}\mathfrak{u}(3)$ será 
\begin{center}
    \begin{tikzpicture}
        \draw[thick, gray](-3,0)--(3,0);
        \draw[thick, gray](-2,3.46410)--(2,-3.46410);
        \draw[thick, gray](-2,-3.46410)--(2,3.46410);
        \draw[very thick, black](2,1.73200)--(0.5,2.598024)--(-2.5,0.866024)--(-2.5,-0.866024)--(-2.5,0.866024)--(-2.5,-0.866024)--(0.5,-2.598024)-- (2,-1.73200)--cycle;
        \draw[very thick, black](0.5,0.8666)--(-1,0)--(0.5,-0.8666)--cycle;
         \filldraw[color=black, fill=white, very thick](2,1.73200) circle (0.1);
         \filldraw[color=black, fill=white, very thick](0.5,2.598024) circle (0.1);
          \filldraw[color=black, fill=white, very thick](-2.5,0.866024) circle (0.1);
        \filldraw[color=black, fill=white, very thick](-2.5,-0.866024) circle (0.1);
         \filldraw[color=black, fill=white, very thick](0.5,-2.598024) circle (0.1);
         \filldraw[color=black, fill=white, very thick](2,-1.73200) circle (0.1);
         \filldraw[color=black, fill=white, very thick](2,0) circle (0.1);
         \filldraw[color=black, fill=white, very thick](-1,1.732) circle (0.1);
         \filldraw[color=black, fill=white, very thick](-1,-1.732) circle (0.1);
         
         \filldraw[color=black, fill=white, very thick](0.5,0.8666) circle (0.1);
          \filldraw[color=black, fill=white, very thick](-1,0) circle (0.1);
           \filldraw[color=black, fill=white, very thick](0.5,-0.8666) circle (0.1);
           
           \node at (2,1.73200) [above]{$1$}; 
           \node at (0.5,2.598024) [above]{$1$};
           \node at (-2.5,0.866024) [above]{$1$};
           \node at (-2.5,-0.866024) [above left]{$1$};
           \node at (0.5,-2.598024) [above]{$1$};
           \node at (2,-1.73200) [above right ]{$1$};
           \node at (2,0) [above right ]{$1$};
           \node at (-1,1.732) [above ]{$1$};
           \node at (-1,-1.732)  [above right ]{$1$};

            \node at (0.5,0.8666)  [above right ]{$2$};
             \node at (-1,0)  [above]{$2$};
              \node at (0.5,-0.8666)  [above right ]{$2$};
   \node at (2,3.46410)[above]{$l$}; 
        \node at (3,0)[below]{$k$};
           
           
           
    \end{tikzpicture}
    
\end{center}
Construcción de la $10$ o sea $(3,0)$ \\ 
El paso uno es dibujar este peso en nuestro diagrama
\begin{center}
    \begin{tikzpicture}
        \draw[thick, gray](-4,0)--(4,0);
        \draw[thick, gray](-2,3.46410)--(2,-3.46410);
        \draw[thick, gray](-2,-3.46410)--(2,3.46410);
         \filldraw[color=black, fill=white, very thick](3,0) circle (0.1);
           \node at (2,3.46410)[above]{$l$}; 
        \node at (3,0)[below]{$k$}; 
    \end{tikzpicture}
\end{center}
El siguiente paso es construir los pesos exteriores por acción del grupo de Weyl
\begin{center}
    \begin{tikzpicture}
        \draw[thick, gray](-3,0)--(3,0);
        \draw[thick, gray](-2,3.46410)--(2,-3.46410);
        \draw[thick, gray](-2,-3.46410)--(2,3.46410);
         \filldraw[color=black, fill=white, very thick](3,0) circle (0.1);
          \filldraw[color=black, fill=white, very thick](-1.5,2.59024) circle (0.1);
                \filldraw[color=black, fill=white, very thick](-1.5,-2.59024) circle (0.1);
                  \node at (2,3.46410)[above]{$l$}; 
        \node at (3,0)[below]{$k$}; 
    \end{tikzpicture}
\end{center}
Construimos el polígono resultante
\begin{center}
    \begin{tikzpicture}
        \draw[thick, gray](-3,0)--(3,0);
        \draw[thick, gray](-2,3.46410)--(2,-3.46410);
        \draw[thick, gray](-2,-3.46410)--(2,3.46410);
        \draw[very thick, black](3,0)--(-1.5,2.59024)--(-1.5,-2.598024)--cycle;
         \filldraw[color=black, fill=white, very thick](3,0) circle (0.1);
          \filldraw[color=black, fill=white, very thick](-1.5,2.59024) circle (0.1);
                \filldraw[color=black, fill=white, very thick](-1.5,-2.59024) circle (0.1);
                  \node at (2,3.46410)[above]{$l$}; 
        \node at (3,0)[below]{$k$}; 
    \end{tikzpicture}
\end{center}
Lo siguiente es construir los pesos debido a la acción de las raíces, es decir, ponemos los pesos que pertenezcan a la red de raíces centrada en cada uno de los pesos 
\begin{center}
    \begin{tikzpicture}
        \draw[thick, gray](-3,0)--(3,0);
        \draw[thick, gray](-2,3.46410)--(2,-3.46410);
        \draw[thick, gray](-2,-3.46410)--(2,3.46410);
        \draw[very thick, black](3,0)--(-1.5,2.59024)--(-1.5,-2.598024)--cycle;
         \filldraw[color=black, fill=white, very thick](3,0) circle (0.1);
          \filldraw[color=black, fill=white, very thick](-1.5,2.59024) circle (0.1);
        \filldraw[color=black, fill=white, very thick](-1.5,-2.59024) circle (0.1);
           \filldraw[color=black, fill=white, very thick](0,-1.732) circle (0.1);
           \filldraw[color=black, fill=white, very thick](0,1.732) circle (0.1);
           \filldraw[color=black, fill=white, very thick](1.5,-0.866024) circle (0.1);
        \filldraw[color=black, fill=white, very thick](1.5, 0.866024) circle (0.1);
        \filldraw[color=black, fill=white, very thick](-1.5, 0.866024) circle (0.1);
         \filldraw[color=black, fill=white, very thick](-1.5, -0.866024) circle (0.1);
         \filldraw[color=black, fill=white, very thick](0,0) circle (0.1);
           \node at (2,3.46410)[above]{$l$}; 
        \node at (3,0)[below]{$k$}; 
    \end{tikzpicture}
\end{center}
Finalmente, tendremos que poner la multiplicidad, como hemos dicho, la externa tendrá siempre multiplicidad 1 porque es una irrep. y el polígono interior por ser el exterior un triangulo entonces tendrá mult. 1
Finalmente encontramos que la $(3,0)$ de $\mathfrak{s}\mathfrak{u}(3)$ será 
\begin{center}
    \begin{tikzpicture}
        \draw[thick, gray](-3,0)--(3,0);
        \draw[thick, gray](-2,3.46410)--(2,-3.46410);
        \draw[thick, gray](-2,-3.46410)--(2,3.46410);
        \draw[very thick, black](3,0)--(-1.5,2.59024)--(-1.5,-2.598024)--cycle;
         \filldraw[color=black, fill=white, very thick](3,0) circle (0.1);
          \filldraw[color=black, fill=white, very thick](-1.5,2.59024) circle (0.1);
        \filldraw[color=black, fill=white, very thick](-1.5,-2.59024) circle (0.1);
           \filldraw[color=black, fill=white, very thick](0,-1.732) circle (0.1);
           \filldraw[color=black, fill=white, very thick](0,1.732) circle (0.1);
           \filldraw[color=black, fill=white, very thick](1.5,-0.866024) circle (0.1);
        \filldraw[color=black, fill=white, very thick](1.5, 0.866024) circle (0.1);
        \filldraw[color=black, fill=white, very thick](-1.5, 0.866024) circle (0.1);
         \filldraw[color=black, fill=white, very thick](-1.5, -0.866024) circle (0.1);
         \filldraw[color=black, fill=white, very thick](0,0) circle (0.1);
          \node at (3,0)[above]{$1$}; 
          \node at (-1.5,2.59024)[above]{$1$}; 
          \node at (-1.5,-2.59024)[above left]{$1$}; 
          \node at (0,-1.732)[above]{$1$};
          \node at (0,1.732)[above]{$1$};
          \node at (1.5,-0.866024)[above]{$1$}; 
          \node at (1.5,0.866024)[above]{$1$};
          \node at (-1.5,-0.866024)[above left]{$1$};
          \node at (-1.5,0.866024)[above left]{$1$};
          \node at (0,0)[above]{$1$};
        \node at (2,3.46410)[above]{$l$}; 
        \node at (3,0)[below]{$k$}; 
          
    \end{tikzpicture}
\end{center}
Por un simple ejercicio de conteo podemos ver que la dimensión de esta representación es $10$ por tanto, esta representación se suele llamar la $10$. \\ 
\newpage
\subsection*{Aplicaciones: Isospin y modelo de Quarks}
Debemos notar una cosa de las representaciones de $\mathfrak{s}\mathfrak{u}(3)$. Hemos visto que el diagrama de pesos de la $(1,0)$ es 
\begin{center}
    \begin{tikzpicture}
        \draw[thick, gray](-3,0)--(3,0);
  \draw[thick, gray](-2,3.46410)--(2,-3.46410);
        \draw[thick, gray](-2,-3.46410)--(2,3.46410);
        \draw[very thick, black](1,0)--(-0.5,0.8660254)--(-0.5,-0.8660254)--cycle;
        \filldraw[color=black, fill=white, very thick](1,0) circle (0.1);
        \filldraw[color=black, fill=white, very thick](-0.5,0.8660254) circle (0.1);
        \filldraw[color=black, fill=white, very thick](-0.5,-0.8660254) circle (0.1);
        \node at (1,0)[below]{$(1,0)$};
        \node at (0,1)[above]{$(-1,1)$};  
        \node at (-1,-1)[below]{$(0,-1)$}; 
        \node at (2,3.46410)[above]{$l$}; 
        \node at (3,0)[below]{$k$}; 
    \end{tikzpicture}
\end{center}

Pero, podríamos haber construido la $(0,1)$. Usando las reglas de construcción antes vistas 
\begin{center}
    \begin{tikzpicture}
        \draw[thick, gray](-3,0)--(3,0);
        \draw[thick, gray](-2,3.46410)--(2,-3.46410);
        \draw[thick, gray](-2,-3.46410)--(2,3.46410);
        \draw[very thick, black](0.5,0.866025)--(-1,0)--(0.5,-0.8660254)--cycle;
        \filldraw[color=black, fill=white, very thick](0.5,0.866025) circle (0.1);
        \filldraw[color=black, fill=white, very thick](-1,0) circle (0.1);
        \filldraw[color=black, fill=white, very thick](0.5,-0.8660254) circle (0.1);
        \node at (0.5,0.866025) [above right]{$(0,1)$};
        \node at (-1,0)[above left]{$(-1,0)$};  
        \node at (0.5,-0.8660254)[above right]{$(1,-1)$}; 
        \node at (2,3.46410)[above]{$l$}; 
        \node at (3,0)[below]{$k$}; 
    \end{tikzpicture}
\end{center}

Observamos que los pesos de la representación $(0,1)$ son exactamente los negativos de la $(1,0)$. En general podemos demostrar que esto no es casualidad y es porque las representaciones $(0,m)$ son representaciones duales, es decir si $(m,0)$ es una representación sobre el espacio $V_m$, la representación $(0,m)$ actuará sobre su dual, $V_m^{*}$. Como ambas representaciones tienen dimensión $3$ para diferenciarlas solemos llamar a la $(0,1)$ la representación conjugada de $(1,0)$ y escribimos $\bar {3}$.

A lo largo del siglo de XX, se fueron descubriendo nuevas partículas, a parte de los piones y las deltas, que no podían explicarse usando solo la simetría de isospin. Gell-Mann, representó las partículas en función de su carga, y estrañeza y para cada spin dibujó los siguientes diagramas en función de la carga y la extrañeza (la extrañeza no es más que otra etiqueta) Encontró que había partículas que eran más ligeras y tenían spin $0$ y se organizaban de acuerdo a un octete
\begin{center}
    \begin{tikzpicture}
        %\draw[thick, gray](-3,0)--(3,0);
        % \draw[thick, gray](-2,3.46410)--(2,-3.46410);
        \node at (4.5,-0.8666) [above right] {$q=1$};
        \node at (3,-1.732) [above right]{$q=0$};
        \node at (1.5,-2.5980) [above right]{$q=-1$};
        
                \draw[thick, gray](-3,3.4640)--(4.5,-0.8666);
        \draw[thick, gray](-3,1.732)--(3,-1.732);
        \draw[thick, gray](-4.5,0.866)--(1.5,-2.5980);
        \draw[thick, gray](-3,-3.4640)--(4.5,0.8666);
        \draw[thick, gray](-3,-1.732)--(3,+1.732);
         \node at (3,+1.732) [above right] {$s=0$};
        \node at (1.5,2.5980) [above right]{$s=1$};
        \node at (4.5,0.8666) [above right]{$s=-1$};
    
        
        \draw[thick, gray](-4.5,-0.866)--(1.5,2.5980);
        \draw[very thick, black](1.5,0.8660254)--(0,1.73200)--(-1.5,0.8660254)--(-1.5,-0.8660254)--(0,-1.73200)--(1.5,-0.8660254)--cycle;
        \filldraw[color=black, fill=white, very thick](1.5,0.8660254) circle (0.1);
         \filldraw[color=black, fill=white, very thick](-1.5,0.8660254) circle (0.1);
        \filldraw[color=black, fill=white, very thick](-1.5,-0.8660254) circle (0.1);
        \filldraw[color=black, fill=white, very thick](1.5,-0.8660254) circle (0.1);
        \filldraw[color=black, fill=white, very thick](0,1.73200) circle (0.1);
        \filldraw[color=black, fill=white, very thick](0,-1.73200) circle (0.1); 
        \filldraw[color=black, fill=white, very thick](0,0) circle (0.1);
        \filldraw[color=black, fill=none, very thick](0,0) circle (0.2);
        \node at (0,1.73200)[above right]{$K^{+}$};
        \node at (-1.5,0.8660254)[above]{$K^0$};
        \node at (-1.5,-0.8660254)[above left]{$\pi^{-}$};
        \node at (0,-1.73200)[below]{$K^{-}$};
        \node at (1.5,-0.8660254)[above right]{$\bar K^{0}$};
         \node at (1.5,0.8660254)[above]{$\pi^{+}$};
         \node at (0,0)[above right]{$\pi^{0}$};
         \node at (0,0)[below left]{$\eta$};
    \end{tikzpicture}
\end{center}
Y también habían partículas más pesadas que podían tener spin $1/2$ o spin $3/2$. Las de spin $1/2$ se organizaban de acuerdo a un octete 
\begin{center}
    \begin{tikzpicture}
        %\draw[thick, gray](-3,0)--(3,0);
        % \draw[thick, gray](-2,3.46410)--(2,-3.46410);
        \node at (4.5,-0.8666) [above right] {$q=1$};
        \node at (3,-1.732) [above right]{$q=0$};
        \node at (1.5,-2.5980) [above right]{$q=-1$};
        
                \draw[thick, gray](-3,3.4640)--(4.5,-0.8666);
        \draw[thick, gray](-3,1.732)--(3,-1.732);
        \draw[thick, gray](-4.5,0.866)--(1.5,-2.5980);
        \draw[thick, gray](-3,-3.4640)--(4.5,0.8666);
        \draw[thick, gray](-3,-1.732)--(3,+1.732);
         \node at (3,+1.732) [above right] {$s=-1$};
        \node at (1.5,2.5980) [above right]{$s=0$};
        \node at (4.5,0.8666) [above right]{$s=-2$};
        
        
        \draw[thick, gray](-4.5,-0.866)--(1.5,2.5980);
        \draw[very thick, black](1.5,0.8660254)--(0,1.73200)--(-1.5,0.8660254)--(-1.5,-0.8660254)--(0,-1.73200)--(1.5,-0.8660254)--cycle;
        \filldraw[color=black, fill=white, very thick](1.5,0.8660254) circle (0.1);
         \filldraw[color=black, fill=white, very thick](-1.5,0.8660254) circle (0.1);
        \filldraw[color=black, fill=white, very thick](-1.5,-0.8660254) circle (0.1);
        \filldraw[color=black, fill=white, very thick](1.5,-0.8660254) circle (0.1);
        \filldraw[color=black, fill=white, very thick](0,1.73200) circle (0.1);
        \filldraw[color=black, fill=white, very thick](0,-1.73200) circle (0.1); 
        \filldraw[color=black, fill=white, very thick](0,0) circle (0.1);
        \filldraw[color=black, fill=none, very thick](0,0) circle (0.2);
        \node at (0,1.73200)[above right]{$p$};
        \node at (-1.5,0.8660254)[above]{$n$};
        \node at (-1.5,-0.8660254)[above left]{$\Sigma^{-}$};
        \node at (0,-1.73200)[below]{$\Xi^{-}$};
        \node at (1.5,-0.8660254)[above right]{$\Xi^{0}$};
         \node at (1.5,0.8660254)[above]{$\Sigma^{+}$};
         \node at (0,0)[above right]{$\Sigma^{0}$};
         \node at (0,0)[below left]{$\Lambda$};
    \end{tikzpicture}
\end{center}
Y las de spin $3/2$ se organizaban de esta forma
\begin{center}
    \begin{tikzpicture}
        %\draw[thick, gray](-4,0)--(4,0);
        \draw[thick, gray](-3,3.4640)--(4.5,-0.8666);
        \draw[thick, gray](-3,1.732)--(3,-1.732);
        \draw[thick, gray](-4.5,0.866)--(1.5,-2.5980);
        %\draw[thick, gray](-3,-1.732)--(0,-3.4640);
        %\draw[thick, gray](-2,3.46410)--(2,-3.46410);
        %\draw[thick, gray](-2,-3.46410)--(2,3.46410);
        %\draw[very thick, black](3,0)--(-1.5,2.59024)--(-1.5,-2.598024)--cycle;
         \draw[thick, gray](-1.5,-3.5)--(-1.5,3.5);
         \draw[thick, gray](0,-3.5)--(0,3.5);
         \draw[thick, gray](1.5,-3.5)--(1.5,3.5);
         \draw[thick, gray](3,-3.5)--(3,3.5);
         
        \node at (-1.5,3.5) [above] {$q=-1$};
        \node at (0,3.5)[above]{$q=-0$};
        \node at (1.5,3.5) [above]{$q=+1$};
        \node at (3,3.5) [above] {$q=+2$};
         
        \node at (4.5,-0.8666) [above right] {$s=0$};
        \node at (3,-1.732) [above right]{$s=-1$};
        \node at (1.5,-2.5980) [above right]{$s=-2$};
       % \node at (0,-3.4640) [above right] {$s=-3$};

        
         \filldraw[color=black, fill=white, very thick](3,0) circle (0.1);
          \filldraw[color=black, fill=white, very thick](-1.5,2.59024) circle (0.1);
        %\filldraw[color=black, fill=white, very thick](-1.5,-2.59024) circle (0.1);
           \filldraw[color=black, fill=white, very thick](0,-1.732) circle (0.1);
           \filldraw[color=black, fill=white, very thick](0,1.732) circle (0.1);
           \filldraw[color=black, fill=white, very thick](1.5,-0.866024) circle (0.1);
        \filldraw[color=black, fill=white, very thick](1.5, 0.866024) circle (0.1);
        \filldraw[color=black, fill=white, very thick](-1.5, 0.866024) circle (0.1);
         \filldraw[color=black, fill=white, very thick](-1.5, -0.866024) circle (0.1);
         \filldraw[color=black, fill=white, very thick](0,0) circle (0.1);
          \node at (3,0)[above right]{$\Delta^{++}$}; 
          \node at (-1.5,2.59024)[above right]{$\Delta^{-}$}; 
          %\node at (-1.5,-2.59024)[above right]{$\Omega^{-}$}; 
          \node at (0,-1.732)[above right]{$\Xi^{* 0}$};
          \node at (0,1.732)[above right]{$\Delta^0$};
          \node at (1.5,-0.866024)[above right]{$\Sigma^{* +}$}; 
          \node at (1.5,0.866024)[above right]{$\Delta^{+}$};
          \node at (-1.5,-0.866024)[above right]{$\Xi^{* -}$};
          \node at (-1.5,0.866024)[above right]{$\Sigma^{* -}$};
          \node at (0,0)[above right]{$\Sigma^{* 0}$};
          
    \end{tikzpicture}
\end{center}
Los bariones de spin $1/2$ nos recuerdan a la representación adjunta o $8$ de $\mathfrak{s}\mathfrak{u}(3)$ pero es que los de spin $3/2$ nos recuerdan a la $10$! Gell-Mann postuló que los bariones estaban compuestos por un trío de una partícula elemental elemental: los quarks. Los quarks son objetos que transforman bajo la representación fundamental de $\mathfrak{s}\mathfrak{u}(3)$, la $3$.  Decimos que hay tres sabores de quarks: up, down y strange. Además son fermiones de spin 1/2. Así, el estado general de un quark, $\ket{\psi}$ pertencerá a $\mathbb{C}^3$ y será una combinación lineal de los tres sabores.\begin{center}
    \begin{tikzpicture}
         \draw[very thick, black,->](0,0)--(3,0);
         \draw[very thick, black,->](0,0)--(0,3);
          \draw[very thick, black,->](0,0)--(-1.5,-1.5);
   \node at (3,0) [above] {down};
   \node at (0,3) [above left] {up};
   \node at (-1.5,-1.5) [above left ] {strange};
   \draw[very thick, red,->](0,0)--(1.5,2);
      \node at (1.5,2) [above right] {$\ket{\psi}$};
    \end{tikzpicture}
\end{center}
O simbólicamente 
$$\ket{\psi}=\alpha \ket{\text{down}} +\beta \ket{\text{up}}+\gamma \ket{\text{strange}}$$
Así, el espacio de los estados de una combinación de $3$ quarks y por tanto será tres copias de este espacio $\mathbb{C}^3\otimes \mathbb{C}^3 \otimes \mathbb{C}^3$ o dicho de otra forma, la representación bajo la que transformarán será la $3\otimes 3 \otimes 3$. Ahora bien, esta representación, como ya vimos, es claramente no irreducible, así pues nos interesará encontrar su descomposición como suma directa de irreps. Al hacer esto, encontraremos los subespacios en los que actua la representación como irrep, y por tanto, encontraremos las partículas que se crean al combinar los quarks.\\ Partimos primero de este resultado:\\
Si $\{e_i\}_{i=1,2,3}$ es una base de $\mathbb{C}^3$, hemos visto ya que además serán vectores de peso con peso $\mu_1= (1,0)$, $\mu_2=(-1,1)$ y  $\mu_3=(0,-1)$, entonces una base de $\mathbb{C}^3\otimes \mathbb{C}^3 \otimes \mathbb{C}^3$ será $\{e_i\otimes e_j\otimes e_k\}_{i,j,k=1,2,3}$ y además cualquier elemento de está base será un vector de pesos con peso $\mu_i+\mu_j+\mu_k$ y pertencerá al espacio de pesos de peso máximo $\mu_1+\mu_1+\mu_1=(3,0)$ Para verlo vemos la actuación de los generadores de Cartan. Igual que antes con $\mathfrak{s}\mathfrak{u}(2)$ el resultado será la suma de pesos de cada representación.
$$H_1 H_2(e_i\otimes e_j\otimes e_k) = (m_i+m_j+m_k)(e_i\otimes e_j\otimes e_k)$$
Así pues, el peso máximo de la representación será la suma de pesos máximos,  de nuevo, como ya vimos en $\mathfrak{s}{u}(2)$.\\ 
Además, como son fermiones de spin $1/2$, podemos calcular la descomposión en irreps del momento angular total $1/2 \otimes 1/2 \otimes 1/2$. De acuerdo con la descomposición C-G será: 
$$1/2 \otimes 1/2 \otimes 1/2= (1\oplus0)\otimes1/2=(1\otimes1/2)\oplus(0\otimes1/2)=3/2\oplus1/2\oplus1/2$$
Así pues, esperamos para los bariones estados de spin total $3/2$ y $1/2$. También por ser Fermiones, por Dirac sabemos que debe existir el anti-quark, el antiquark transformará de acuerdo a la $\bar {3} $. Que actuará en el espacio de base $\theta_j$ tal que $\theta_j(e_i) =\delta_{ji}$. Podemos encontrar los pesos de las distintas combinaciones de productos tensoriales de $e_i$ con $\theta_j$ que serán:

$$e_1 \otimes \theta_1 \rightarrow (1,0)+(0,1) =(1,1) $$
$$e_1 \otimes \theta_2 \rightarrow (1,0) + (-1,0) =(0,0)$$
$$e_1 \otimes \theta_3 \rightarrow (1,0) + (1,-1) =(2,-1)$$
$$e_2 \otimes \theta_1 \rightarrow (-1,1) + (1,1) =(-1,2)$$
$$e_2 \otimes \theta_2 \rightarrow (-1,1) + (-1,0) =(-2,1)$$
$$e_2 \otimes \theta_3 \rightarrow (-1,1) + (1,-1) =(0,0)$$
$$e_3 \otimes \theta_1 \rightarrow (0,-1) + (0,1) =(0,0)$$
$$e_3 \otimes \theta_2 \rightarrow (0,-1) + (-1,0) =(-1,-1)$$
$$e_3 \otimes \theta_3 \rightarrow (0,-1) + (1,-1) =(1,-2)$$
El peso $(0,0)$ tiene multiplicidad $3$ y el resto de pesos forman un hexágono con multiplicidad $1$
\begin{center}
    \begin{tikzpicture}
        \draw[thick, gray](-3,0)--(3,0);
        \draw[thick, gray](-2,3.46410)--(2,-3.46410);
        \draw[very thick, black](1.5,0.8660254)--(0,1.73200)--(-1.5,0.8660254)--(-1.5,-0.8660254)--(0,-1.73200)--(1.5,-0.8660254)--cycle;
        \filldraw[color=black, fill=white, very thick](1.5,0.8660254) circle (0.1);
         \filldraw[color=black, fill=white, very thick](-1.5,0.8660254) circle (0.1);
        \filldraw[color=black, fill=white, very thick](-1.5,-0.8660254) circle (0.1);
        \filldraw[color=black, fill=white, very thick](1.5,-0.8660254) circle (0.1);
        \filldraw[color=black, fill=white, very thick](0,1.73200) circle (0.1);
        \filldraw[color=black, fill=white, very thick](0,-1.73200) circle (0.1); 
        \filldraw[color=black, fill=white, very thick](0,0) circle (0.1);
        \node at (0,1.73200)[above right]{$1$};
        \node at (-1.5,0.8660254)[above left]{$1$};
        \node at (-1.5,-0.8660254)[above left]{$1$};
        \node at (0,-1.73200)[below]{$1$};
        \node at (1.5,-0.8660254)[above right]{$1$};
         \node at (1.5,0.8660254)[above right]{$1$};
         \node at (0,0)[below left]{$3$};
           \draw[thick, gray](-2,3.46410)--(2,-3.46410);
        \draw[thick, gray](-2,-3.46410)--(2,3.46410);
         \node at (2,3.46410)[above]{$l$}; 
        \node at (3,0)[below]{$k$};
    \end{tikzpicture}
\end{center}
Por tanto, vemos que se trata de la adjunta + la trivial que tiene peso $(0,0)$ decimos entonces que
$$(1,0)\otimes (0,1) =(1,1)\oplus (0,0)$$
O como suele ser más común, en términos de dimensiones.
$$3 \otimes \bar{3} =8\oplus1 $$
Así, las partículas más ligeras de spin $0$, los mesones, serían pares quark-antiquark. \\
Siendo esto un éxito nos podemos ahora preguntar, cuantas combinaciones posibles de $3$ quarks existen. Esperamos obtener dos grandes grupos: las de spin $1/2$ y las de spin $3/2$. Podemos encontrar los pesos de las distintas combinaciones de $e_i,e_j,e_k$ con sus multiplicidades. Ahora tendremos un montón, exactamente $27$ pesos distintos.
$$e_1 \otimes e_1 \otimes e_1 \rightarrow (1,0)+(1,0)+(1,0) =(3,0) $$
$$e_1 \otimes e_1 \otimes e_2 \rightarrow (1,0) + (1,0)+(-1,1) =(1,1)$$
$$e_1 \otimes e_1 \otimes e_3 \rightarrow (1,0) + (1,0)+(0,-1) =(2,-1)$$
$$e_1 \otimes e_2 \otimes e_1 \rightarrow (1,0) + (-1,1)+(1,0) =(1,1)$$
$$e_1 \otimes e_3 \otimes e_1 \rightarrow (1,0) + (0,-1)+(1,0) =(2,-1)$$
$$e_1 \otimes e_2 \otimes e_2 \rightarrow (1,0) + (-1,-1)+(-1,-1) =(-1,-2)$$
$$e_1 \otimes e_2 \otimes e_3 \rightarrow (1,0) + (-1,-1)+(0,-1) =(0,-2)$$
$$e_1 \otimes e_3 \otimes e_3 \rightarrow (1,0) + (-1,-1)+(0,-1) =(0,-2)$$
$$ \vdots$$


y las representamos en el diagrama obteniendo

\begin{center}
    \begin{tikzpicture}
        \draw[thick, gray](-3,0)--(3,0);
        \draw[thick, gray](-2,3.46410)--(2,-3.46410);
        \draw[thick, gray](-2,-3.46410)--(2,3.46410);
        \draw[very thick, black](3,0)--(-1.5,2.59024)--(-1.5,-2.598024)--cycle;
         \filldraw[color=black, fill=white, very thick](3,0) circle (0.1);
          \filldraw[color=black, fill=white, very thick](-1.5,2.59024) circle (0.1);
        \filldraw[color=black, fill=white, very thick](-1.5,-2.59024) circle (0.1);
           \filldraw[color=black, fill=white, very thick](0,-1.732) circle (0.1);
           \filldraw[color=black, fill=white, very thick](0,1.732) circle (0.1);
           \filldraw[color=black, fill=white, very thick](1.5,-0.866024) circle (0.1);
        \filldraw[color=black, fill=white, very thick](1.5, 0.866024) circle (0.1);
        \filldraw[color=black, fill=white, very thick](-1.5, 0.866024) circle (0.1);
         \filldraw[color=black, fill=white, very thick](-1.5, -0.866024) circle (0.1);
         \filldraw[color=black, fill=white, very thick](0,0) circle (0.1);
          \node at (3,0)[above]{$1$}; 
          \node at (-1.5,2.59024)[above]{$1$}; 
          \node at (-1.5,-2.59024)[above left]{$1$}; 
          \node at (0,-1.732)[above]{$3$};
          \node at (0,1.732)[above]{$3$};
          \node at (1.5,-0.866024)[above]{$3$}; 
          \node at (1.5,0.866024)[above]{$3$};
          \node at (-1.5,-0.866024)[above left]{$3$};
          \node at (-1.5,0.866024)[above left]{$3$};
          \node at (0,0)[above]{$6$};     
            \draw[thick, gray](-2,3.46410)--(2,-3.46410);
        \draw[thick, gray](-2,-3.46410)--(2,3.46410);
         \node at (2,3.46410)[above]{$l$}; 
        \node at (3,0)[below]{$k$};
    \end{tikzpicture}
\end{center}
Observamos primeramente que el polígono exterior tiene pesos de multiplicidad mayor que uno. Esto es porque, como esperábamos la representación no es irreducible. Elijamos por ejemplo el peso $(3,0)$ sabemos que este será el peso máximo de una representación, que será la $(3,0)$. Podemos construir la representación, que ya lo hemos hecho, y luego <<restarla>> así nos queda que $(1,0)\otimes (1,0) \otimes (1,0) = (3,0) \oplus$...
\begin{center}
    \begin{tikzpicture}
        \draw[thick, gray](-3,0)--(3,0);
        \draw[thick, gray](-2,3.46410)--(2,-3.46410);
        \draw[very thick, black](1.5,0.8660254)--(0,1.73200)--(-1.5,0.8660254)--(-1.5,-0.8660254)--(0,-1.73200)--(1.5,-0.8660254)--cycle;
        \filldraw[color=black, fill=white, very thick](1.5,0.8660254) circle (0.1);
         \filldraw[color=black, fill=white, very thick](-1.5,0.8660254) circle (0.1);
        \filldraw[color=black, fill=white, very thick](-1.5,-0.8660254) circle (0.1);
        \filldraw[color=black, fill=white, very thick](1.5,-0.8660254) circle (0.1);
        \filldraw[color=black, fill=white, very thick](0,1.73200) circle (0.1);
        \filldraw[color=black, fill=white, very thick](0,-1.73200) circle (0.1); 
        \filldraw[color=black, fill=white, very thick](0,0) circle (0.1);
        \node at (0,1.73200)[above right]{$2$};
        \node at (-1.5,0.8660254)[above left]{$2$};
        \node at (-1.5,-0.8660254)[above left]{$2$};
        \node at (0,-1.73200)[below]{$2$};
        \node at (1.5,-0.8660254)[above right]{$2$};
         \node at (1.5,0.8660254)[above right]{$2$};
         \node at (0,0)[below left]{$5$};
           \draw[thick, gray](-2,3.46410)--(2,-3.46410);
        \draw[thick, gray](-2,-3.46410)--(2,3.46410);
         \node at (2,3.46410)[above]{$l$}; 
        \node at (3,0)[below]{$k$};
    \end{tikzpicture}
\end{center}
De nuevo, uno puede hacer lo mismo que antes, tomamos un peso como por ejemplo $(1,1)$ este será el peso máximo de una representación. Es fácil ver que tenemos en este caso $2$ copias de esta representación más la representación trivial de peso $(0,0)$ por tanto, escribimos que 
$$(1,0)\otimes (1,0) \otimes (1,0) = (3,0) \oplus (1,1) \oplus (1,1) \oplus (0,0)$$
O como es más común, según la dimensión 
$$3\otimes 3\otimes 3 = 10\oplus 8\oplus 8 \oplus 1$$ por tanto, que la distribución de bariones tuviera esas simetrías no era casualidad. A los bariones de spin $3/2$ les faltaba una partícula para ser la $10$. Así, Gell-Mann postuló la existencia de esta partícula, el barión $\Omega^{-}$ prediciendo que tendría extrañeza $-3$ y carga $-1$. Usando otros métodos, que están fuera del objetivo del video, estimó también su masa, lo que se conoce como la fórmula de Gell-Mann-Okubo. 2 años más tarde, en el 1964, esta partícula fue descubierta.
\begin{center}
    \begin{tikzpicture}
            \draw[very thick, black](3,0)--(-1.5,2.59024)--(-1.5,-2.598024)--cycle;
        %\draw[thick, gray](-4,0)--(4,0);
        \draw[thick, gray](-3,3.4640)--(4.5,-0.8666);
        \draw[thick, gray](-3,1.732)--(3,-1.732);
        \draw[thick, gray](-4.5,0.866)--(1.5,-2.5980);
        \draw[thick, gray](-3,-1.732)--(0,-3.4640);
        %\draw[thick, gray](-2,3.46410)--(2,-3.46410);
        %\draw[thick, gray](-2,-3.46410)--(2,3.46410);
        %\draw[very thick, black](3,0)--(-1.5,2.59024)--(-1.5,-2.598024)--cycle;
         \draw[thick, gray](-1.5,-3.5)--(-1.5,3.5);
         \draw[thick, gray](0,-3.5)--(0,3.5);
         \draw[thick, gray](1.5,-3.5)--(1.5,3.5);
         \draw[thick, gray](3,-3.5)--(3,3.5);
         
        \node at (-1.5,3.5) [above] {$q=-1$};
        \node at (0,3.5)[above]{$q=-0$};
        \node at (1.5,3.5) [above]{$q=+1$};
        \node at (3,3.5) [above] {$q=+2$};
         
        \node at (4.5,-0.8666) [above right] {$s=0$};
        \node at (3,-1.732) [above right]{$s=-1$};
        \node at (1.5,-2.5980) [above right]{$s=-2$};
        \node at (0,-3.4640) [above right] {$s=-3$};

        
         \filldraw[color=black, fill=white, very thick](3,0) circle (0.1);
          \filldraw[color=black, fill=white, very thick](-1.5,2.59024) circle (0.1);
        \filldraw[color=black, fill=white, very thick](-1.5,-2.59024) circle (0.1);
           \filldraw[color=black, fill=white, very thick](0,-1.732) circle (0.1);
           \filldraw[color=black, fill=white, very thick](0,1.732) circle (0.1);
           \filldraw[color=black, fill=white, very thick](1.5,-0.866024) circle (0.1);
        \filldraw[color=black, fill=white, very thick](1.5, 0.866024) circle (0.1);
        \filldraw[color=black, fill=white, very thick](-1.5, 0.866024) circle (0.1);
         \filldraw[color=black, fill=white, very thick](-1.5, -0.866024) circle (0.1);
         \filldraw[color=black, fill=white, very thick](0,0) circle (0.1);
          \node at (3,0)[above right]{$\Delta^{++}$}; 
          \node at (-1.5,2.59024)[above right]{$\Delta^{-}$}; 
          \node at (-1.5,-2.45)[above right]{$\Omega^{-}$}; 
          \node at (0,-1.54)[above right]{$\Xi^{* 0}$};
          \node at (0,1.732)[above right]{$\Delta^0$};
          \node at (1.5,-0.60)[above right]{$\Sigma^{* +}$}; 
          \node at (1.5,0.866024)[above right]{$\Delta^{+}$};
          \node at (-1.5,-0.866024)[above right]{$\Xi^{* -}$};
          \node at (-1.5,0.866024)[above right]{$\Sigma^{* -}$};
          \node at (0,0)[above right]{$\Sigma^{* 0}$};
          
    \end{tikzpicture}
\end{center}
\section{$SL(2,\mathbb{C})$ y el grupo de Poincaré}
(esto debe ir antes, después de $SU(2)$) \\
Empezaremos caracterizando el grupo de Lorentz y sus subgrupos. El grupo de Lorentz $O(1,3)$ es el grupo de todas las trasformaciones que dejan invariantes el producto escalar en un espacio métrico de Minkowski de métrica
$$\eta_{\mu \nu}=\begin{pmatrix}
	1 & 0 & 0 &0 \\ 0 &-1 &0 &0 \\ 0 & 0 & -1 & 0 \\ 0 & 0 & 0& -1
\end{pmatrix}$$
Supongamos $\Lambda \in O(1,3)$
$$x^\mu\rightarrow x'^\mu=\Lambda^\mu_\nu x^\nu$$ 
El producto escalar transforma entonces como
$$x^\mu\eta_{\mu \nu}x^\nu\rightarrow x'^\sigma\eta_{\sigma \rho}x'^\rho=(x^\mu\Lambda^\sigma_\mu)\eta_{\sigma \rho}(\Lambda^\rho_\nu x^\nu)=x^\mu\eta_{\mu \nu}x^\nu$$
Como la última igualdad debe ser cierta $\forall x^\mu \in \mathbb{M}_4$ encontramos que la identidad definitoria de las transformaciones de Lorentz es
$$\Lambda^\sigma_\mu\eta_{\sigma \rho}\Lambda^\rho_\nu=\eta_{\mu \nu}$$
O en forma matricial
$$\Lambda^T\eta\Lambda=\eta$$
Es decir, son las matrices ortogonales en el espacio de Mikowski, lo que hace que ahora se entienda por qué el nombre de $O(1,3)$.\\
Es importante notar ahora dos cosas:\\
La primera es que el determinante de esta transformación no es necesariamente $1$
$$\det\Lambda\det\eta\det\Lambda=\det\eta$$
$$\det\Lambda^2=1$$
$$\det\Lambda=\pm 1$$
La segunda observación es que pasa con la componente $\mu=\nu=0$
$$\Lambda^\sigma_0\eta_{\sigma \rho}\Lambda^\rho_0=\eta_{00}$$
$$\Lambda^\sigma_0\eta_{\sigma \rho}\Lambda^\rho_0=(\Lambda^0_0)^2-\sum_{i}(\Lambda^i_0)^2=1$$
Por tanto
$$\Lambda^0_0=\pm\sqrt{1+\sum_{i}(\Lambda^i_0)^2}$$
Tenemos entonces 4 posibilidades porque tenemos 4 combinaciones de signos por estas dos condiciones. \\
Como $\Lambda^0_0$ actuará en la componente temporal, no tiene sentido (físico) que sea negativo, tampoco tiene sentido que el espacio se invierta (lo que sería determinante -1). Si el determinante de la transformacion es $+1$ entonces el subgrupo se llama Grupo de Lorentz propio $SO(1,3)$ Si tomamos el caso de la componente $\Lambda^0_0>0$ entonces este subgrupo de $O(1,3)$ se llama Grupo de Lorentz propio ortocrono $SO(1,3)^\uparrow$, en el que centraremos nuestro estudio. También existen los operadores:
$$\Lambda_p=\begin{pmatrix}
	1 & 0 & 0 &0 \\ 0 &-1 &0 &0 \\ 0 & 0 & -1 & 0 \\ 0 & 0 & 0& -1
\end{pmatrix}$$
de paridad  y
$$\lambda_T=\begin{pmatrix}
	-1 & 0 & 0 &0 \\ 0 &1 &0 &0 \\ 0 & 0 & 1 & 0 \\ 0 & 0 & 0& 1
\end{pmatrix}$$
de inversión temporal.\\
Podemos entender el grupo de Lorentz $O(1,3)$ como el conjunto de subgrupos:
$$O(1,3)=\{SO(1,3)^\uparrow,\Lambda_pSO(1,3)^\uparrow,\Lambda_TSO(1,3)^\uparrow,\Lambda_p\Lambda_TSO(1,3)^\uparrow\}$$
Por lo que solo debemos centrarnos en encontrar las representaciones de $SO(1,3)^\uparrow$.\\\\
Pasemos ahora a intentar dar una representación.\\
Por las condiciones definitorias del grupo, para describir una T. de Lorentz general necesitamos seis parámetros. Por tanto, si encontramos seis generadores linealmente independientes habremos encontrado una base para el álgebra de este grupo.\\
Lo primero que vemos es que los elementos de $SO(3)$ cumplen la condición para la parte espacial
$$R^T I R=I$$
$$\det(R)=+1$$
Por tanto si dejamos el tiempo intacto, la transformación de Lorentz será
$$\Lambda_{espacio}=\begin{pmatrix}
	1 & \\ & R_{3\times 3}
\end{pmatrix}$$
Los correspondientes generadores deben ser análogos a los de la representación tridimensional de $SO(3)$
$$J_i=\begin{pmatrix}
	0 & \\ & J_i^{(3 dim)}
\end{pmatrix}$$
Con esto ya hemos encontrado $3$ de las $6$ matrices de la base. Para encontrar transformaciones que involucren el tiempo y el espacio analicemos una transformación infinitesimal tal que $\mathcal{O}(\varepsilon^2)=0$
$$\Lambda^\mu_\rho=\delta^\mu_\rho+\varepsilon K^\mu_\rho$$
Si introducimos esta transformación en la condición para ser de Lorentz
$$\Lambda^\mu_\rho\eta_{\mu \nu}\Lambda^\nu_\sigma=\eta_{\rho \sigma}$$
$$(\delta^\mu_\rho+\varepsilon K^\mu_\rho)\eta_{\mu \nu}(\delta^\nu_\sigma+\varepsilon K^\nu_\sigma)=\eta_{\rho\sigma}$$
$$\eta_{\rho\sigma}+\varepsilon K^\mu_\rho\eta_{\mu \sigma}+\varepsilon K^\nu_\sigma \eta_{\rho \nu}
+ \mathcal{O}(\varepsilon^2)=\eta_{\rho\sigma}$$
Por tanto
$$K^\mu_\rho\eta_{\mu \sigma}+K^\nu_\sigma\eta_{\rho \nu}=0$$
O en forma matricial
$$K^T\eta=-\eta K$$
Obteniendo por fin una condición definitoria de las matrices generadoras de Boost
(Poner cuales son las matrices...)
\subsection{El álgebra de Lie $\mathfrak{so(1,3)}^\uparrow$}
Podemos encontrar las relaciones de conmutación de este álgebra por cálculo directo
$$[J_i,J_j]=i\varepsilon_{ijk}J_k,$$ $$[J_i,K_j]=i\varepsilon_{ijk}K_k,$$ $$[K_i,K_j]=-i\varepsilon_{ijk}J_k$$
Podemos definir ahora dos operadores que si conmuten entre si como
$$N_i^{\pm}= \frac{1}{2}(J_i\pm iK_i)$$
Encontrando las siguientes relaciones de conmutación
$$[N_i^+,N_j^+]=i\varepsilon_{ijk}N_k^+$$
$$[N_i^-,N_j^-]=i\varepsilon_{ijk}N_k^-$$
$$[N_i^+,N_j^-]=0$$
¡Estas son precisamente las relaciones de conmutación de $\mathfrak{su(2)}$! Por tanto, hemos descubierto que el álgebra de Lie de $SO(1,3)^\uparrow$ consiste en dos copias de $\mathfrak{su(2)}$ Esto a parte de ser una bonita propiedad son muy buenas noticias porque ya sabemos construir todas las representaciones irreducibles de $\mathfrak{su(2)}$. \\ Sin embargo, el grupo $SO(1,3)$ no es simplemente conexo y recordemos que para grupos que no son simplemente conexos, no hay una inyectividad entre las representaciones irreducibles del álgebra y las representaciones grupo. Por tanto, en vez de encontrar las representaciones irreducibles del álgebra de Lie del grupo de Lorentz, encontramos las representaciones irreducibles del recubrimiento doble del grupo de Lorentz.\\
(Es importante en este mundo haber visto la idea de recubridores de grupos)\\
El grupo recubridor del grupo de Lorentz puede demostrarse ser $SL(2,\mathbb{C})$ y es un recubrimiento doble porque el núcleo es de orden $2$.
Cada representación irreducible del $\mathfrak{su(2)}$ puede identificarse por el esclar $j$ del operador de Casimir en $SU(2)$. Sabemos entonces que podemos etiquetar las representaciones irreducibles del grupo recubridor del grupo de Lorentz por dos enteros o mitad de enteros $j_1,j_2$. Es decir, que deberemos analizar las $(j_1,j_2)$ representaciones con $j_1,j_2=0,1/2,1,..$.\\
Es común y conveniente escribir el álgebra de Lorentz  usando las matrices $M_{\mu \nu}$ definidas por
$$J_i=\frac{1}{2}\varepsilon_{ijk}M_{jk}$$
$$K_i=M_{0i}$$
Con esta definición el álgebra de Lorentz viene definida por la única relación de conmutación
$$[M_{\mu_\nu},M_{\rho \sigma}]=i(\eta_{\mu \rho}M_{\nu\sigma}-\eta_{\mu \sigma}M_{\nu\rho}-\eta_{\nu\rho}M_{\mu\sigma}+\eta_{\nu \sigma}M_{\mu\rho})$$
\subsubsection{Representación $(0,0)$}
La primera representación es para $SU(2)$ es la trivial porque el espacio vectorial es unidimensional para las dos copias de $\mathfrak{su(2)}$. Los generadores deberán ser, por tanto, matrices $1\times1$ y las únicas matrices que cumplen las relaciones de conmutacion son trivialmente 0
$$N_i^+=N_i^-=0$$
Por tanto, la representación $(0,0)$ actua en objetos que no cambian bajo transformaciones de Lorentz. Esta representación se llama representación escalar del grupo de Lorentz.
\subsubsection{Representación $(1/2,0)$}
En esta representación vamos a usar la representación bidimensional de la copia de $\mathfrak{su(2)}$ $$N_i^+=\frac{\sigma_i}{2}$$ donde $\sigma_i$ son las matrices de Pauli. Y una representación unidimensinal $$N_i^-=0$$ 
Recordemos que esto se extrae de la relación $\text{dimension}=2j+1$\\
Por la definición de $N^-$ 
$$N_i^-=\frac{1}{2}(J_i-iK_i)=0$$
$$J_i=iK_i$$
Ahora para $N^+$
$$N_i^+=\frac{1}{2}(J_i+iK_i)=\frac{1}{2}(iK_i+iK_i)=iK_i$$
Con esto es facil ver que
$$iK_i=\frac{\sigma_i}{2}$$
por tanto
$$K_i=\frac{\sigma_i}{2i}=\frac{-i}{2}\sigma_i$$
Y por otro lado como 
$$J_i=iK_i=\frac{-i^2}{2}\sigma_i=\frac{1}{2}\sigma_i $$
Lo importante es ver ahora que las transformaciones de Lorentz vendrán representadas por matrices $2\times2$ complejas. Estas transformaciones no pueden, por tanto, actuar en cuadrivectores del espacio de Minkowski. Debe existir un objeto de dos componentes en los que actue esta transformación. Estos objetos se llaman espinores levógiros (Left-chiral spinors):
$$\begin{pmatrix}
	(\chi_L)_1 \\ (\chi_L)_2
\end{pmatrix}$$
De momento los espinores no son más que objetos de dos componentes. De momento la definición que daremos de los espinores es que son objetos que transforman bajo la representación $(1/2,0)$ del grupo $O(1,3)$. 
\subsubsection{Representación $(0,1/2)$}
(Equivalente a la otra, pero escribir)
Concluimos que no son el mismo objeto, aunque son muy similares. Llamaremos a estos objetos espinores dextrógiros o espinores de Weyl.\\
\subsubsection{Notación de Van der Waerden}
Los espinores levógiros se definiran por un índice bajo sin punto 
$$\chi_L=\chi_a$$
Los espinores dextrógiros tendran un indice arriba con punto
$$\chi_R=\chi^{\dot{a}}$$
Introducimos ahora la métrica espinorial, lo que nos permitirá transformar de espinores R a espinores L y viceversa. Definimos la métrica espinorial como
$$\epsilon^{ab}=\begin{pmatrix}
	0 & 1 \\ -1 & \\
\end{pmatrix}$$
Vamos a definir tambien lo siguiente:
$$\chi_L^C=\epsilon\chi_L^*$$
Donde $+$ se refiere, por supuesto, a tomar el conjugado complejo. Este objeto tomará sentido más adelante.\\
Vamos a ver como $\chi^C_L$ transforma bajo transformaciones de Lorentz. Es importante primero tener en centa las identidades:
$$ (-\epsilon)(\epsilon)=1$$
$$(\epsilon)\sigma_i^*(-\epsilon)=-\sigma_i$$
Ahora sí, veamos que una transformación de Lorentz a este último objeto resulta
$$X'^C_L=\epsilon (X'_L)^*=\epsilon(e^{\frac{\vec{\phi}}{2}\cdot\vec{\sigma}}\chi_L)^*=\epsilon e^{\frac{\vec{\phi}}{2}\cdot\vec{\sigma}}(-\epsilon)(\epsilon)\chi_L^*)=e^{-\frac{\vec{\phi}}{2}\cdot\vec{\sigma}}\epsilon \chi_L^*=e^{-\frac{\vec{\phi}}{2}\cdot\vec{\sigma}}X_L^C
$$
Que es como transforman los R espinores. \\\\
¿Hablar de los Lagrangianos generales de cada cosa?

$$\mathfrak{s}\mathfrak{o}(1,3)_\mathbb{C}\cong \mathfrak{s}\mathfrak{u}(2)_ \mathbb{C} \oplus \mathfrak{s}\mathfrak{u}(2)_\mathbb{C}  \cong \mathfrak{s}\mathfrak{l}(2,\mathbb{C})\otimes \mathfrak{s}\mathfrak{l}(2,\mathbb{C})  $$
\section{Teorías gauge}
En este nuevo hilo, vamos a revisar rápidamente como se reescriben las ecuaciones de Maxwell en el espacio de Minkowski y luego vamos a analizarlas detenidamente para darnos cuenta de que presentan ciertas estructuras que ya conocemos. \\\\
Este hilo va para toda la gente que esté cursando electrodinámica, principalmente, o cualquier asignatura de mecánica teórica que les haya generalizado las ecuaciones de Maxwell.\\\\
Es, a su vez, la cuarta parte de la saga de grupos de Lie, así que para entender todo en más detalle, mejor haber visto aunque sea el primer hilo, de todos modos, las definiciones necesarias las menciono ahora: 
def de grup, algebra de lie etc.\\
\\\\
Las ecuaciones de Maxwell en el vacío son una serie de ecuaciones diferenciales parciales acopladas que describen todos los fenómenos electromagnéticos. Y son:
\begin{equation}
    \nabla \cdot \vec E =\frac{\rho}{\varepsilon_0}
\end{equation}
\begin{equation}
    \nabla \cdot \vec B =0
\end{equation}
\begin{equation}
    \nabla \times \vec E =-\frac{\partial \vec B}{\partial t}
\end{equation}
\begin{equation}
    \nabla \times \vec B =\mu_0 \vec J +\mu_0\varepsilon_0 \frac{\partial \vec E}{\partial t}
\end{equation}
Hay varias cosas que podemos comentar de estas ecuaciones. En cierto sentido parecen muy simétricas pero en realidad no lo son. Esto es debido principalmente al hecho experimental de que solo existen las cargas eléctricas. Si eliminamos las cargas del juego las ecuaciones de Maxwell rezan:
\begin{equation}
    \nabla \cdot \vec E =0
\end{equation}
\begin{equation}
    \nabla \cdot \vec B =0
\end{equation}
\begin{equation}
    \nabla \times \vec E +\frac{\partial \vec B}{\partial t}=0
\end{equation}
\begin{equation}
    \nabla \times \vec B -\mu_0\varepsilon_0 \frac{\partial \vec E}{\partial t}=0
\end{equation}
Que son también llamadas ecuaciones de Maxwell homogéneas. De estas cuatro ecuaciones fue un hito histórico el hecho (atribuido a Maxwell) de que se pueden resumir en dos ecuaciones de ondas
\newpage
\begin{equation}
    \nabla^2 \vec E =\mu_0\varepsilon_0\frac{\partial^2 \vec E}{\partial t^2}
\end{equation}

\begin{equation}
    \nabla^2 \vec B =\mu_0\varepsilon_0\frac{\partial^2 \vec B}{\partial t^2}
\end{equation}
Que describen una onda formada por los campos $\vec E$ y $\vec B$ perpendiculares entre sí y que se propaga por el vacío con velocidad
\begin{equation}
    c=\frac{1}{\sqrt{\mu_0\varepsilon_0}}
\end{equation}
Lo que hoy llamamos onda electromagnética, todos conocemos, pero no amamos etc. etc. Pero otra cosa, no tan conocida de las ecuaciones de Maxwell homogéneas es que, efectivamente, presentan una simetría. Esta simetría se llama dualidad electromagnética y está dada por la transformación 
$$\{\vec E/c\mapsto \vec B, \vec B \mapsto -\vec E/c\}$$
Esta transformación además, nos <<separa>> las cuatro ecuaiones en dos pares y nos pasa de un par de ecuaciones al otro. Partamos del par que, en el caso no homogéneo tiene cargas y corrientes 
\begin{equation}
    \nabla \cdot \vec E =0
\end{equation}
\begin{equation}
    \nabla \times \vec B -\mu_0\varepsilon_0 \frac{\partial \vec E}{\partial t}=0
\end{equation}
Si ahora aplicamos la dualidad EM tendremos las ecuaciones
\begin{equation}
    \nabla \cdot \vec B =0
\end{equation}
\begin{equation}
    -\nabla \times \vec E -\frac{\partial \vec B}{\partial t}=0
\end{equation}
Que son el otro par de ecuaciones de Maxwell. En efecto, si introducimos la cantidad compleja
\begin{equation}
    \vec{ \mathcal{E}}=\vec E + i\vec B
\end{equation}
La dualidad electromagnética es simplemente realizar la transformación 
\begin{equation}
    \vec {\mathcal{E}}\mapsto -i\vec {\mathcal{E}}
\end{equation}
Y las ecuaciones de Maxwell homogéneas quedarían reducidas a dos 
\begin{equation}
    \nabla \cdot \vec{\mathcal{E}}=0 
\end{equation}
\begin{equation}
    \nabla \times \vec{\mathcal{E}} =i\frac{\partial \vec{\mathcal{E}}}{\partial t}
\end{equation}
La dualidad electromagnética sugiere fuertemente que las ecuaciones de Maxwell deberían ser dos. Además, por la primera ecuación si existe el vector complejo $\vec {\mathcal{A}}$ tal que 
\begin{equation}
    \vec{\mathcal{E}}= \nabla \times \vec {\mathcal{A}}
\end{equation}
La primera ecuación es tautológica y las ecuaciones de Maxwell se reducen a una 
\begin{equation}
    \nabla\times (\nabla\times  \vec {\mathcal{A}}) =\nabla (\nabla \cdot \vec {\mathcal{A}}) -\nabla^2 \vec {\mathcal{A}} =i\frac{\partial }{\partial t}( \nabla \times \vec {\mathcal{A}})
\end{equation}
Sin embargo, nuestra fantasía acaba cuando recordamos que en nuestro mundo existen las cargas y las corrientes. Si queremos introducirlas en nuestras ecuaciones de Maxwell complejas debemos introducirlas de formas complejas como
\begin{equation}
    \rho =\rho_e+ i\rho_m
\end{equation}
\begin{equation}
    \vec j =\vec j_e+i\vec j_m
\end{equation}
Donde el subíndice $e$ índia electricas y $m$ magnéticas. Con esto, las ecuaciones de Maxwell complejas quedan 
\begin{equation}
    \nabla \cdot \vec{\mathcal{E}}=\rho 
\end{equation}
\begin{equation}
    \nabla \times \vec{\mathcal{E}} =i\left(\frac{\partial \vec{\mathcal{E}}}{\partial t}+\vec j\right)
\end{equation}
Sin embargo no parece que existan cargas ni corrientes magnéticas porque no existen los monopolos magnéticos así que esto es pura especulación. Sin embargo, podemos usar la dualidad electromanética a  nuestro favor. \\
Las ecuaciones de Maxwell en forma covariante en el espacio de Minkowski se escriben:
\begin{equation}
    \begin{split}
        \partial_\mu F^{\mu \nu}=\mu_0 \mathcal{J}^\nu  \\
        \frac{1}{2}\partial_\mu \epsilon^{\mu \nu \alpha \beta} F_{\alpha \beta}= 0
    \end{split}
\end{equation}
Donde $F^{\mu \nu}$ es el llamado tensor campo electromagnético y $\mathcal{J}^\nu$ es la generalización de la corriente al espacio de Minkowski,la $4-$corriente. Vamos a detenernos un rato en el tensor electromagnético. Existen dos formas distintas pero equivalentes de definirlo. La primera es, simplemente, como la matriz 
\begin{equation}
    [F^{\mu \nu}] =\begin{pmatrix}
        0 & -E_x/c & -E_y/c & -E_z/c \\
        E_x/c & 0 & -B_z & B_y \\ 
        E_y/c & B_z& 0& -B_x \\
        E_z/c &-B_y & B_x&0 
    \end{pmatrix}
\end{equation}
Que nos relaciona los campos eléctrico y magnético. Si ahora definimos el tensor campo electromagnético dual (de Hodge) 
\begin{equation}
    \mathcal{F}{\mu \nu} = \frac{1}{2}\epsilon^{\mu \nu \alpha \beta } F_{\alpha \beta}
\end{equation} 
Se puede ver que la relación entre $F^{\mu \nu}$ y $\mathcal{F}^{\mu \nu}$ es precisamente la transformación dual que comentábamos antes $\{\vec E/c\mapsto \vec B, \vec B \mapsto -\vec E/c\}$
\begin{equation}
    [\mathcal{F}^{\mu \nu}] =\begin{pmatrix}
        0 & -B_x & -B_y & -B_z\\
        B_x & 0 & E_z/c& -E_y/c \\ 
        B_y & -E_z/c& 0& E_x/c \\
        B_z &E_y/c &-E_x/c&0 
    \end{pmatrix}
\end{equation}
Como veíamos antes, la transformación dual nos pasaba del par de ecuaciones inhomogéneas al par de ecuaciones homogéneas, y viceversa, la idea entonces de estas ecuaciones es que la primera nos da las ecuaciones inhomogéneas de maxwell, luego realizamos una transformación dual a estas ecuaciones y pasamos a las segundas, pero como están en ausencia de cargas, simplemente hacemos que la corriente en esta ecuación sea nula y ya. Hasta aquí todo bien, quedaría demostrar que efectivamente recuperamos las ecuaciones de Maxwell, pero esto es solo calculito. \\ La idea que podemos empezar a olernos, de alguna forma, es que la primera ecuación de Maxwell contiene a la otra, de manera que solo necesitaríamos una ecuación. Sin embargo, no parece ser el caso, no hay forma de deducir la segunda de la primera. De momento, claro, de eso va el hilo.\\\\
Todo cambia cuando, definimos $F$ en términos de un potencial vector. Primero, vamos a ver qué es eso del potencial vector. \\
Ya comentábamos antes en las ecuaciones de Maxwell complexificadas que si existía $\vec {\mathcal{A}}$ tal que 
\begin{equation}
    \vec{\mathcal{E}}= \nabla \times \vec {\mathcal{A}}
\end{equation} 
la primera ecuación era tautológica y las ecuaciones de Maxwell se reducían a una. Por supuesto, este vector ya puede definirse para las cuatro ecuaciones. De la divergencia nula del campo magnético se deduce que es el rotacional de algún vector $\vec A$ 
\begin{equation}
    \vec B =\nabla \times \vec A
\end{equation}
Y la ley de Faraday-Lenz si sustituimos se deduce
\begin{equation}
    \nabla \times \left( \vec E +\frac{\partial \vec A}{\partial t}\right) =0
\end{equation}
Como su rotacional es cero debe ser el gradiente de algo 
\begin{equation}
    \vec E+ \frac{\partial \vec A}{\partial t}=-\nabla \phi 
\end{equation}
Esto es para el espacio euclídeo de toda la vida. Si pasamos a Minkowski el potencial vector admite una forma covariante el $4-$vector:
\begin{equation}
    A^\mu =(\phi/c, \vec A)
\end{equation}
Usando el potencial vector, el tensor campo electromagnético se define como
\begin{equation}
    F^{\mu \nu}= \partial^\mu A^\nu -\partial^\nu A^\mu 
\end{equation}
Se puede calcular explícitamente que obtenemos la misma matriz que antes. Sin embargo, ahora va a pasar una cosa, al igual que antes que la definición del potencial vector complexificado nos hacía la primera ecuación de Maxwell una tautología, lo mismo va a pasar ahora. Esto es porque  definir el tensor campo electromagnético en términos del potencial vector implica la identidad de Bianchi 
\begin{equation}
    \partial_\lambda F_{\mu \nu}+\partial_\mu F_{\nu \lambda}+ \partial_\nu F_{\lambda \mu}=0
\end{equation}
También suele escribirse con la notación de antisimetrización de los índices  
\begin{equation}
    \partial_{[\lambda} F_{\mu \nu]} = 0
\end{equation} 
Donde esto solo es una forma compacta de poner 
\begin{equation}
    \partial_{[\lambda} F_{\mu \nu]} =\frac{1}{3!}\left(\partial_\lambda F_{\mu \nu}+\partial_\nu F_{\lambda \mu}+ \partial_\mu F_{\nu \lambda}-\partial_\lambda F_{\nu \mu}-\partial_\nu F_{\mu \lambda}-\partial_\mu F_{\lambda \nu} 
    \right)
\end{equation}
La identidad de Bianchi se da y se da solo si definimos $F$ en términos de $A$ porque 
\begin{equation*}
    \begin{split}
      &  \partial_\lambda (\partial_\mu A_\nu -\partial_\nu A_\mu)+\partial_\mu (\partial_\nu A_\lambda -\partial_\lambda A_\nu)+\partial_\nu (\partial_\lambda A_\mu -\partial_\mu A_\lambda)  \\ 
    & = \partial_\lambda\partial_\mu A_\nu -\partial_\lambda\partial_\nu A_\mu +\partial_\mu\partial_\nu A_\lambda-\partial_\mu \partial_\lambda\partial_\nu +\partial_\nu\partial_\lambda A_\mu -\partial_\nu\partial_\mu A_\lambda= 0
    \end{split}
\end{equation*}
La identidad de Bianchi da trivialidades si se repiten índices, pero si no, es muy fácil ver que se obtienen las dos ecuaciones de Maxwell homogéneas. \\
Si  $\mu,\nu,\lambda \neq 0$ Tenemos un único caso
\begin{equation*}
\begin{split}
        & \partial_1 F_{23}+\partial_2 F_{31}+\partial_3 F_{12}= \\ 
    & \frac{\partial F_{23}}{\partial x^1} + \frac{\partial F_{31}}{\partial x^2}+\frac{\partial F_{12}}{\partial x^3}=0 \\
    & \frac{\partial (-B_x)}{\partial x} +\frac{\partial (-B_y)}{\partial y} +\frac{\partial (-B_z)}{\partial z}= 0 
\end{split}
\end{equation*}
Lo que implica 
\begin{equation}
    \nabla \cdot \vec B =0
\end{equation}
Si uno de los índices es 0 y los otros no. Supongamos $\mu=0, \nu =1, \lambda=2$ 
\begin{equation*}
\begin{split}
        & \partial_2 F_{01}+\partial_0 F_{12}+\partial_1 F_{20}= \\ 
    & \frac{\partial F_{01}}{\partial x^2} + \frac{\partial F_{12}}{\partial x^0}+\frac{\partial F_{20}}{\partial x^1}=0 \\
    & \frac{\partial (E_x/c)}{\partial y} +\frac{\partial (-B_z)}{\partial (ct)} +\frac{\partial (-E_y/c)}{\partial x}= 0  \\
    & \frac{\partial E_y}{\partial x}-\frac{\partial E_x}{\partial y}=-\frac{\partial B_z}{\partial t}
\end{split}
\end{equation*}
Si lo hacemos para las tres combinaciones posibles obtenemos las tres componentes de la ecuación de Faraday-Lenz
\begin{equation}
    \nabla \times \vec E =-\frac{\partial \vec B}{\partial t}
\end{equation}
Por tanto, la identidad de Bianchi implica la ecuación homogénea de maxwell 
 \begin{equation}
        \frac{1}{2}\partial_\mu \epsilon^{\mu \nu \alpha \beta} F_{\alpha \beta}= 0
\end{equation}
deSin embargo, la identidad de Bianchi está dada solo por definir $F$ en términos de $A$ por tanto esta ecuación es tautológica y las ecuaciones de Maxwell en forma covariante, por tanto, se reducen a una sola
\begin{equation}
        \partial_\mu F^{\mu \nu}=\mu_0 \mathcal{J}^\nu  
\end{equation}
Lo que quiero haceros ver ahora, amigos mios, es que la identidad de Bianchi se da porque estamos trabajando con elementos evaluados en un álgebra de Lie. Vamos a verlo.\\
Definamos el siguiente operador que, de momento, ni nos interesa sobre qué actua 
\begin{equation}
    D_\mu =\partial_\mu +A_\mu 
\end{equation}
Si ahora hacemos el conmutador de el consigo mismo vemos
\begin{equation}
    [D_\mu, D_\nu]=D_\mu D_\nu +D_\nu D_\mu = \partial_\mu A_\nu +\partial_\nu A_\mu +[A_\mu, A_\nu]
\end{equation}
Por tanto
\begin{equation}
    [D_\mu ,D_\nu]=F_{\mu \nu}+[A_\mu,A_\nu]
\end{equation}
Obtenemos entonces que
\begin{equation}
    [D_\mu, D_\nu] = F_{\mu \nu} 
\end{equation}
Si $[A_\mu,A_\nu]=0$ o sea si el álgebra de Lie subyacente es conmutativa.  Vamos a decir que esto es así aunque no lo vamos a demostrar mucho de momento. Bueno. Si ahora consideramos el conmutador
\begin{equation}
    [D_\lambda, F_{\mu\nu}]=\partial_\lambda F_{\mu \nu} + A_\lambda F_{\mu \nu} -F_{\mu \nu}\partial_\lambda -A_\lambda F_{\mu \nu}.
\end{equation}
Fijémonos que el término $-F_{\mu\nu} \partial_\lambda$ es nulo esto es porque si actua sobre una base $\partial_\alpha$
\begin{equation}
    -F_{\mu\nu} \partial_\lambda (\partial_\alpha) =-\partial_\lambda(F_{\mu\nu} \partial_\alpha)+\partial_\alpha\partial_\lambda F_{\mu\nu} = 0
\end{equation}
Por tanto
\begin{equation}
    [D_\lambda, F_{\mu\nu}]=\partial_\lambda F_{\mu \nu} 
\end{equation}
Ahora, la proposición 
\begin{equation*}
    [D_\lambda, F_{\mu\nu}] +[D_\nu, F_{\lambda\mu}] +[D_\mu F_{\nu\lambda}] = 0
\end{equation*}
Es equivalente a la identidad de Bianchi 
\begin{equation}
    \partial_\lambda F_{\mu \nu}+\partial_\mu F_{\nu \lambda}+ \partial_\nu F_{\lambda \mu}=0
\end{equation}
Sin embargo, fijaos que como $F_{\mu\nu}=[D_\mu, D_\nu]$ la suma anterior es 
\begin{equation*}
    [D_\lambda,[D_\mu,D_\nu]]+[D_\nu, [D_\lambda, D_\mu]]+ [D_\mu, [D_\nu, D_\lambda]]=0
\end{equation*}
Como la identidad de Bianchi está dada  (dada la existencia del campo $A$, como ya hemos demostrado) y estas dos proposiciones son equivalentes, entonces la suma anterior es, efectivamente nula. \\
Sin embargo, nos topamos ahora con algo muy interesante, puesto que la anterior proposición no es más que la identidad de Jacobi! Una propiedad que viene dada por definición en todas las álgebras de Lie. Por tanto, el resultado que podemos empezar a ver es que, efectivamente, $A$ es un elemento de un álgebra de Lie pero ¿Cómo? ¿Qué tiene qué ver aquí un álgebra de Lie? ¿Es que tenemos, acaso alguna simetría? \\
Habéis podido observar que he evitado hablar dela simetría más importante del electromagnetismo, la simetría gauge. Esta simetría nos dice que la transformación 
\begin{equation}
    A'^\mu =A^\mu +\partial^\mu \psi 
\end{equation}
Es una simetría del campo electromagnético en tanto que el tensor $F^{\mu \nu}$ queda igual y por tanto se obtienen los mismos campos. Vamos a ver, realmente, qué es la simetría gauge.\\
\section*{Fibrados}
Vamos a considerar un tipo de espacios muy especiales, estos espacios se llaman Fibrados $G-$principales o simplemente $G-$fibrados. No vamos a dar ninguna teoría de fibrados aquí, solo vamos a dar las ideas más sencillas.\\
La idea del fibrado es que dada una variedad $M$ (recordad que una variedad es un espacio localmente homeomorfo a $\mathbb{R}^n$ o sea, que localmente parece plano) podemos, a cada punto de la variedad, asignarle otro espacio $E_p$, de  manera que existe una apliación proyección $\pi :E_p\mapsto p$. El espacio $E_p$ se llama fibra de $E$ en $p$ Donde $E$ es el llamado espacio total que simplemente se define como la unión de todas las fibras:
\begin{equation}
    E=\bigcup_{p\in M}E_p
\end{equation}
Como la aplicación proyección debe existir para cada punto entonces normalmente se dice que un fibrado es $\pi :E\rightarrow M$ y decimos que tenemos un fibrado $E$ sobre $M$. \\
Un ejemplo muy sencillo de fibrado es el cilindro. El cilindro es quivalente a asignar un segmento a cada punto del círculo decimos entonces que el cilindro es un fibrado sobre $S^1$, el círculo. En este caso, el cilindro simplemente se escribe $S^1\times I $  y como sepuede escribir así lo llamamos fibrado trivial ¿cual es la diferencia entre el espacio producto $S^1\times I$ y un fibrado? Que el fibrado nos permite hacerlo \textit{localmente}. Por ejemplo, la cinta de Mobius que tanto le gusta a internet, es casi un cilindro, ocurre que cuando <<se cierra el cilindro>> los puntos opuestos se identifican, le damos un giro y lo cerramos. La banda de Mobius es un fibrado sobre $S^1$ pero no es trivial. Sin embargo, localmente, esto es, si cogemos una sección del círculo, la cinta de mobius parece un cilindro. Esta es la gracia de los fibrados. Finalmente, si la fibra es un espacio vectorial, entonces decimos que es un fibrado vectorial. \\
Los fibrados, además, nos dan en física una bonita forma de generalizar los campos. Sea $\pi: E\rightarrow M$ un fibrado, una sección del fibrado es una aplicación $s:\mathcal{M}\rightarrow E$ tal que para cualquier $p\in M$ $s(p)\in E_p$. En otras palabras, una sección asigna a cada punto de la variedad un vector en la fibra de ese punto. \\ Por ejemplo, supongamos el fibrado 
\begin{equation}
    \pi:  {\mathbb{C}} \rightarrow\mathbb{R} 
\end{equation}
Las secciones serán las funciones 
$$\psi: \mathbb{R}\rightarrow \mathbb{C}$$
Que es, por ejemplo, una función de onda unidimensional. La gracia del fibrado es que permite que en vez de $\mathbb{R}$ tengamos cualquier variedad así, por ejemplo podríamos tener 
\begin{equation}
    \pi:  {\mathbb{C}} \rightarrow M 
\end{equation}
Las secciones de este fibrado serán pues las aplicaciones 
Las secciones serán las funciones 
$$\psi: M\rightarrow \mathbb{C}$$
Que podemos, por ejemplo, entender como funciones de ondas definidas por la variedad $M$. Las secciones parecen una forma complicada de hablar de una función de toda la vida, y lo es, lo que ocurrirá es qué paga cuando tratamos de <<derivarlas>> .\\\\
Supongamos un fibrado vectorial $\pi : E\rightarrow M$ trivial que podemos escribir trivialmente como $M\times V$ donde $V$ es el espacio vectorial de alguna representación $\Pi$ de algún grupo de Lie $G$ Recordad que esto es el homomorfismo $\Pi: G\rightarrow GL(V)$ llamamos a este tipo de fibrados $G-$fibrados. Esto implica, por tanto, que podemos <<movernos por la fibra>> usando aplicacones $\Pi(g): V \rightarrow V$ . Esta definción tiene un montón de cosas más y condiciones de compatibilidad y cosas que vamos, fijate tu, a ignorar completamente. \\\\ En las teorías gauge tendremos que los llamados, campos de materia, esto es, los que no son bosones gauge, serán secciones de $G-$fibrados.  
Con esto, damos la siguiente definición:\\
Sea $\pi: E\rightarrow M$ un $G-$fibrado, una transformación gauge es una transformación  $E_p\rightarrow E_p$ que varia suavemente conforme nos movemos de punto $p \in M$  que es un elemento del grupo $G$. \\   
Supongamos el espacio vectorial $V$ al que es isomorfa la fibra. Supongamos una base de este espacio vectorial $\{e_i\}$, así que un vector cualquier de $V$ será $v=v^ie_i$ por supuesto usando la notación de índices sumados. Ahora, tomemos un elemento $\Pi (g)$ tenemos una transformación entre vectores. La base cambia como
$$ e_{i'}=\Pi(g)^{i}_{\ i'} e_{i}$$
Y las componentes como
$$ v^{i'}=(\Pi(g)^{-1})^{i'}_{\ i} v^{i}$$
Realizar este tipo de transformaciones es lo que llamaremos una transformación de gauge, de acuerdo con la anterior definición. \\\
Para acabar, y ya veréis como en seguida volvemos a lo de toda la vida, vamos a ver como derivan las secciones. Esto es, como varian los vectores definidos en las fibras conforme nos movemos por la variedad y, por tanto, nos movemos de fibra. El problema es, precisamente, este. Si cada fibra es un espacio vectorial, lo que va a pasar es que tenemos espacios vectoriales distintos involucrados, los elementos de $E_p$ y los elementos de $E_q$ son distintos y por tanto no podemos operar entre ellos. La derivada es comparar (pensar en la definición de toda la vida) dos cantidades, por tanto, tenemos un problema. Pero no os preocupéis porque nos va a salvar un invento, la derivada covariante y la conexión. La derivda covariante es una generalización de la derivada direccional para secciones, recordemos que la derivada direccional era ver como variaba una función siguiendo una dirección dada por un campo vectorial $\vec v$. Ahora, vamos a definir la derivada covariante, esto es, sea $v$ campo vectorial y $s$ una sección entonces vamos a hablar del operador 
\begin{equation}
    D_v s
\end{equation}
Escribiremos $v=v^\mu \partial_\mu$ y $s=s^i e_i(x^\mu)$. y $D_\mu =D_{\partial \mu}$ Donde estoy queriendo expresar de la base de secciones depende del punto, pero voy a dejar de expresarlo.\\
Para cualesquiera $\mu, i$ podemos expresar $D_\mu e_i(x^\mu)$ como una combinación lineal de secciones $e_i$  
\begin{equation}
    D_\mu e_i =A^j _{\ \mu i}(x^\nu)e_j(x^\nu)
\end{equation}
Los coeficientes $A^j _{\ \mu i}$ se llaman coeficientes de la conexión. Si ahora hacemos actuar $D_\mu $ sobre una forma general lo vamos a definir como
\begin{equation}
    D_\mu s = D_\mu (s^i e_i) =((\partial_\mu s^i)+ A^i_{\ \mu j}s^j)e_i 
\end{equation}
Si suprimimos la sección y las bases podemos expresarlo como el operador
\begin{equation}
    D_\mu =\partial_\mu + A_\mu 
\end{equation}
Si estamos en un $G-$fibrado, la conexión $A_\mu$ esta debe estar evaluada en su álgebra de Lie. Hay varias maneras de ver por qué esto es así. Podemos pensar $A$ como una apliación que come direcciones (un vector $v$) y que actua entonces en las secciones. Si lo hacemos sobre las bases tenemos
\begin{equation}
    A(\partial _\mu) e_j =A^i_{\ \mu j}e_i
\end{equation}
Así que $A$ se encarga de las variaciones infinitesimales de las secciones, esto es precisamente, lo que caracteriza un elemento de un álgebra de Lie. \\

Supongamos ahora que realizamos una transformación gauge, esto es, que realizamos una transformación sobre las secciones, la derivada covariante no puede cambiar entonces tendremos
\begin{equation}
    D_\mu s = D_\mu s' =D_\mu (\Pi(g) s) =D'(s)
\end{equation}
aLo que me dará una nueva derivada covariante $D'$ que será igual que la original, pero me cambiará la conexión. Tenemos entonces, por no marranear mucho vamos a escribir $g$ en vez de $\Pi(g)$ y lo vamos a poner todo matricialmente tenemos entonces
$$(\partial_\mu (s)+A_\mu s)e_i = (\partial_\mu (s')+A_\mu 's')e_{i'} = (\partial_\mu (s')+A_\mu 's')g e_i $$
Tenemos entonces la igualdad en forma matricial
$$\partial_\mu (s)+A_\mu s = g (\partial_\mu (g^{-1}s)+A_\mu ' g^{-1}s)$$
Por derivada del producto tenemos
$$\partial_\mu (s)+A_\mu s = g((\partial_\mu g^{-1})s+ g^{-1} \partial_\mu s+A_\mu ' g^{-1}s)$$
De donde 
$$\partial_\mu (s)+A_\mu s = g(\partial_\mu g^{-1})s+  \partial_\mu s+ g A_\mu ' g^{-1}s $$
Como esto debe ser para todo vector vemos entonces que
\begin{equation}
    A_\mu =g\partial_\mu g^{-1} +g A'_\mu g^{-1} 
\end{equation}
En los libros sale del revés
\begin{equation}
    A'_\mu =g\partial_\mu g^{-1} +g A_\mu g^{-1} 
\end{equation}
Que es lo que llamamos una transformación gauge. \\
Vamos a volver al electromagntismo. Para el electromagnetismo tenemos que decir que tendremos un $U(1)-$fibrado $M\times \mathbb{C}\rightarrow M$ en otras, palabras, la fibra sobre cualquier punto $p$ son los complejos $\mathbb{C}$ que es la representación estándar de $U(1)$, como Si definimos una derivada covariante con la conexión $A_\mu$ esta debe estar evaluada en su álgebra de Lie.
El álgebra de Lie de $U(1)$ son simplemente los imaginarios puros 
\begin{equation}
    \mathfrak{u}(1)=\{ix| x\in \mathbb{R}\}
\end{equation}
Por tanto seremos descuidados con nuestra notación y simplemente diremos que 
$A_\mu \mapsto i A_\mu$
Claro, $g=e^{i\theta(x)}$ La transformación gauge del potencial nos queda
\begin{equation}
    iA'_\mu = e^{i \theta(x^\mu)}\partial_\mu e^{-i\theta(x^\mu)} + e^{i \theta(x^\mu)}iA_\mu e^{-i\theta(x^\mu)} 
\end{equation}
Realizando la derivada y como $A_\mu$ es un número real nos queda
\begin{equation}
    iA'_\mu = -i\partial_\mu (\theta) +i A_\mu 
\end{equation}
Por tanto
\begin{equation}
    A'_\mu =A_\mu - \partial_\mu \theta
\end{equation}
Que es una transformación gauge de toda la vida. Pero a partir de ahora deberemos llamarla transformación gauge $U(1)$\\\\
\noindent 
First of all, $a_i \in {\pm 1}$ so they are real. 

Suppose $\alpha \in \mathbb{C}$  is a root of a polynomial  $$p(z)=\sum_i a_i z^i$$ then 
$$p(\alpha) = \sum_i a_i \alpha^i=0$$ 
Y can take the complex conjugate to all the equation and since the $a_k\in \mathbb{R}$ and get
$$\sum_i a_i \alpha^{*i}=0$$
so the complex conjugate must be a root.

To the $-z$ to be also a root (and therefore $-z^*$ and completing the 4-fold sym.)  it is defined $X_n\subset \mathbb{C}$ as the is the set of roots of some little wood polynomial with $n$ nonzero terms so suppose u have a littlewood polynomial $p(z)$ defined as above then u can change sign of all the coefficients and will still be a polynomial whose roots will lie in $X_n$



\begin{equation}
    \xymatrix{
&Cl(n-1) \ar@{^(->}[r], \ar@{->}[d] & Cl(n) \ar@{->}[d]& \\ 
&Spin(n)\ar@{->}[d]^{\rho}    \ar@{^(->}[r] & Pin(n) \ar@{->}[d]& \\
\mathfrak{s}\mathfrak{o}(n) \ar@{->}[ur]^{\exp} \ar@{->}[d]^{\pi}   &SO(n) \ar@{^(->}[r], \ar@{->}[l], \ar@{->}[dr]^{\Pi}   & O(n) &\\
\mathfrak{g}\mathfrak{l}(\mathbb{R},n) \ar@{->}[rr]& & GL(\mathbb{R},n) \ar@{->}[r] & \mathbb{R}^n \ar@{^(->}@/_{5mm}/[uuul]
}
\end{equation}
\newpage
\begin{equation*}
    \text{Spin}(n+1,1)\cong SL(2,\mathbb{K}) 
\end{equation*}
Where $\mathbb{K}=\mathbb{R},\mathbb{C}, \mathbb{H}, \mathbb{O}$ and $n=\dim \mathbb{K}=1, 2, 4, 8$
\\\\
In particular 
\begin{equation*}
    \text{Spin}(3,1)\cong SL(2,\mathbb{C}) 
\end{equation*}
Gives the double cover of the Lorentz Group $SO(1,3)$
\begin{equation}
    \xymatrix{
&Cl(n-1) \ar@{^(->}[r], \ar@{->}[d] & Cl(n) \ar@{->}[d]& \\ 
&Spin(n)\ar@{->}[d]^{\rho}    \ar@{^(->}[r] & Pin(n) \ar@{->}[d]& \\
\mathfrak{s}\mathfrak{o}(n) \ar@{->}[ur]^{\exp} \ar@{->}[d]^{\pi}   &SO(n) \ar@{^(->}[r], \ar@{->}[l], \ar@{->}[dr]^{\Pi}   & O(n) &\\
\mathfrak{g}\mathfrak{l}(\mathbb{R},n) \ar@{->}[rr]& & GL(\mathbb{R},n)
}
\end{equation}
\newpage

\end{document}
